% Diálogo familiar, debate sobre la violencia de género et traduction : SER y ESTAR — Espagnol Tle A — Cours signé OnBuch+
% Programme officiel MINESEC. Compiler avec : tectonic E02-dialogo-familiar-ser-estar.tex
\newcommand{\DOCMATIERE}{Espagnol}
\newcommand{\DOCNIVEAU}{Tle A}
\newcommand{\DOCMODULE}{Módulo 1 : Vie familiale et intégration sociale}
\newcommand{\DOCLECON}{E02}
\newcommand{\DOCTITRE}{Diálogo familiar, debate sobre la violencia de género et traduction : SER y ESTAR}
% OnBuch+ — préambule « cours signés » ENRICHI (compilé avec tectonic / XeLaTeX)
% Système visuel « Pop » d'OnBuch+ : fond crème (jamais blanc), bordures encre
% épaisses, ombres « dures » décalées, titres Archivo Black en capitales.
% Version MATHÉMATIQUES : graphes (pgfplots/tikz), boîtes pédagogiques
% (définition, propriété, méthode, attention, le savais-tu, exercice résolu,
% exercice, TP, bilan), page de garde et signature OnBuch+ ancrée en bas.
%
% Macros attendues dans main.tex AVANT \input{preamble} :
%   \DOCMATIERE  \DOCNIVEAU  \DOCMODULE  \DOCLECON (numéro)  \DOCTITRE
\documentclass[11pt]{article}

\usepackage{fontspec}
\usepackage[spanish,french]{babel} % français = langue principale ; l'espagnol via \esp{..} / espagnol
\usepackage[a4paper,margin=2cm,top=3.4cm,bottom=2.8cm,headheight=1.4cm,footskip=1.3cm]{geometry}
\usepackage{amsmath,amssymb,amsfonts}
\usepackage{mathtools} % \xmapsto, \coloneqq... souvent utilisés par les modèles
\usepackage{cancel} % simplifications barrées (bilans de forces)
% \abs{z} (module, valeur absolue) et \norm{u} (norme), utilisés par les modèles
\providecommand{\abs}{}\let\abs\relax\DeclarePairedDelimiter{\abs}{\lvert}{\rvert}
\providecommand{\norm}{}\let\norm\relax\DeclarePairedDelimiter{\norm}{\lVert}{\rVert}
\usepackage[table]{xcolor} % [table] : fournit aussi \rowcolors (alternance de lignes)
\usepackage{pagecolor}
\usepackage{tikz}
\usetikzlibrary{arrows.meta,positioning,calc,shapes.geometric,shapes.misc,shapes.arrows,decorations.pathmorphing,decorations.pathreplacing,decorations.markings,patterns,shadows,mindmap,trees,fit,backgrounds,matrix,intersections,angles,quotes}
\usepackage{pgfplots}
\usepackage[version=4]{mhchem} % équations nucléaires \ce{^{235}_{92}U}
\usepackage[european,siunitx]{circuitikz} % schémas électriques
\pgfplotsset{compat=1.18}
\usepgfplotslibrary{fillbetween,groupplots} % aires entre courbes (calcul intégral)
\usepackage{enumitem}
\usepackage{fancyhdr}
% [most] charge toutes les bibliothèques tcolorbox (listings, minted, raster,
% documentation...) et gonfle inutilement la mémoire TeX sur un document long
% (~300 boîtes) : on ne charge que ce qui est réellement utilisé.
\usepackage[skins,breakable]{tcolorbox}
\usepackage{graphicx}
\usepackage{colortbl}
\usepackage{booktabs}
\usepackage{tabularx}
\usepackage{diagbox} % case d'en-tête diagonale des tableaux à double entrée
% Colonnes extensibles centrée / gauche / droite (C, L, R), souvent utilisées par les modèles
\newcolumntype{C}{>{\centering\arraybackslash}X}
\newcolumntype{L}{>{\raggedright\arraybackslash}X}
\newcolumntype{R}{>{\raggedleft\arraybackslash}X}
\usepackage{array}
\usepackage{multicol}
\usepackage{siunitx}
% parse-numbers=false : accepte aussi \SI{2,5 \times 10^{-3}}{...}
\sisetup{locale=FR,output-decimal-marker={,},group-minimum-digits=4,parse-numbers=false}
\DeclareSIUnit\ohm{\ensuremath{\symup{\Omega}}} % U+2126 absent de la police
\DeclareSIUnit\division{div} % réglages d'oscilloscope (ms/div, V/div)
\DeclareSIUnit\tr{tr} % tours (vitesse de rotation, ex. \SI{1500}{\tr\per\minute})
\DeclareSIUnit\molar{M} % concentration molaire (ex. \SI{3}{\molar}, \SI{1}{\micro\molar})
\DeclareSIUnit\an{an} % année (le modèle confond parfois avec \year, un compteur TeX) — ex. \SI{5}{\kilo\meter\per\an}
\DeclareSIUnit\FCFA{FCFA} % franc CFA (ex. \SI{500}{\FCFA\per\kilo\gram})
\DeclareSIUnit\degreeBrix{\textdegree Brix} % teneur en sucre d'un sirop/jus (réfractométrie)
\DeclareSIUnit\month{mois} % le modèle utilise parfois \month, un compteur TeX, comme unité de durée
\usepackage{qrcode}
% \degree : commande fréquemment utilisée par les modèles pour un angle ou une
% température en dehors de \SI{}{} (non fournie par les paquets chargés).
\providecommand{\degree}{^{\circ}}
% \qed (fin de démonstration), utilisé par les modèles sans amsthm
\providecommand{\qed}{\hfill$\square$}
% \barre : double barre des valeurs interdites dans un tableau de variations (tkz-tab)
\providecommand{\barre}{\ensuremath{\|}}
% environnement proof (amsthm non chargé) utilisé par les modèles
\ifdefined\proof\else\newenvironment{proof}[1][Démonstration]{\par\noindent\textit{#1.}\ }{\qed\par}\fi
\usepackage{lastpage}
\usepackage{hyperref}
\hypersetup{hidelinks}

% ============================================================
% Palette « Pop » OnBuch+ — mêmes hex que la classe P (plus_ui.dart)
% ============================================================
\definecolor{poporange}{HTML}{F59321}
\definecolor{popdark}{HTML}{E07A0C}
\definecolor{popcream}{HTML}{FFF6E8}
\definecolor{popink}{HTML}{1C1714}
\definecolor{poppurple}{HTML}{7A5AE0}
\definecolor{popgreen}{HTML}{2E9E6A}
\definecolor{poppink}{HTML}{E0457B}
\definecolor{popblue}{HTML}{2F7DE1}
\definecolor{popgold}{HTML}{C98A12}
\definecolor{popmuted}{HTML}{8A7F76}
\definecolor{popline}{HTML}{EADFD2}
\definecolor{popgray}{HTML}{8A7F76}
\colorlet{popgrey}{popgray}
% Alias tolérants (le modèle nomme parfois les couleurs différemment)
\colorlet{popwhite}{white}
\colorlet{popblack}{popink}
\colorlet{popred}{poppink}
\colorlet{popyellow}{popgold}
% Teintes claires (fonds de tableaux/encadrés) — le modèle nomme parfois
% les couleurs avec un suffixe L (ex. popblueL) sans les avoir définies.
\colorlet{poporangeL}{poporange!20}
\colorlet{popdarkL}{popdark!20}
\colorlet{poppurpleL}{poppurple!20}
\colorlet{popgreenL}{popgreen!20}
\colorlet{poppinkL}{poppink!20}
\colorlet{popblueL}{popblue!20}
\colorlet{popgoldL}{popgold!20}
\colorlet{popredL}{popred!20}
\colorlet{popgrayL}{popgray!20}
% Teintes claires pour fonds de tableaux / zones
\colorlet{popblueL}{popblue!10!white}
\colorlet{poporangeL}{poporange!14!white}
\colorlet{popgreenL}{popgreen!12!white}
\colorlet{poppurpleL}{poppurple!12!white}
\colorlet{poppinkL}{poppink!12!white}
\colorlet{popgoldL}{popgold!14!white}

\pagecolor{popcream}

% ============================================================
% Typographie
% ============================================================
\defaultfontfeatures{Ligatures=TeX}
\setmainfont{PlusJakartaSans-Medium.ttf}[Path=../../../fonts/,BoldFont=PlusJakartaSans-Bold.ttf]
% Maths en sans-serif (Fira Math) avec chiffres et lettres droites de Plus
% Jakarta, pour que les formules se fondent dans le texte.
\usepackage[math-style=ISO,bold-style=ISO]{unicode-math}
\setmathfont{FiraMath-Regular.otf}[Path=../../../fonts/]
\setmathfont{PlusJakartaSans-Medium.ttf}[Path=../../../fonts/,range={up/num,up/{Latin,latin},\mathup}]
\usepackage{icomma} % virgule décimale sans espace en mode math
% Caractères mathématiques Unicode absents des polices de texte (Plus Jakarta,
% Archivo Black) : rendus par leur équivalent mathématique, en texte comme en
% formule — sinon ils s'affichent en carré vide (ex. « Arithmétique dans ℤ »).
\usepackage{newunicodechar}
\newunicodechar{ℝ}{\ensuremath{\mathbb{R}}}
\newunicodechar{ℤ}{\ensuremath{\mathbb{Z}}}
\newunicodechar{ℂ}{\ensuremath{\mathbb{C}}}
\newunicodechar{ℕ}{\ensuremath{\mathbb{N}}}
\newunicodechar{ℚ}{\ensuremath{\mathbb{Q}}}
\newunicodechar{𝔻}{\ensuremath{\mathbb{D}}}
\newunicodechar{α}{\ensuremath{\alpha}}
\newunicodechar{β}{\ensuremath{\beta}}
\newunicodechar{γ}{\ensuremath{\gamma}}
\newunicodechar{δ}{\ensuremath{\delta}}
\newunicodechar{ε}{\ensuremath{\varepsilon}}
\newunicodechar{θ}{\ensuremath{\theta}}
\newunicodechar{λ}{\ensuremath{\lambda}}
\newunicodechar{μ}{\ensuremath{\mu}}
\newunicodechar{σ}{\ensuremath{\sigma}}
\newunicodechar{τ}{\ensuremath{\tau}}
\newunicodechar{φ}{\ensuremath{\varphi}}
\newunicodechar{ψ}{\ensuremath{\psi}}
\newunicodechar{ω}{\ensuremath{\omega}}
\newunicodechar{Σ}{\ensuremath{\Sigma}}
% majuscules produites par \MakeUppercase dans les titres (α-aminés -> Α)
\newunicodechar{Α}{\ensuremath{\alpha}}
\newunicodechar{Β}{\ensuremath{\beta}}
\newunicodechar{Γ}{\ensuremath{\Gamma}}
\newunicodechar{Δ}{\ensuremath{\Delta}}
\newunicodechar{Ε}{\ensuremath{\varepsilon}}
\newunicodechar{Θ}{\ensuremath{\Theta}}
\newunicodechar{Λ}{\ensuremath{\Lambda}}
\newunicodechar{Μ}{\ensuremath{\mu}}
\newunicodechar{Τ}{\ensuremath{\tau}}
\newunicodechar{Φ}{\ensuremath{\Phi}}
\newunicodechar{Ψ}{\ensuremath{\Psi}}
\newunicodechar{Ω}{\ensuremath{\Omega}}
\newunicodechar{↦}{\ensuremath{\mapsto}}
\newunicodechar{∈}{\ensuremath{\in}}
\newunicodechar{∉}{\ensuremath{\notin}}
\newunicodechar{∧}{\ensuremath{\wedge}}
\newunicodechar{⇔}{\ensuremath{\Leftrightarrow}}
\newunicodechar{⟺}{\ensuremath{\Longleftrightarrow}}
\newunicodechar{⇒}{\ensuremath{\Rightarrow}}
\newunicodechar{⟹}{\ensuremath{\Longrightarrow}}
\newunicodechar{∘}{\ensuremath{\circ}}
\newunicodechar{∪}{\ensuremath{\cup}}
\newunicodechar{∩}{\ensuremath{\cap}}
\newunicodechar{≡}{\ensuremath{\equiv}}
\newunicodechar{⊂}{\ensuremath{\subset}}
\newunicodechar{⊥}{\ensuremath{\perp}}
\newunicodechar{∃}{\ensuremath{\exists}}
\newunicodechar{∀}{\ensuremath{\forall}}
\newunicodechar{∛}{\ensuremath{\sqrt[3]{\,}}}
\newunicodechar{∜}{\ensuremath{\sqrt[4]{\,}}}
\newunicodechar{⃗}{\ensuremath{{}^{\to}}}
\newunicodechar{ⁿ}{\ifmmode^{n}\else\textsuperscript{\ensuremath{n}}\fi}
\newunicodechar{ˣ}{\ifmmode^{x}\else\textsuperscript{\ensuremath{x}}\fi}
\newunicodechar{ᵇ}{\ifmmode^{b}\else\textsuperscript{\ensuremath{b}}\fi}
\newunicodechar{ʳ}{\ifmmode^{r}\else\textsuperscript{\ensuremath{r}}\fi}
\newunicodechar{ᵘ}{\ifmmode^{u}\else\textsuperscript{\ensuremath{u}}\fi}
\newunicodechar{ʸ}{\ifmmode^{y}\else\textsuperscript{\ensuremath{y}}\fi}
\newunicodechar{ᵃ}{\ifmmode^{a}\else\textsuperscript{\ensuremath{a}}\fi}
\newunicodechar{ᵉ}{\ifmmode^{e}\else\textsuperscript{\ensuremath{e}}\fi}
\newunicodechar{ᵏ}{\ifmmode^{k}\else\textsuperscript{\ensuremath{k}}\fi}
\newunicodechar{ᵖ}{\ifmmode^{p}\else\textsuperscript{\ensuremath{p}}\fi}
\newunicodechar{ᴺ}{\ifmmode^{N}\else\textsuperscript{\ensuremath{N}}\fi}
\newunicodechar{ᶜ}{\ifmmode^{c}\else\textsuperscript{\ensuremath{c}}\fi}
\newunicodechar{ʰ}{\ifmmode^{h}\else\textsuperscript{\ensuremath{h}}\fi}
\newunicodechar{ᵠ}{\ifmmode^{q}\else\textsuperscript{\ensuremath{q}}\fi}
\newunicodechar{ₙ}{\ifmmode_{n}\else\textsubscript{\ensuremath{n}}\fi}
\newunicodechar{ₐ}{\ifmmode_{a}\else\textsubscript{\ensuremath{a}}\fi}
\newunicodechar{ᵢ}{\ifmmode_{i}\else\textsubscript{\ensuremath{i}}\fi}
\newunicodechar{ᵣ}{\ifmmode_{r}\else\textsubscript{\ensuremath{r}}\fi}
\newunicodechar{ₖ}{\ifmmode_{k}\else\textsubscript{\ensuremath{k}}\fi}
\newunicodechar{ᵦ}{\ifmmode_{\beta}\else\textsubscript{\ensuremath{\beta}}\fi}
\newunicodechar{ₓ}{\ifmmode_{x}\else\textsubscript{\ensuremath{x}}\fi}
\newfontfamily\popdisplay{ArchivoBlack-Regular.ttf}[Path=../../../fonts/]
\newfontfamily\poplogo{SpaceGrotesk-Bold.ttf}[Path=../../../fonts/]
\newfontfamily\popsemi{PlusJakartaSans-SemiBold.ttf}[Path=../../../fonts/]
\newfontfamily\popxbold{PlusJakartaSans-ExtraBold.ttf}[Path=../../../fonts/]
\newfontfamily\popxboldls{PlusJakartaSans-ExtraBold.ttf}[Path=../../../fonts/,LetterSpace=3.0]

\setlist[enumerate]{leftmargin=1.7em,itemsep=5pt,topsep=4pt}
\setlist[itemize]{leftmargin=1.7em,itemsep=4pt,topsep=4pt,label={\color{poporange}\textbullet}}
\setlength{\parindent}{0pt}
\setlength{\parskip}{5pt}
\renewcommand{\arraystretch}{1.25}

% pgfplots : style Pop par défaut
\pgfplotsset{
  every axis/.append style={axis line style={popink,line width=1pt},tick style={popink},
    grid style={popline,line width=0.5pt},label style={font=\small},tick label style={font=\footnotesize}},
  popcurve/.style={poporange,line width=1.8pt,no markers,smooth},
}
% Même style disponible dans un \draw TikZ simple (hors pgfplots)
\tikzset{popcurve/.style={poporange,line width=1.8pt}}
\tikzset{popgrid/.style={popline,very thin}}
\colorlet{popcurve}{poporange} % le nom du style est parfois utilisé comme couleur

% ============================================================
% Pill / squiggle
% ============================================================
\newcommand{\poppill}[3]{%
  \tikz[baseline=-0.55ex]\node[draw=popink,line width=1.2pt,rounded corners=2.6pt,%
    fill=#2,inner xsep=6.5pt,inner ysep=3pt]{%
    \popxboldls\fontsize{7}{7}\selectfont\color{#3}\MakeUppercase{#1}};%
}
\newcommand{\popsquiggle}[1]{%
  \tikz[baseline=-0.4ex]{
    \draw[#1,line width=2.6pt,line cap=round]
      plot [smooth] coordinates {(0,0.09) (0.34,0.01) (0.68,0.21) (1.02,0.01) (1.36,0.10)};
  }%
}
% Mot-clé surligné (marqueur orange sous le texte)
\newcommand{\cle}[1]{{\popxbold\color{popdark}#1}}

% ============================================================
% En-tête / pied
% ============================================================
\pagestyle{fancy}
\fancyhf{}
\renewcommand{\headrulewidth}{0pt}
\renewcommand{\footrulewidth}{0pt}
\lhead{\raisebox{-0.15ex}{\poplogo\fontsize{13}{13}\selectfont\color{popink}OnBuch\color{poporange}+}}
\rhead{\poppill{\DOCMATIERE\ \textbullet\ \DOCNIVEAU\ \textbullet\ Leçon \DOCLECON}{white}{popink}}
\lfoot{\popxbold\fontsize{7.6}{7.6}\selectfont\color{popmuted}\MakeUppercase{Cours signé OnBuch+}\ \textbullet\ {\popsemi\DOCTITRE}}
\rfoot{\tikz[baseline=-0.6ex]\node[rounded corners=8.5pt,fill=popink,minimum height=17pt,minimum width=17pt,inner xsep=5pt,inner sep=0]{\popxbold\fontsize{8}{8}\selectfont\color{popcream}\thepage\,/\,\pageref*{LastPage}};}

% ============================================================
% Titres
% ============================================================
\newcommand{\chaptitre}[2]{%
  \begin{tcolorbox}[enhanced,colback=poporange,colframe=popink,boxrule=2.6pt,arc=11pt,
      drop shadow=popink,left=15pt,right=15pt,top=12pt,bottom=11pt,before skip=3pt,after skip=18pt]
    \poppill{#2}{popink}{popcream}\par\vspace{9pt}
    {\raggedright\hyphenpenalty=10000\exhyphenpenalty=10000\popdisplay\fontsize{22}{24}\selectfont\color{popcream}\MakeUppercase{#1}\par}
  \end{tcolorbox}
}
% Section numérotée : pastille orange avec le numéro + titre Archivo
\newcounter{popsec}
\newcommand{\coursec}[1]{%
  \stepcounter{popsec}\setcounter{popsub}{0}%
  \par\vspace{16pt}\noindent
  \begin{minipage}{\linewidth}
  \tikz[baseline=-0.7ex]\node[draw=popink,line width=1.6pt,fill=poporange,rounded corners=4pt,minimum size=22pt,inner sep=0]{\popdisplay\fontsize{12}{12}\selectfont\color{popcream}\thepopsec};%
  \hspace{8pt}{\popdisplay\fontsize{15}{17}\selectfont\color{popink}\MakeUppercase{#1}}\par
  \vspace{4pt}\noindent{\color{poporange}\rule{36pt}{3.6pt}}
  \end{minipage}\par\nopagebreak\vspace{8pt}
}
\newcounter{popsub}
\newcommand{\courssub}[1]{%
  \stepcounter{popsub}%
  \par\vspace{9pt}\noindent{\popxbold\fontsize{11.5}{13}\selectfont\color{popink}{\color{poporange}\thepopsec.\thepopsub}\hspace{6pt}#1}\par\nopagebreak\vspace{4pt}
}

% ============================================================
% Boîtes pédagogiques — même recette : fond blanc, bordure épaisse colorée,
% ombre dure encre, badge plein. Titre optionnel [#1].
% ============================================================
\tcbset{popbase/.style={breakable,enhanced,colback=white,boxrule=2.3pt,arc=9pt,
  shadow={3pt}{-3pt}{0pt}{popink},left=14pt,right=14pt,top=11pt,bottom=11pt,before skip=11pt,after skip=15pt,
  fontupper=\popsemi\fontsize{10.6}{14.5}\selectfont\color{popink},
  fontlower=\popsemi\fontsize{10.6}{14.5}\selectfont\color{popink}}}
% Coche dessinée (le glyphe ✓ n'existe pas dans les polices du document)
\newcommand{\popcheck}[1]{\tikz[baseline=-0.5ex]\draw[#1,line width=1.6pt,line cap=round,line join=round](-0.13,0.0)--(-0.04,-0.09)--(0.13,0.1);}
\providecommand{\coche}{\popcheck{popgreen}}
\providecommand{\resultat}[1]{\par\smallskip\noindent\fcolorbox{popgreen}{popgreenL}{\parbox{\dimexpr\linewidth-2\fboxsep-2\fboxrule\relax}{\textbf{Résultat.} #1}}\par\smallskip}
\newcommand{\popboxhead}[3]{\poppill{#1}{#2}{white}\ifx\relax#3\relax\else\hspace{6pt}{\popxbold\fontsize{10.6}{12}\selectfont\color{popink}#3}\fi\par\vspace{7pt}}

\newtcolorbox{aretenir}[1][]{popbase,colframe=popgreen,before upper={\popboxhead{À retenir}{popgreen}{#1}}}
% Texte en espagnol (document de lecture, dialogue, poème, extrait) : cadre
% d'étude. Un titre optionnel (source, auteur) s'affiche dans l'en-tête.
\newtcolorbox{textebox}[1][]{popbase,colframe=popblue,before upper={\popboxhead{Texto}{popblue}{#1}\begin{otherlanguage}{spanish}},after upper={\end{otherlanguage}}}
% Marques ✓ / ✗ (forme correcte / fautive) souvent utilisées par les modèles
\providecommand{\cross}{\textcolor{poppink}{\ensuremath{\times}}}
\providecommand{\xmark}{\cross}\providecommand{\crossmark}{\cross}\providecommand{\cmark}{\ensuremath{\checkmark}}
% Ligne de réponse / justification à compléter (inventée par les modèles)
\providecommand{\justification}[1]{\par\noindent\textit{Justification~:} \dotfill\par}
% \textipa (package tipa non chargé) : la police ne couvre pas l'API -> simple passage
\providecommand{\textipa}[1]{#1}
% Mot ou phrase en espagnol à l'intérieur d'un paragraphe français
\newcommand{\esp}[1]{\foreignlanguage{spanish}{\emph{#1}}}
\newtcolorbox{exemplebox}[1][]{popbase,colframe=poporange,before upper={\popboxhead{Exemple}{poporange}{#1}}}
\newtcolorbox{definition}[1][]{popbase,colframe=popblue,before upper={\popboxhead{Définition}{popblue}{#1}}}
\newtcolorbox{propriete}[1][]{popbase,colframe=poppurple,before upper={\popboxhead{Propriété}{poppurple}{#1}}}
\newtcolorbox{methode}[1][]{popbase,colframe=popgold,before upper={\popboxhead{Méthode}{popgold}{#1}}}
\newtcolorbox{attention}[1][]{popbase,colframe=poppink,before upper={\popboxhead{Attention}{poppink}{#1}}}
\newtcolorbox{experience}[1][]{popbase,colframe=popblue,colback=popblueL,before upper={\popboxhead{Expérience}{popblue}{#1}}}
% « Le savais-tu ? » : carte violette pleine (contexte, culture, Cameroun)
\newtcolorbox{savaistu}[1][]{popbase,colback=poppurple,colframe=popink,
  fontupper=\popsemi\fontsize{10.4}{14.2}\selectfont\color{white},
  before upper={\poppill{Le savais-tu\,?}{white}{poppurple}\ifx\relax#1\relax\else\hspace{6pt}{\popxbold\color{white}#1}\fi\par\vspace{7pt}}}
% Exercice résolu : énoncé en haut, \tcblower puis la solution sur fond crème
\newcounter{popexo}\newcounter{popexr}
\newtcolorbox{exoresolu}[1][]{popbase,colframe=poporange,colbacklower=popcream,
  segmentation style={popink,line width=1.2pt,dashed},
  before upper={\stepcounter{popexr}\popboxhead{Exercice résolu \thepopexr}{poporange}{#1}},
  before lower={\poppill{Solution}{popink}{popcream}\par\vspace{6pt}}}
% Exercice (non résolu en place) — la correction est dans la partie Corrigés
\newtcolorbox{exercice}[1][]{popbase,colframe=popink,
  before upper={\stepcounter{popexo}\popboxhead{Exercice \thepopexo}{popink}{#1}}}
% Corrigé
\newtcolorbox{corrige}[1][]{popbase,colframe=popgreen,colback=popgreenL,
  before upper={\popboxhead{Corrigé}{popgreen}{#1}}}
% Objectifs / prérequis / bilan
\newtcolorbox{objectifs}{popbase,colframe=popink,colback=poporangeL,
  before upper={\popboxhead{Objectifs de la leçon}{popink}{}}}
\newtcolorbox{prerequis}{popbase,colframe=popink,colback=popblueL,
  before upper={\popboxhead{Prérequis}{popblue}{}}}
\newtcolorbox{bilan}{popbase,colframe=popink,colback=white,boxrule=2.8pt,
  before upper={\popboxhead{Fiche bilan}{poporange}{L'essentiel à connaître pour l'examen}}}
% Figure encadrée avec légende
\newtcolorbox{popfigure}[1][]{enhanced,breakable=false,colback=white,colframe=popline,boxrule=1.4pt,arc=8pt,
  left=10pt,right=10pt,top=10pt,bottom=8pt,before skip=10pt,after skip=12pt,halign=center,
  lower separated=false,halign lower=center,
  fontlower=\popsemi\fontsize{9}{11.5}\selectfont\color{popmuted},
  #1}
\newcommand{\legende}[1]{\par\vspace{6pt}{\popsemi\fontsize{9}{11.5}\selectfont\color{popmuted}#1}}

% ============================================================
% Signature OnBuch+
% ============================================================
\newcommand{\poplogoinline}[1]{{\poplogo\fontsize{#1}{#1}\selectfont\color{popink}OnBuch\color{poporange}+}}
\newcommand{\signatureonbuch}{%
  \begin{tcolorbox}[enhanced,colback=popink,colframe=popink,boxrule=2.2pt,arc=10pt,
      drop shadow=poporange,left=14pt,right=14pt,top=10pt,bottom=10pt,before skip=0pt,after skip=0pt]
    \tikz[baseline=-0.7ex]\node[circle,fill=popgreen,draw=popcream,line width=1.3pt,minimum size=22pt,inner sep=0]{\popcheck{white}};%
    \hspace{9pt}\begin{minipage}[c]{0.62\linewidth}
      {\popxbold\fontsize{12}{13}\selectfont\color{popcream}Signé\ \textbullet\ {\poplogo OnBuch\color{poporange}+}}\par\vspace{2pt}
      {\raggedright\popsemi\fontsize{8}{10.5}\selectfont\color{popline}\MakeUppercase{Équipe pédagogique OnBuch+}\\\MakeUppercase{Conforme au programme officiel MINESEC}\par}
    \end{minipage}\hfill
    {\poplogo\fontsize{22}{22}\selectfont\color{popcream}OnBuch\color{poporange}+}
  \end{tcolorbox}%
}

% ============================================================
% Page de garde
% #1 titre · #2 matière · #3 niveau · #4 module · #5 n° leçon · #6 durée ·
% #7 description / objectifs (texte ou itemize)
% ============================================================
\newcommand{\pagegarde}[7]{%
  \thispagestyle{empty}
  \newgeometry{a4paper,margin=1.7cm,top=1.6cm,bottom=1.4cm}%
  \begin{tikzpicture}[remember picture,overlay]
    \fill[poporange!18] ([xshift=-3.2cm,yshift=-3.6cm]current page.north east) circle (4.6cm);
    \fill[poppurple!14] ([xshift=2.4cm,yshift=7.6cm]current page.south west) circle (3.8cm);
  \end{tikzpicture}%
  {\poplogo\fontsize{30}{30}\selectfont\color{popink}OnBuch\color{poporange}+}\hfill
  \raisebox{0.6ex}{\poppill{Cours signé \textbullet\ Premium}{popink}{popcream}}\par
  \vspace{1pt}\popsquiggle{poppurple}\par
  \vspace{8pt}
  \begin{tcolorbox}[enhanced,colback=poporange,colframe=popink,boxrule=2.8pt,arc=13pt,
      drop shadow=popink,left=18pt,right=18pt,top=12pt,bottom=13pt,after skip=10pt]
    \poppill{#2 \textbullet\ #3}{popink}{popcream}\hspace{5pt}\poppill{#4}{white}{popink}\par\vspace{8pt}
    {\popdisplay\fontsize{14}{14}\selectfont\color{popink}LEÇON #5}\par\vspace{4pt}
    {\raggedright\hyphenpenalty=10000\exhyphenpenalty=10000\popdisplay\fontsize{25}{27}\selectfont\color{popcream}\MakeUppercase{#1}\par}
    \vspace{8pt}
    {\popxbold\fontsize{9.5}{9.5}\selectfont\color{popink}\MakeUppercase{Durée conseillée : #6}}
  \end{tcolorbox}
  \begin{tcolorbox}[enhanced,colback=white,colframe=popink,boxrule=2.2pt,arc=10pt,
      drop shadow=popink,left=16pt,right=16pt,top=10pt,bottom=10pt,after skip=8pt]
    \poppill{Dans cette leçon}{poppurple}{white}\par\vspace{5pt}
    \popsemi\fontsize{9.3}{12.6}\selectfont\color{popink}#7
  \end{tcolorbox}
  \noindent\begin{minipage}[t]{0.487\linewidth}
    \begin{tcolorbox}[enhanced,colback=popgreen,colframe=popink,boxrule=2.2pt,arc=10pt,
        drop shadow=popink,left=11pt,right=11pt,top=10pt,bottom=10pt,height=3.1cm,valign=center]
      \tcbox[colback=white,colframe=popink,boxrule=1.3pt,arc=5pt,boxsep=0pt,nobeforeafter,
        left=3pt,right=3pt,top=3pt,bottom=3pt,valign=center]{\qrcode[height=1.9cm]{https://play.google.com/store/apps/details?id=com.luvvix.onbuch&hl=fr}}%
      \hspace{8pt}\begin{minipage}[c]{0.5\linewidth}
        {\popdisplay\fontsize{11}{12.5}\selectfont\color{white}\MakeUppercase{Télécharge OnBuch}}\par\vspace{4pt}
        {\raggedright\popsemi\fontsize{8.4}{11}\selectfont\color{white}Gratuit \textbullet\ Google Play\\Tuteur IA, annales, concours, résultats\par}
      \end{minipage}
    \end{tcolorbox}
  \end{minipage}\hfill
  \begin{minipage}[t]{0.487\linewidth}
    \begin{tcolorbox}[enhanced,colback=poppurple,colframe=popink,boxrule=2.2pt,arc=10pt,
        drop shadow=popink,left=11pt,right=11pt,top=10pt,bottom=10pt,height=3.1cm,valign=center]
      \tikz[baseline=-0.7ex]\node[circle,fill=white,draw=popink,line width=1.3pt,minimum size=1.6cm,inner sep=0]{\popdisplay\fontsize{24}{24}\selectfont\color{poppurple}+};%
      \hspace{8pt}\begin{minipage}[c]{0.6\linewidth}
        {\popdisplay\fontsize{11}{12.5}\selectfont\color{white}\MakeUppercase{Passe à OnBuch+}}\par\vspace{4pt}
        {\raggedright\popsemi\fontsize{8.4}{11}\selectfont\color{white}Tous les cours signés, Tuteur IA illimité, fiches PDF — onglet OnBuch+ de l'app\par}
      \end{minipage}
    \end{tcolorbox}
  \end{minipage}
  \vspace{8pt}
  \popcontenu
  \vfill
  \signatureonbuch
  \restoregeometry
  \newpage
}

% Rangée « Ce cours contient » (page de garde)
\newcommand{\poptile}[3]{% #1 couleur #2 gros texte #3 légende
  \begin{tcolorbox}[enhanced,colback=#1,colframe=popink,boxrule=2pt,arc=9pt,shadow={2.5pt}{-2.5pt}{0pt}{popink},
      left=6pt,right=6pt,top=9pt,bottom=9pt,halign=center,valign=center,height=1.75cm,width=\linewidth,nobeforeafter]
    {\popdisplay\fontsize{12.5}{13}\selectfont\color{white}\MakeUppercase{#2}}\par\vspace{3pt}
    {\centering\popsemi\fontsize{7.8}{9.5}\selectfont\color{white}#3\par}
  \end{tcolorbox}}
\newcommand{\popcontenu}{%
  \par\noindent{\popxboldls\fontsize{7.5}{7.5}\selectfont\color{popmuted}CE COURS SIGNÉ CONTIENT}\par\vspace{7pt}
  \noindent\begin{minipage}[t]{0.235\linewidth}\poptile{popblue}{Cours}{complet, illustré, pas à pas}\end{minipage}\hfill
  \begin{minipage}[t]{0.235\linewidth}\poptile{poporange}{Méthodes}{et exercices résolus}\end{minipage}\hfill
  \begin{minipage}[t]{0.235\linewidth}\poptile{poppink}{Exercices}{type examen + corrigés}\end{minipage}\hfill
  \begin{minipage}[t]{0.235\linewidth}\poptile{popgreen}{Fiche bilan}{l'essentiel à retenir}\end{minipage}\par}

% Sommaire
\newtcolorbox{sommaire}{enhanced,colback=white,colframe=popink,boxrule=2.2pt,arc=10pt,
  drop shadow=popink,left=18pt,right=18pt,top=13pt,bottom=13pt,before skip=6pt,after skip=20pt,
  before upper={\poppill{Sommaire}{poppurple}{white}\par\vspace{9pt}\popsemi\fontsize{10.6}{14.5}\selectfont\color{popink}}}

% Signature de fin de cours — poussée tout en bas de la dernière page
\newcommand{\findecours}{%
  \par\vfill\nopagebreak
  \begin{center}\popsquiggle{poporange}\end{center}\vspace{4pt}
  \signatureonbuch
}
\AtBeginDocument{\def\llbracket{\mathopen{[\mkern-3.5mu[}}\def\rrbracket{\mathclose{]\mkern-3.5mu]}}} % intervalles d'entiers (glyphe ⟦ absent de la police)
% Glyphes absents de Fira Math : redéfinis à partir de glyphes disponibles
\AtBeginDocument{%
  \def\setminus{\mathbin{\backslash}}\let\smallsetminus\setminus
  \def\vdots{\mathord{\vbox{\baselineskip3.2pt\lineskiplimit0pt\kern4pt\hbox{.}\hbox{.}\hbox{.}}}}%
  \def\ddots{\mathinner{\mkern1mu\raise6.5pt\hbox{.}\mkern2mu\raise3.6pt\hbox{.}\mkern2mu\raise0.7pt\hbox{.}\mkern1mu}}%
  \def\triangle{\mathord{\scalebox{0.9}{$\symup{\Delta}$}}}%
  \def\nabla{\mathord{\raisebox{\depth}{\scalebox{1}[-1]{$\symup{\Delta}$}}}}%
  \def\checkmark{\popcheck{popdark}}%
  \def\star{\mathbin{\ast}}\def\bigstar{\mathord{\ast}}%
  \def\diamond{\mathbin{\scalebox{0.6}{$\Diamond$}}}%
  \def\bigcup{\mathop{\mathchoice{\raisebox{-0.3em}{\scalebox{1.7}{$\cup$}}}{\scalebox{1.25}{$\cup$}}{\cup}{\cup}}\displaylimits}%
  \def\bigcap{\mathop{\mathchoice{\raisebox{-0.3em}{\scalebox{1.7}{$\cap$}}}{\scalebox{1.25}{$\cap$}}{\cap}{\cap}}\displaylimits}%
  \def\lesseqqgtr{\mathrel{\lessgtr}}\def\bigoplus{\mathop{\oplus}\displaylimits}%
}
\providecommand{\ssi}{si et seulement si} % abréviation fréquente
\AtBeginDocument{\providecommand{\wideparen}{\overparen}} % arc de cercle

\begin{document}

\pagegarde{Diálogo familiar, debate sobre la violencia de género et traduction : SER y ESTAR}{Espagnol}{Tle A}{Módulo 1 : Vie familiale et intégration sociale}{E02}{2 h}{Leçon mixte du Module 1 (Vie familiale et intégration sociale) : à partir d'un dialogue familial et d'un débat sur la violence de genre, les élèves apprennent à communiquer oralement sur la vie familiale et maîtrisent l'emploi de SER et ESTAR pour traduire « être » du français vers l'espagnol.
\vspace{4pt}
\begin{multicols}{2}\small
\begin{itemize}
\item À la fin de cette leçon, je sais comprendre un dialogue familial en espagnol et en dégager les informations essentielles.
\item À la fin de cette leçon, je sais écouter activement, relancer et reformuler dans une conversation familiale.
\item À la fin de cette leçon, je sais exprimer mon opinion et argumenter dans un débat sur les violences faites aux femmes.
\item À la fin de cette leçon, je sais utiliser le lexique de la famille, des relations et de la violence de genre.
\item À la fin de cette leçon, je sais conjuguer SER et ESTAR au présent et distinguer leurs emplois fondamentaux.
\item À la fin de cette leçon, je sais expliquer les changements de sens avec SER/ESTAR + adjectif (ser listo / estar listo...).
\end{itemize}\end{multicols}}

\begin{sommaire}
\begin{multicols}{2}
\begin{enumerate}
\item Texte support : un dialogue familial
\item Lexique thématique : famille et violence de genre
\item Méthode du commentaire et commentaire modèle
\item Grammaire : SER et ESTAR — emplois fondamentaux
\item SER/ESTAR + adjectif : changements de sens et expressions figées
\item Traduction guidée : thème et version
\item Activité d'intégration
\item Méthodes et exercices résolus
\item Exercices
\item Corrigés des exercices
\item Fiche bilan
\end{enumerate}
\end{multicols}
\end{sommaire}

\begin{objectifs}
\begin{itemize}
\item À la fin de cette leçon, je sais comprendre un dialogue familial en espagnol et en dégager les informations essentielles.
\item À la fin de cette leçon, je sais écouter activement, relancer et reformuler dans une conversation familiale.
\item À la fin de cette leçon, je sais exprimer mon opinion et argumenter dans un débat sur les violences faites aux femmes.
\item À la fin de cette leçon, je sais utiliser le lexique de la famille, des relations et de la violence de genre.
\item À la fin de cette leçon, je sais conjuguer SER et ESTAR au présent et distinguer leurs emplois fondamentaux.
\item À la fin de cette leçon, je sais expliquer les changements de sens avec SER/ESTAR + adjectif (ser listo / estar listo...).
\item À la fin de cette leçon, je sais utiliser les expressions figées avec ESTAR (estar de acuerdo, estar de moda...).
\item À la fin de cette leçon, je sais traduire du français vers l'espagnol des phrases contenant le verbe « être » sans erreur de choix.
\end{itemize}
\end{objectifs}

\begin{prerequis}
\begin{itemize}
\item Le présent de l'indicatif des verbes réguliers et irréguliers courants (Première)
\item Le lexique de base de la famille et de la description physique et morale (Seconde/Première)
\item Les structures d'expression de l'opinion vues en Première (creo que, me parece que)
\item L'accord de l'adjectif en genre et en nombre
\item La notion de thème et de version (traduction française ↔ espagnole)
\end{itemize}
\end{prerequis}

\newpage

\coursec{Texte support : un dialogue familial}

C'est dimanche à Yaoundé. Chez les Ngo Bella, comme chaque semaine, toute la famille est réunie autour de la table : le poulet DG fume encore, les enfants se disputent la dernière part d'alloco, et la télévision diffuse les informations en fond sonore. Aujourd'hui, la conversation s'anime : les études de la cadette, le mariage du cousin, mais aussi un sujet grave évoqué au journal télévisé — les violences faites aux femmes, un fléau dénoncé aussi bien au Cameroun qu'en Espagne ou au Mexique. Chacun prend la parole, écoute, relance, reformule, donne son avis.

Voilà exactement la compétence que cette leçon va construire : comment \cle{écouter}, \cle{relancer} et \cle{reformuler} dans un dialogue familial en espagnol ? Comment donner son avis et argumenter dans un débat ? Et comment traduire correctement le verbe « être », qui se dit tantôt \esp{ser}, tantôt \esp{estar} ? Cette leçon vous donne toutes les clés.

\courssub{Lecture du texte : « Una discusión en familia »}

Avant toute explication, lisons le dialogue. Lisez-le deux fois : une première fois pour le sens général, une deuxième fois en soulignant les opinions de chaque personnage.

\begin{textebox}[Una discusión en familia]
Es domingo en Yaoundé. La familia Ngo Bella está reunida alrededor de la mesa: los abuelos, los padres y los tres hijos. La televisión está encendida y todos escuchan las noticias.

---Mamá, ¿has oído las noticias? Una vecina ha denunciado a su marido por malos tratos ---dice Amina, la hija mayor.

---Es terrible, hija. La violencia de género es un problema de toda la sociedad ---responde la madre---. Ya veo que este tema te preocupa mucho.

---Claro que sí. En mi facultad hablamos mucho de eso. Yo creo que hay que educar a los chicos desde pequeños ---dice el padre, que es profesor.

---¿De verdad piensas eso, papá? ---pregunta Amina, sorprendida.

---Sí, hija. Es decir, la educación empieza en casa y en la escuela.

---¿Y tú qué piensas, abuelo? ---pregunta la madre, que quiere incluir a todos en la conversación.

---Pienso que antes no se hablaba de esto, pero me alegro de que hoy las mujeres denuncien ---contesta el abuelo---. En otras palabras, los tiempos cambian, y cambian para bien.

---O sea, ¿estás de acuerdo con nosotros? ---reformula Amina.

---Totalmente de acuerdo. Además, tu suegra me dijo ayer que en su barrio hay una asociación que apoya a las mujeres maltratadas.

---¡Qué bien! Entonces, ¿por qué no hablamos con la nuera del tío Etienne? Ella es abogada y puede ayudarnos.

---Buena idea ---dice el padre---. Pero ahora, hablemos también de tus estudios: ¿sigues queriendo estudiar derecho?

---Sí, papá. Quiero ser abogada para defender a las víctimas.

---Estamos muy orgullosos de ti, hija ---concluye la madre, sonriendo.
\end{textebox}

\begin{definition}[Glossaire du texte]
\begin{itemize}
\item \esp{los suegros} : les beaux-parents ; \esp{el suegro} : le beau-père ; \esp{la suegra} : la belle-mère.
\item \esp{la nuera} : la belle-fille ; \esp{el yerno} : le gendre.
\item \esp{maltratar} : maltraiter ; \esp{los malos tratos} : les mauvais traitements.
\item \esp{la pareja} : le couple, le/la partenaire.
\item \esp{denunciar} : dénoncer, porter plainte contre.
\item \esp{apoyar} : soutenir.
\item \esp{discutir} : discuter, mais aussi se disputer (voir la boîte Attention).
\item \esp{estar de acuerdo} : être d'accord.
\item \esp{alegrarse de que + subjonctif} : se réjouir que.
\end{itemize}
\end{definition}

\begin{savaistu}[La famille, institution centrale]
Dans le monde hispanophone comme au Cameroun, le repas du dimanche en famille est un moment sacré de conversation. En Espagne, la \esp{sobremesa} désigne le temps passé à discuter à table après le repas, parfois des heures ! C'est souvent là que naissent les grands débats familiaux. En Guinée équatoriale, pays hispanophone voisin du Cameroun, les mêmes sujets de société — éducation, mariage, droits des femmes — animent les discussions familiales.
\end{savaistu}

\courssub{Compréhension globale : qui ? où ? de quoi ?}

Avant d'entrer dans le détail, posons-nous les quatre questions fondamentales de tout dialogue.

\begin{methode}[Comprendre un dialogue en quatre étapes]
\begin{enumerate}
\item \textbf{Identifier les personnages et le lieu} : qui parle ? où se déroule la scène ? Ici : la famille Ngo Bella (les grands-parents, les parents, les enfants), à table, un dimanche à Yaoundé.
\item \textbf{Repérer le thème} : de quoi parle-t-on ? Ici : un fait divers sur les violences conjugales, puis les études d'Amina.
\item \textbf{Noter l'opinion de chacun} : chaque personnage a une position. La mère dénonce un problème de société ; le père insiste sur l'éducation ; le grand-père constate le changement des mentalités ; Amina veut devenir avocate.
\item \textbf{Relever les connecteurs et marqueurs d'interaction} : \esp{es decir}, \esp{o sea}, \esp{además}, \esp{entonces}, \esp{¿y tú qué piensas?} structurent l'échange.
\end{enumerate}
\end{methode}

\begin{exemplebox}[Réponses à la compréhension globale]
\begin{itemize}
\item \textbf{¿Quiénes hablan?} \esp{Hablan los miembros de la familia Ngo Bella: Amina, sus padres y su abuelo.} (Les membres de la famille Ngo Bella : Amina, ses parents et son grand-père.)
\item \textbf{¿Dónde están?} \esp{Están en casa, alrededor de la mesa, un domingo.} (Ils sont à la maison, autour de la table, un dimanche.)
\item \textbf{¿De qué hablan?} \esp{Hablan de la violencia de género y de los estudios de Amina.} (Ils parlent des violences faites aux femmes et des études d'Amina.)
\item \textbf{¿Cuál es el conflicto?} Il n'y a pas de dispute ouverte, mais un débat d'idées : comment lutter contre les violences conjugales ? Par l'éducation (le père), par la dénonciation (le grand-père), par le droit (Amina).
\end{itemize}
\end{exemplebox}

\courssub{Compréhension fine : les opinions de chaque personnage}

\begin{exercice}[Questions de détail]
Répondez en espagnol, par des phrases complètes.
\begin{enumerate}
\item ¿Qué ha hecho la vecina de la familia?
\item Según la madre, ¿de quién es el problema de la violencia de género?
\item ¿Qué solución propone el padre?
\item ¿Qué piensa el abuelo de la situación actual?
\item ¿Qué profesión quiere ejercer Amina y por qué?
\end{enumerate}
\end{exercice}

\begin{corrige}[Exercice « Questions de détail »]
\begin{enumerate}
\item \esp{La vecina ha denunciado a su marido por malos tratos.} (La voisine a porté plainte contre son mari pour mauvais traitements.)
\item \esp{Según la madre, es un problema de toda la sociedad.} (Selon la mère, c'est un problème de toute la société.)
\item \esp{El padre propone educar a los chicos desde pequeños.} (Le père propose d'éduquer les garçons dès leur plus jeune âge.)
\item \esp{El abuelo piensa que antes no se hablaba de esto, pero se alegra de que hoy las mujeres denuncien.} (Le grand-père pense qu'autrefois on n'en parlait pas, mais il se réjouit qu'aujourd'hui les femmes portent plainte.)
\item \esp{Amina quiere ser abogada para defender a las víctimas.} (Amina veut devenir avocate pour défendre les victimes.)
\end{enumerate}
\end{corrige}

\courssub{Les marqueurs d'écoute active et de relance}

Un dialogue n'est pas une succession de monologues : chaque locuteur montre qu'il écoute et invite l'autre à parler. Relevons ces marqueurs dans le texte.

\begin{center}
\begin{tabularx}{\linewidth}{|X|X|X|}
\hline
\rowcolor{popblueL} Marqueur & Fonction & Traduction \\
\hline
\esp{¿De verdad?} & réaction, écoute active & « Vraiment ? » \\
\hline
\esp{Ya veo} & montrer qu'on comprend & « Je vois » \\
\hline
\esp{Claro que sí} & acquiescer avec force & « Bien sûr que oui » \\
\hline
\esp{¿Y tú qué piensas?} & relancer, donner la parole & « Et toi, qu'en penses-tu ? » \\
\hline
\esp{Totalmente de acuerdo} & exprimer l'accord & « Tout à fait d'accord » \\
\hline
\esp{¡Qué bien!} & exprimer l'enthousiasme & « C'est génial ! » \\
\hline
\end{tabularx}
\end{center}

\begin{exemplebox}[Exemples traduits]
\esp{¿Y tú qué piensas?} = « Et toi, qu'en penses-tu ? »

\esp{Me alegro de que denuncien.} = « Je me réjouis qu'elles portent plainte. » (Attention : après \esp{alegrarse de que}, l'espagnol exige le \cle{subjonctif} : \esp{denuncien}, et non l'indicatif.)

\esp{Ya veo que este tema te preocupa.} = « Je vois que ce sujet te préoccupe. »
\end{exemplebox}

\courssub{La reformulation : dire autrement}

\begin{definition}[La reformulation]
La \cle{reformulation} consiste à redire avec d'autres mots ce que vient de dire l'interlocuteur, pour vérifier qu'on a bien compris et pour faire avancer l'échange. Les marqueurs typiques sont : \esp{es decir} (c'est-à-dire), \esp{o sea} (c'est-à-dire, familier), \esp{en otras palabras} (en d'autres termes), \esp{¿quieres decir que...?} (tu veux dire que...?).
\end{definition}

Dans le dialogue, la reformulation fonctionne en boucle :
\begin{itemize}
\item Le père : \esp{Hay que educar a los chicos desde pequeños. Es decir, la educación empieza en casa.} — il précise sa propre idée.
\item Amina : \esp{O sea, ¿estás de acuerdo con nosotros?} — elle vérifie qu'elle a compris le grand-père.
\item Le grand-père : \esp{En otras palabras, los tiempos cambian.} — il résume sa pensée.
\end{itemize}

\begin{exoresolu}[Reformuler les propos d'un personnage]
Énoncé : reformulez en espagnol, avec \esp{en otras palabras}, la phrase de la mère : \esp{La violencia de género es un problema de toda la sociedad.}
\tcblower
Solution : \esp{En otras palabras, la violencia contra las mujeres no es un asunto privado, sino un problema que concierne a todos.}

Traduction : « En d'autres termes, la violence contre les femmes n'est pas une affaire privée, mais un problème qui concerne tout le monde. »

Méthode : on reprend l'idée centrale (problème de toute la société) et on l'exprime avec un vocabulaire différent (\esp{asunto privado}, \esp{concierne a todos}), sans changer le sens.
\end{exoresolu}

\begin{attention}[Faux ami : « discutir »]
\esp{Discutir} ne signifie pas seulement « discuter » au sens d'échanger des idées : très souvent, il veut dire \textbf{se disputer}. De même, \esp{la discusión} peut être « le débat » ou « la dispute » selon le contexte.

\esp{Mis padres están discutiendo otra vez.} = « Mes parents sont encore en train de se disputer. »

\esp{Tuvimos una discusión interesante en clase.} = « Nous avons eu un débat intéressant en classe. »

Pour dire simplement « discuter, parler de », préférez \esp{hablar de} : \esp{Hablamos de tus estudios.} = « Nous parlons de tes études. »
\end{attention}

\courssub{La famille Ngo Bella : arbre généalogique et vocabulaire}

Pour bien situer les personnages, voici l'arbre généalogique de la famille, avec les étiquettes espagnoles à connaître.

\begin{popfigure}
\begin{tikzpicture}[
level distance=1.6cm,
level 1/.style={sibling distance=4.5cm},
level 2/.style={sibling distance=2.6cm},
every node/.style={draw=popblue, thick, rounded corners, fill=popblueL, align=center, font=\small},
edge from parent/.style={draw=popdark, thick}]
\node[align=center]{los abuelos\\ \esp{el abuelo y la abuela}}
  child { node[align=center]{los padres\\ \esp{el padre y la madre}}
    child { node[align=center]{\esp{Amina}\\ la hija mayor}
      edge from parent node[draw=none, fill=none, left, font=\scriptsize]{} }
    child { node[align=center]{el hijo\\ \esp{el hermano}} }
    child { node[align=center]{la hija menor\\ \esp{la cadette}} }
  };
\end{tikzpicture}
\legende{Arbre généalogique de la famille Ngo Bella : \esp{los abuelos} (les grands-parents), \esp{los padres} (les parents), \esp{los hijos} (les enfants).}
\end{popfigure}

\begin{center}
\begin{tabularx}{\linewidth}{|X|X|X|}
\hline
\rowcolor{popblueL} Français & Español & Remarque \\
\hline
le beau-père / la belle-mère & \esp{el suegro / la suegra} & aussi « les beaux-parents » : \esp{los suegros} \\
\hline
la belle-fille & \esp{la nuera} & la femme du fils \\
\hline
le gendre & \esp{el yerno} & le mari de la fille \\
\hline
le couple, le/la partenaire & \esp{la pareja} & toujours féminin, même pour un homme \\
\hline
maltraiter & \esp{maltratar} & \esp{los malos tratos} : les mauvais traitements \\
\hline
dénoncer, porter plainte & \esp{denunciar} & \esp{denunciar a alguien} \\
\hline
soutenir & \esp{apoyar} & \esp{apoyar a las víctimas} \\
\hline
\end{tabularx}
\end{center}

\courssub{Synthèse : les trois gestes du bon dialogue}

\begin{popfigure}
\begin{tikzpicture}[mindmap, grow cyclic,
every node/.style={concept, align=center, font=\small},
root concept/.append style={concept color=popblue, text=white, font=\bfseries},
level 1/.append style={level distance=3.4cm, sibling angle=120},
level 1 concept/.append style={concept color=poporange, text=popdark}]
\node {El diálogo familiar}
  child { node[align=center]{escuchar\\ \esp{ya veo}\\ \esp{¿de verdad?}} }
  child { node[align=center]{relanzar\\ \esp{¿y tú qué piensas?}\\ \esp{¿estás de acuerdo?}} }
  child { node[align=center]{reformular\\ \esp{es decir}\\ \esp{o sea}\\ \esp{en otras palabras}} };
\end{tikzpicture}
\legende{Carte mentale : les trois compétences du dialogue familial — écouter, relancer, reformuler.}
\end{popfigure}

\begin{aretenir}[À retenir de cette section]
\begin{itemize}
\item Un dialogue familial repose sur trois gestes : \cle{écouter} (\esp{ya veo}, \esp{¿de verdad?}), \cle{relancer} (\esp{¿y tú qué piensas?}), \cle{reformuler} (\esp{es decir}, \esp{o sea}, \esp{en otras palabras}).
\item Pour comprendre un dialogue : identifier les personnages et le lieu, repérer le thème, noter l'opinion de chacun, relever les marqueurs d'interaction.
\item Le vocabulaire de la famille politique (\esp{suegros}, \esp{nuera}, \esp{yerno}) et du débat social (\esp{maltratar}, \esp{denunciar}, \esp{apoyar}) est indispensable pour le commentaire et le débat.
\item Attention au faux ami \esp{discutir} : souvent « se disputer ».
\end{itemize}
\end{aretenir}

\begin{experience}[Jeu de rôle : à la table des Ngo Bella]
Par groupes de quatre, rejouez la scène : un élève est Amina, un autre la mère, un autre le père, un autre le grand-père. Remplacez le fait divers par une actualité de votre choix (le mariage du cousin, les examens, un match de football des Lions Indomptables). Obligation : chaque personnage doit utiliser au moins une fois un marqueur d'écoute (\esp{ya veo}, \esp{¿de verdad?}), une relance (\esp{¿y tú qué piensas?}) et une reformulation (\esp{o sea}...). La classe écoute et coche les marqueurs entendus.
\end{experience}

\coursec{Lexique thématique : famille et violence de genre}

Après avoir analysé le dialogue familial de la section précédente, où l'on a pu observer les interactions quotidiennes entre parents et enfants, nous devons désormais outiller l'élève pour qu'il puisse \textbf{nommer} les acteurs, \textbf{décrire} les relations, \textbf{dénoncer} les dysfonctionnements graves et \textbf{argumenter} lors d'un échange oral ou écrit. Ce lexique est la « boîte à outils » indispensable pour réussir les épreuves de compréhension et d'expression du baccalauréat.

\courssub{Champ 1 : La famille et les liens (La familia y los vínculos)}

Observons d'abord comment le dialogue introduisait les membres de la famille : \esp{« Mi padre y mi madre discuten »} (« Mon père et ma mère discutent »), \esp{« Mis hermanos gemelos »} (« Mes frères jumeaux »). L'espagnol regroupe souvent par le masculin pluriel (genre non marqué).

\begin{center}
\begin{tabularx}{\linewidth}{|X|X|X|}
\hline
\rowcolor{popblueL}
\textbf{Español} & \textbf{Français} & \textbf{Remarque / Contexte camerounais} \\
\hline
los padres & les parents (père + mère) & \textbf{Attention :} ne signifie pas « les pères ». Pour « les pères », dire \esp{los progenitores} o \esp{los padres varones}. \\
la madre / el padre & la mère / le père & \\
los abuelos & les grands-parents & \esp{el abuelo}, \esp{la abuela}. \\
los hijos & les enfants (fils + filles) & \esp{el hijo}, \esp{la hija}. \\
los hermanos & les frères et sœurs & \esp{el hermano}, \esp{la hermana}. \\
los gemelos & les jumeaux / les jumelles & \esp{gemelo/a} (vrais jumeaux), \esp{mellizo/a} (faux jumeaux). \\
el marido / la esposa & le mari / l'épouse & Aussi \esp{el esposo} / \esp{la esposa}. \esp{La mujer} = la femme (et l'épouse). \\
el hogar & le foyer, la maison & Charge affective forte : \esp{un hogar unido} (un foyer uni). \\
la convivencia & la cohabitation, la vie commune & \esp{Las reglas de convivencia} (règles de vie). \\
la crianza & l'éducation, l'élevage (enfants) & \esp{La crianza de los hijos} (éducation des enfants). \\
los parientes & les parents (au sens large : oncles, cousins) & \textbf{Faux ami :} \esp{pariente} $\neq$ parent (père/mère). \esp{Un familiar} = un membre de la famille. \\
\hline
\end{tabularx}
\end{center}

\begin{exemplebox}[Exemples contextualisés]
\esp{En mi hogar, la convivencia es pacífica.} « Dans mon foyer, la cohabitation est pacifique. » \\
\esp{Mis padres educan a mis hermanos gemelos con mucho cariño.} « Mes parents élèvent mes frères jumeaux avec beaucoup d'affection. » \\
\esp{Tengo muchos parientes en Douala, pero mi familia nuclear vive en Yaoundé.} « J'ai beaucoup de parents (oncles, cousins) à Douala, mais ma famille nucléaire vit à Yaoundé. »
\end{exemplebox}

\begin{attention}[Pièges majeurs pour francophones]
\begin{itemize}
    \item \textbf{\cle{los padres}} = les parents (le couple parental). Jamais « les pères ».
    \item \textbf{\cle{pariente}} = membre de la famille élargie (cousin, oncle). \textbf{\cle{padre/madre}} = parent direct.
    \item \textbf{\cle{familiar}} (adj. ou n.) = de famille / membre de la famille. \esp{Un familiar} = un parent (au sens large). Ne pas confondre avec l'adjectif français « familier » (connu, intime) qui se dit \esp{conocido} o \esp{íntimo}.
    \item \cle{esposo/esposa} vs \cle{marido/mujer} : \esp{esposo} est plus formel/juridique ; \esp{marido/mujer} plus courant à l'oral.
\end{itemize}
\end{attention}

\courssub{Champ 2 : Les relations familiales (Las relaciones familiares)}

Le dialogue montrait des tensions (\esp{« discuten »}) mais aussi du soutien. Voici les verbes clés pour décrire la qualité des liens. \textbf{Attention :} certains sont pronominaux (réciproques), d'autres transitifs directs ou prépositionnels.

\begin{center}
\begin{tabularx}{\linewidth}{|X|X|X|X|}
\hline
\rowcolor{poporangeL}
\textbf{Verbe / Locution} & \textbf{Sens} & \textbf{Type} & \textbf{Structure / Exemple} \\
\hline
llevarse bien/mal con & s'entendre bien/mal avec & Pron. réciproque & \esp{Me llevo bien con mi suegra.} (Je m'entends bien avec ma belle-mère.) \\
apoyarse (en) & se soutenir (mutuellement) & Pron. réciproque & \esp{Los hermanos se apoyan mucho.} (Les frères se soutiennent beaucoup.) \\
confiar en & avoir confiance en & Trans. prépositionnel & \esp{Confío plenamente en mi padre.} (J'ai entièrement confiance en mon père.) \\
pelearse (con) & se disputer, se battre (avec) & Pron. réciproque & \esp{Los niños se pelean por los juguetes.} (Les enfants se disputent pour les jouets.) \\
reconciliarse (con) & se réconcilier (avec) & Pron. réciproque & \esp{Tras la discusión, se reconciliaron.} (Après la dispute, ils se sont réconciliés.) \\
educar & éduquer, élever & Trans. direct & \esp{Los padres educan a sus hijos.} (Les parents éduquent leurs enfants.) \\
respetar & respecter & Trans. direct & \esp{Hay que respetar a los mayores.} (Il faut respecter les aînés.) \\
\hline
\end{tabularx}
\end{center}

\begin{propriete}[Grammaire : Verbes pronominaux réciproques]
Les verbes \esp{llevarse}, \esp{apoyarse}, \esp{pelearse}, \esp{reconciliarse} s'emploient au \textbf{pronominal réciproque} (action mutuelle entre plusieurs sujets).
\begin{itemize}
    \item Forme : \esp{Nos llevamos bien} (Nous nous entendons bien / Nous nous entendons bien \textbf{l'un l'autre}).
    \item Accord du participe passé : Il s'accorde avec le sujet si le pronom \esp{se} est COD (ex: \esp{Se han apoyado} -- ils se sont soutenus), mais \textbf{pas} s'il est COI (ex: \esp{Se han hablado} -- ils se sont parlé) ni s'il fait partie du lexème verbal (ex: \esp{Se han llevado bien}).
    \item \esp{Se pelean} = Ils se battent \textbf{entre eux}. \esp{Pelean a alguien} = Ils battent quelqu'un (transitif direct, non pronominal).
\end{itemize}
\end{propriete}

\begin{exemplebox}[Variété des registres]
\esp{Mi tía y yo nos llevamos como el perro y el gato.} « Ma tante et moi, on s'entend comme chien et chat » (idiome : très mal). \\
\esp{En los momentos difíciles, la familia se apoya.} « Dans les moments difficiles, la famille se soutient. » \\
\esp{No te pelees con tu hermano, reconciliaos.} « Ne te dispute pas avec ton frère, réconciliez-vous. » (Impératif \emph{vosotros} : \esp{no os peleéis, reconciliaos} ; \emph{ustedes} : \esp{no se peleen, reconcíliense}).
\end{exemplebox}

\courssub{Champ 3 : La violence de genre (La violencia de género)}

Ce vocabulaire est technique, juridique et social. Il ne s'agit pas de « disputes conjugales » mais d'un phénomène structurel. En Espagne, le terme légal est \esp{violencia de género} (loi organique 1/2004). Au Cameroun, on parle souvent de \esp{violencias basadas en el género} (VBG) dans les ONG et les textes officiels (MINPROFF).

\begin{definition}[La violencia de género]
La \cle{violencia de género} (ou \esp{violencia machista}) est toute violence exercée contre une personne \textbf{en raison de son sexe ou genre}, qui cause ou peut causer un dommage physique, sexuel, psychologique ou économique. Elle touche majoritairement les femmes par des hommes (partenaires ou ex-partenaires). Ce n'est pas un « conflit de couple », c'est une violation des droits humains.
\end{definition}

\begin{center}
\begin{tabularx}{\linewidth}{|X|X|X|}
\hline
\rowcolor{poppurpleL}
\textbf{Término (Español)} & \textbf{Français} & \textbf{Précision juridique / sociale} \\
\hline
la violencia de género & la violence de genre / conjugale & Terme légal en Espagne. Inclut violence physique, psychologique, sexuelle, économique. \\
los malos tratos & les mauvais traitements & Terme juridique ancien, encore très utilisé. \esp{Delito de malos tratos}. \\
el maltratador / la maltratadora & l'agresseur / l'auteure des violences & \esp{El agresor} est plus général. \esp{Maltratador} implique une relation de pouvoir/dominance. \\
la víctima & la victime & Personne qui subit. On dit \esp{mujer víctima de violencia de género}. \\
el acoso & le harcèlement & \esp{Acoso sexual}, \esp{acoso psicológico}, \esp{ciberacoso} (cyberharcèlement). \\
el abuso & l'abus, l'abus sexuel & \esp{Abuso sexual} (sans pénétration), \esp{agresión sexual} (avec pénétration/violence). \\
la agresión & l'agression & \esp{Agresión física}, \esp{agresión sexual}. \\
el feminicidio / el femicidio & le féminicide & Meurtre d'une femme \textbf{parce qu'elle est femme}. Terme légal dans 18 pays d'AmLat. \\
la orden de protección & l'ordonnance de protection & Mesure judiciaire urgente (éloignement, interdiction de communiquer). \\
\hline
\end{tabularx}
\end{center}

\begin{exemplebox}[Phrases authentiques style presse/rapport]
\esp{La ley integral contra la violencia de género protege a las víctimas.} « La loi intégrale contre la violence de genre protège les victimes. » \\
\esp{El maltratador fue condenado a tres años de cárcel por malos tratos habituales.} « L'agresseur a été condamné à trois ans de prison pour mauvais traitements habituels. » \\
\esp{El acoso psicológico es invisible pero destruye la autoestima.} « Le harcèlement psychologique est invisible mais détruit l'estime de soi. »
\end{exemplebox}

\courssub{Champ 4 : Agir et réagir : verbes d'action civique (Actuar y reaccionar)}

Face à la violence, le lexique passe de la description à l'action citoyenne et institutionnelle.

\begin{center}
\begin{tabularx}{\linewidth}{|X|X|X|}
\hline
\rowcolor{popgreenL}
\textbf{Verbe} & \textbf{Français} & \textbf{Construction / Exemple} \\
\hline
denunciar & dénoncer, porter plainte & \esp{Denunciar a su agresor.} (Porter plainte contre son agresseur.) \textbf{Transitif direct.} \\
proteger & protéger & \esp{Proteger a las víctimas.} (Protéger les victimes.) \\
prevenir & prévenir (avant l'acte) & \esp{Prevenir la violencia con educación.} (Prévenir la violence par l'éducation.) \\
sensibilizar & sensibiliser & \esp{Sensibilizar a la sociedad.} (Sensibiliser la société.) \\
luchar contra & lutter contre & \esp{Luchamos contra el machismo.} (Nous luttons contre le machisme.) \\
condenar & condamner (moral/juridique) & \esp{La sociedad condena la violencia.} (La société condamne la violence.) \\
apoyar a & soutenir, appuyer & \esp{Apoyar a las mujeres maltratadas.} (Soutenir les femmes battues.) \\
\hline
\end{tabularx}
\end{center}

\begin{aretenir}[Distinction \cle{prevenir} vs \cle{denunciar}]
\begin{itemize}
    \item \esp{Prevenir} = agir \textbf{avant} (éducation, campagnes). \esp{La prevención es clave.}
    \item \esp{Denunciar} = agir \textbf{après} le fait (acte judiciaire/police). \esp{Presentar una denuncia.}
    \item \esp{Proteger} = mesure immédiate (hébergement, ordonnance).
\end{itemize}
\end{aretenir}

\begin{exemplebox}[Engagement associatif]
\esp{Las ONG en Camerún trabajan para sensibilizar a los jóvenes en los colegios.} « Les ONG au Cameroun travaillent à sensibiliser les jeunes dans les lycées. » \\
\esp{Si eres testigo, denuncia: el silencio es cómplice.} « Si tu es témoin, dénonce : le silence est complice. » \\
\esp{El Estado debe proteger y apoyar a las víctimas, no solo condenar al agresor.} « L'État doit protéger et soutenir les victimes, pas seulement condamner l'agresseur. »
\end{exemplebox}

\courssub{Champ 5 : Exprimer son opinion et argumenter dans un débat (Opinar y argumentar)}

Le programme exige une \textbf{compétence orale} : participer à un débat guidé. Il faut structurer sa parole : \textbf{Opinion $\rightarrow$ Argument $\rightarrow$ Exemple $\rightarrow$ Conclusion}.

\begin{center}
\begin{tabularx}{\linewidth}{|X|X|X|}
\hline
\rowcolor{poppinkL}
\textbf{Fonction} & \textbf{Español} & \textbf{Français / Nuance} \\
\hline
\multicolumn{3}{|c|}{\textbf{Donner son avis}} \\
\hline
Opinion personnelle & \esp{En mi opinión, ...} / \esp{Desde mi punto de vista, ...} & « À mon avis / De mon point de vue » (formel). \\
Accord & \esp{Estoy de acuerdo (con...)} & « Je suis d'accord (avec...) ». \textbf{Pas *\esp{soy de acuerdo*}.} \\
Désaccord & \esp{No estoy de acuerdo / Discrepo.} & « Je ne suis pas d'accord / Je diverge. » \\
\hline
\multicolumn{3}{|c|}{\textbf{Structurer le discours}} \\
\hline
Premier argument & \esp{Por un lado, ...} & « D'un côté / Premièrement. » \\
Argument opposé & \esp{Por otro lado, ...} & « De l'autre côté / D'autre part. » \\
Contraste & \esp{Sin embargo, / No obstante, ...} & « Cependant / Néanmoins » (plus fort que \esp{pero}, suivi de virgule). \\
Ajout & \esp{Además, / Es más, ...} & « De plus / Qui plus est. » \\
\hline
\multicolumn{3}{|c|}{\textbf{Justifier (Cause / Conséquence)}} \\
\hline
Cause & \esp{Porque / Ya que / Puesto que ...} & « Parce que / Comme / Puisque » (\esp{ya que}, \esp{puesto que} souvent en tête de phrase). \\
Conséquence & \esp{Por eso / Por lo tanto / Por consiguiente ...} & « C'est pourquoi / Par conséquent / Donc. » \\
Renforcement & \esp{En efecto, ...} & « En effet / Effectivement » (confirme ce qui précède). \\
\hline
\end{tabularx}
\end{center}

\begin{methode}[Participer à un débat guidé (4 étapes)]
\begin{enumerate}
    \item \textbf{ÉCOUTER} activement : identifier la thèse de l'interlocuteur (\esp{¿Qué opina usted?}).
    \item \textbf{RÉAGIR} : valider ou contester poliment. \esp{Estoy de acuerdo contigo porque...} / \esp{No estoy de acuerdo, ya que...}
    \item \textbf{ARGUMENTER} : développer avec connecteurs. \esp{Por un lado... Por otro lado... Además... En efecto...}
    \item \textbf{CONCLURE} : synthétiser. \esp{En conclusión, creo que hay que...} / \esp{Para terminar, propongo...}
\end{enumerate}
\end{methode}

\begin{popfigure}
\begin{tikzpicture}[
    mindmap,
    grow cyclic,
    every node/.style=concept,
    concept color=popblue,
    level 1/.append style={level distance=4.5cm, sibling angle=90, concept color=popblue!70},
    level 2/.append style={level distance=3cm, sibling angle=45, concept color=poporange!70},
    level 3/.append style={level distance=2.5cm, sibling angle=30, concept color=popgreen!70},
    text=white,
    font=\small\bfseries
]
\node {El Debate}
    child [concept color=poporange] { node {Opinar}
        child { node {Estoy de acuerdo} }
        child { node {No estoy de acuerdo} }
        child { node {En mi opinión} }
        child { node {Desde mi punto de vista} }
    }
    child [concept color=popgreen] { node {Argumentar}
        child { node {Porque / Ya que} }
        child { node {Por un lado} }
        child { node {Por otro lado} }
        child { node {Además / Es más} }
        child { node {Sin embargo} }
    }
    child [concept color=poppurple] { node {Reaccionar}
        child { node {Denunciar} }
        child { node {Apoyar a la víctima} }
        child { node {Sensibilizar} }
        child { node {Proponer soluciones} }
    }
    child [concept color=poppink] { node {Concluir}
        child { node {En conclusión} }
        child { node {Por lo tanto} }
        child { node {Para terminar} }
    };
\end{tikzpicture}
\legende{Carte mentale : Structure lexicale d'un débat oral sur la violence de genre.}
\end{popfigure}

\begin{savaistu}[25 Noviembre : Día Internacional]
Le \textbf{25 novembre} est la \esp{Día Internacional de la Eliminación de la Violencia contra la Mujer}. Instituée par l'ONU en 1999, elle commémore l'assassinat des sœurs Mirabal (République dominicaine, 1960). En Espagne et dans tout le monde hispanophone, c'est une journée de \textbf{mobilisation massive} : manifestations (\esp{manifestaciones}), minutes de silence (\esp{minutos de silencio}), bâtiments illuminés en violet (\esp{iluminados de morado}), lectures de manifestes. Au Cameroun, les associations féministes et le MINPROFF organisent des marches à Yaoundé et Douala ce jour-là. Le violet (\esp{morado}) est la couleur symbolique du féminisme hispanique.
\end{savaistu}

\courssub{Texte intégrateur : « La charla en el comedor »}

Le texte suivant met en scène une famille camerounaise (père, mère, fils aîné, fille) discutant d'un fait divers local. Il réinvestit \textbf{tous les champs lexicaux} vus ci-dessus.

\begin{textebox}[Dialogue familial : débat sur un fait de violence de genre]
-- \textbf{Padre:} \esp{¿Han visto las noticias? En el barrio de Mvog-Ada, un hombre ha sido detenido por malos tratos a su esposa. La víctima está en el hospital.} \\
-- \textbf{Madre:} \esp{¡Dios mío! Otra vez la violencia de género. \textbf{Por un lado}, me duele el alma; \textbf{por otro lado}, me da rabia la impunidad.} \\
-- \textbf{Hijo (Kevin, 17 ans):} \esp{\textbf{En mi opinión}, el problema es cultural. \textbf{Porque} muchos hombres creen que \textbf{la mujer} es su propiedad. \textbf{Además}, no denuncian por miedo.} \\
-- \textbf{Hija (Grace, 15 ans):} \esp{\textbf{Estoy de acuerdo con Kevin}. \textbf{Puesto que} la educación \textbf{en el hogar} falla. \textbf{Sin embargo}, \textbf{los padres} debemos \textbf{educar} en el respeto. \textbf{Ya que} si no \textbf{prevenimos}, luego hay que \textbf{denunciar}.} \\
-- \textbf{Padre:} \esp{\textbf{En efecto}. \textbf{Los parientes} y los vecinos a veces callan. \textbf{Hay que} \textbf{sensibilizar} a la comunidad. \textbf{Por lo tanto}, si vemos \textbf{acoso} o \textbf{agresión}, debemos \textbf{apoyar a la víctima} y \textbf{denunciar al maltratador}.} \\
-- \textbf{Madre:} \esp{Exacto. \textbf{La convivencia} sana se basa en \textbf{respetar} y \textbf{confiar}. \textbf{No en} \textbf{pelearse} ni \textbf{controlar}. \textbf{Por eso}, esta noche, \textbf{Kevin, Grace}, hablemos de \textbf{la crianza} que queréis para vuestro futuro.}
\end{textebox}

\begin{center}
\textbf{Glossaire des mots difficiles (contexte) :}
\end{center}
\begin{itemize}
    \item \esp{Mvog-Ada} : Quartier populaire de Yaoundé (contexte camerounais réaliste).
    \item \esp{Detenido} : Interpellé, arrêté (participe de \esp{detener}).
    \item \esp{Impunidad} : Impunité (fait de ne pas être puni).
    \item \esp{Propiedad} : Propriété, bien matériel.
    \item \esp{Callan} : Se taisent (verbe \esp{callar}, indicatif présent 3e pers. pl.).
    \item \esp{Controlar} : Contrôler (au sens psychologique : surveillance, isolement).
    \item \esp{Hablemos} : Parlons (impératif \emph{nosotros} / subjonctif présent 1re pers. pl.).
    \item \esp{Queréis} : Vous voulez (indicatif présent \emph{vosotros}, usage péninsulaire standard au programme).
\end{itemize}

\begin{exoresolu}[Traduction thème/version : Lexique famille \& violence]
\textbf{Énoncé :} Traduire en espagnol les phrases suivantes en mobilisant le vocabulaire des 5 champs.
\begin{enumerate}
    \item Mes parents (père et mère) s'entendent bien avec mes grands-parents.
    \item La violence de genre est un fléau ; il faut sensibiliser les jeunes pour la prévenir.
    \item L'agresseur a été condamné car la victime a osé porter plainte.
    \item À mon avis, l'éducation à la maison est la clé ; cependant, l'État doit protéger les victimes.
    \item D'un côté, les voisins se taisent ; de l'autre, les associations luttent contre le machisme.
\end{enumerate}

\tcblower
\textbf{Solution détaillée :}
\begin{enumerate}
    \item \esp{Mis \textbf{padres} se \textbf{llevan bien} con mis \textbf{abuelos}.} (\cle{padres} = parents ; \cle{llevarse bien con} + pronominal réciproque).
    \item \esp{La \textbf{violencia de género} es un \textbf{flagelo} (ou \textbf{lacra}, \textbf{plaga}) ; \textbf{hay que sensibilizar} a los jóvenes para \textbf{prevenirla}.} (\cle{hay que} + infinitif = il faut ; \cle{prevenirla} = la prévenir, pronom enclitique obligatoire).
    \item \esp{El \textbf{maltratador} (ou \textbf{agresor}) fue \textbf{condenado} porque la \textbf{víctima} se \textbf{atrevió a denunciar}.} (\cle{fue condenado} = passif passé simple \emph{ser} + participe ; \cle{atreverse a} = oser, pronominal).
    \item \esp{\textbf{En mi opinión}, la \textbf{educación en el hogar} es la \textbf{clave} ; \textbf{sin embargo}, el Estado debe \textbf{proteger a las víctimas}.} (\cle{sin embargo} = cependant, suivi de virgule ; \cle{debe} = doit, indicatif obligation).
    \item \esp{\textbf{Por un lado}, los vecinos \textbf{callan} (ou \textbf{guardan silencio}) ; \textbf{por otro lado}, las asociaciones \textbf{luchan contra el machismo}.} (\cle{por un lado / por otro lado} = structure fixe ; \cle{luchar contra} + nom).
\end{enumerate}
\end{exoresolu}

\courssub{Exercices d'appropriation}

\begin{exercice}[Complétion lexicale et choix grammatical]
Complétez les phrases avec le mot juste du champ lexical (forme correcte).
\begin{enumerate}
    \item Mi tío y mi tía son mis \_\_\_\_\_\_\_\_\_\_\_\_. (parientes / padres)
    \item \_\_\_\_\_\_\_\_\_\_\_\_ la ley, el maltratador tiene prohibido acercarse a la víctima. (Segun / Según)
    \item No \_\_\_\_\_\_\_\_\_\_\_\_ con tu hermana; \_\_\_\_\_\_\_\_\_\_\_\_ pronto. (pelearse / reconciliarse - impératif \emph{vosotros})
    \item \_\_\_\_\_\_\_\_\_\_\_\_ mi punto de vista, \_\_\_\_\_\_\_\_\_\_\_\_ es la mejor arma contra la violencia. (En / Desde ; la educación / la denuncia)
    \item Los vecinos \_\_\_\_\_\_\_\_\_\_\_\_ a la policía cuando oyeron la \_\_\_\_\_\_\_\_\_\_\_\_. (denunciaron / callaron ; agresión / pelea)
\end{enumerate}

\end{exercice}

\begin{exercice}[Production écrite guidée : « Carta al director »]
\textbf{Consigne :} Vous êtes élève en Terminale A au Lycée de Nkolbisson (Yaoundé). Écrivez une lettre formelle (120-150 mots) au directeur du journal \emph{Mutations} pour réagir à un article sur une agression conjugale.
\textbf{Contraintes :}
\begin{itemize}
    \item Utilisez au moins \textbf{10 mots} du lexique des 5 champs (surlignez-les).
    \item Structurez : Salutation $\rightarrow$ Opinion (\esp{En mi opinión}) $\rightarrow$ 2 arguments (\esp{Porque}, \esp{Además}) $\rightarrow$ Contre-argument/Contraste (\esp{Sin embargo}) $\rightarrow$ Proposition (\esp{Propongo que...} + subjonctif) $\rightarrow$ Formule de politesse.
    \item Soignez \textbf{ser/estar} (révision Section 4 à venir) : \esp{La violencia \textbf{es} un problema estructural} vs \esp{La víctima \textbf{está} traumatizada}.
\end{itemize}
\end{exercice}

\coursec{Méthode du commentaire et commentaire modèle}

Dans la section précédente, nous avons constitué le lexique de la famille et de la violence de genre : \esp{la pareja, los malos tratos, denunciar, la igualdad, el maltratador}... Ces mots ne sont pas seulement utiles pour parler : ils vont maintenant nous servir à \cle{analyser un texte}. Au baccalauréat, l'épreuve de commentaire de texte exige de passer du vocabulaire à l'argumentation. C'est l'objet de cette section : apprendre à construire un commentaire rigoureux, puis lire un commentaire modèle rédigé sur notre texte support.

\courssub{Le texte support : un dialogue familial}

Relisons d'abord le texte sur lequel portera tout notre travail de commentaire. Il s'agit d'un dialogue familial original, situé à Yaoundé, où trois générations abordent le problème des violences faites aux femmes.

\begin{textebox}[Un diálogo en familia — « Hablar es el primer paso »]
Es domingo por la tarde en Yaoundé. Después de la comida, la familia Etoundi se reúne en el salón. La hija mayor, Carine, que estudia en el liceo, propone un tema serio.

— Abuelo, mamá, quiero hablaros de algo importante. En el barrio, una vecina sufre malos tratos de su marido y nadie dice nada.

— Hija —responde el abuelo—, antes estas cosas no se hablaban. El hombre mandaba en casa y punto.

— Pero abuelo, ¡es terrible que todavía haya mujeres que callen por miedo! Me alegro de que cada vez más mujeres denuncien a sus maltratadores.

— Tu hermana tiene razón —interviene la madre—. Cuando yo era joven, no existían asociaciones de ayuda. Hoy hay centros que escuchan a las víctimas.

— ¿De verdad crees que las cosas han cambiado tanto? —pregunta el abuelo.

— Sí, y me alegro de que los jóvenes como tú, Carine, habléis de este problema sin vergüenza. Antes era un tabú.

— ¿Y qué podemos hacer nosotros? —insiste Carine.

— Podemos empezar por escuchar a nuestra vecina y acompañarla a denunciar. El silencio es el mejor aliado del maltratador.

El abuelo sonríe. Comprende que los tiempos han cambiado y que el diálogo en familia es el primer paso para construir una sociedad más justa, donde hombres y mujeres vivan con la misma dignidad.
\end{textebox}

\textbf{Glossaire :} \esp{los malos tratos} = les mauvais traitements ; \esp{mandar} = commander, décider ; \esp{y punto} = et un point c'est tout ; \esp{callar} = se taire ; \esp{denunciar} = dénoncer, porter plainte ; \esp{el maltratador} = le maltraitant, l'auteur des violences ; \esp{la vergüenza} = la honte ; \esp{un tabú} = un tabou ; \esp{acompañarla} = l'accompagner ; \esp{el aliado} = l'allié ; \esp{la dignidad} = la dignité.

\textbf{Questions de compréhension :}
\begin{enumerate}
\item ¿Qué problema plantea Carine a su familia ?
\item ¿Cómo reacciona el abuelo al principio del diálogo ?
\item ¿Qué cambios sociales menciona la madre ?
\item ¿Qué solución concreta propone la familia al final ?
\end{enumerate}

\begin{corrige}[Questions de compréhension]
\begin{enumerate}
\item \esp{Carine plantea el problema de una vecina que sufre malos tratos de su marido y del silencio de todo el barrio.} « Carine pose le problème d'une voisine qui subit des mauvais traitements de la part de son mari, et du silence de tout le quartier. »
\item \esp{Al principio, el abuelo defiende la tradición: dice que antes estas cosas no se hablaban y que el hombre mandaba en casa.} « Au début, le grand-père défend la tradition : il dit qu'avant on ne parlait pas de ces choses et que l'homme décidait à la maison. »
\item \esp{La madre menciona la existencia de asociaciones de ayuda y de centros que escuchan a las víctimas, que no existían cuando ella era joven.} « La mère mentionne l'existence d'associations d'aide et de centres qui écoutent les victimes, ce qui n'existait pas quand elle était jeune. »
\item \esp{La familia propone escuchar a la vecina y acompañarla a denunciar a su maltratador.} « La famille propose d'écouter la voisine et de l'accompagner pour porter plainte contre son maltraitant. »
\end{enumerate}
\end{corrige}

\courssub{Rappel : la structure du commentaire de texte au baccalauréat}

\begin{definition}[Le commentaire de texte]
Le \cle{commentaire de texte} est un exercice d'analyse : il ne s'agit pas de raconter le texte, mais d'expliquer \textbf{ce qu'il dit} (le thème), \textbf{comment il le dit} (les procédés linguistiques) et \textbf{pourquoi c'est important} (la problématique et l'ouverture). En espagnol : \esp{el comentario de texto}.
\end{definition}

Le commentaire se compose de trois parties obligatoires, comme en français :

\begin{center}
\begin{tabularx}{\linewidth}{|X|X|X|}
\hline
\rowcolor{popblueL} Partie & Contenu & Durée indicative \\
\hline
\esp{Introducción} & Présenter le texte : type, source, thème, personnages ; annoncer la problématique & 4 à 6 lignes \\
\hline
\esp{Análisis (desarrollo)} & Analyser en 2 ou 3 paragraphes : thème, personnages, procédés linguistiques, avec de courtes citations & 15 à 20 lignes \\
\hline
\esp{Conclusión} & Bilan de l'analyse + ouverture vers un problème plus large & 4 à 6 lignes \\
\hline
\end{tabularx}
\end{center}

\begin{popfigure}
\begin{tikzpicture}[
  node distance=0.55cm and 1.2cm,
  caja/.style={draw=popblue, fill=popblueL, rounded corners, align=center, text width=4.6cm, inner sep=6pt, font=\small},
  detalle/.style={draw=popmuted, fill=popcream, rounded corners, align=left, text width=6.4cm, inner sep=6pt, font=\small},
  flecha/.style={-{Stealth[length=3mm]}, thick, popdark}
]
\node[caja, align=center] (intro){\textbf{1. Introducción}\\ Tipo de texto, tema,\\ personajes, problemática};
\node[caja, below=of intro, align=center] (anal){\textbf{2. Análisis}\\ Tema y problemática\\ Personajes y puntos de vista\\ Procedimientos lingüísticos};
\node[caja, below=of anal, align=center] (concl){\textbf{3. Conclusión}\\ Balance + apertura};
\node[detalle, right=of intro] (d1) {«Este texto es un diálogo familiar... Trata de... Los personajes son...»};
\node[detalle, right=of anal, align=center] (d2){Ideas + citas cortas :\\ «es terrible que... callen»,\\ «me alegro de que... denuncien»};
\node[detalle, right=of concl] (d3) {«En conclusión... Este texto nos invita a reflexionar sobre...»};
\draw[flecha] (intro) -- (anal);
\draw[flecha] (anal) -- (concl);
\draw[flecha, popmuted, dashed] (intro.east) -- (d1.west);
\draw[flecha, popmuted, dashed] (anal.east) -- (d2.west);
\draw[flecha, popmuted, dashed] (concl.east) -- (d3.west);
\end{tikzpicture}
\legende{Organigramme : la structure du commentaire de texte et le contenu de chaque partie.}
\end{popfigure}

\courssub{Étape 1 : dégager le thème et la problématique}

Avant d'écrire, il faut répondre à deux questions : \textbf{de quoi parle le texte ?} (le thème) et \textbf{quelle question pose-t-il ?} (la problématique).

Ici, le thème est double : le \cle{dialogue familial} comme espace de parole, et les \cle{violences faites aux femmes} comme problème de société. La problématique peut se formuler ainsi : \esp{¿Cómo puede el diálogo en familia contribuir a la lucha contra la violencia de género?} (« Comment le dialogue en famille peut-il contribuer à la lutte contre la violence de genre ? »).

\begin{exemplebox}[Formules utiles pour annoncer le thème et la problématique]
\begin{itemize}
\item \esp{El texto plantea el problema de la violencia de género.} « Le texte pose le problème de la violence de genre. »
\item \esp{El tema central es el diálogo familiar como espacio de palabra.} « Le thème central est le dialogue familial comme espace de parole. »
\item \esp{El autor quiere mostrar que el silencio es el mejor aliado del maltratador.} « L'auteur veut montrer que le silence est le meilleur allié du maltraitant. »
\item \esp{La cuestión que se plantea es la siguiente: ¿han cambiado realmente las mentalidades?} « La question qui se pose est la suivante : les mentalités ont-elles vraiment changé ? »
\end{itemize}
\end{exemplebox}

\courssub{Étape 2 : analyser les personnages et leurs points de vue}

Un dialogue se commente d'abord par ses \cle{voix}. Chaque personnage représente un point de vue, souvent lié à sa génération :

\begin{center}
\begin{tabularx}{\linewidth}{|X|X|X|}
\hline
\rowcolor{popgreenL} Personaje & Punto de vista & Cita que lo demuestra \\
\hline
El abuelo & La tradition : avant, on se taisait ; puis il évolue & «antes estas cosas no se hablaban» \\
\hline
La madre & Le témoin du changement : les associations n'existaient pas & «hoy hay centros que escuchan a las víctimas» \\
\hline
Carine & La jeunesse engagée : il faut parler et dénoncer & «me alegro de que cada vez más mujeres denuncien» \\
\hline
\end{tabularx}
\end{center}

L'analyse doit dégager l'\textbf{évolution des mentalités} : le grand-père passe de la défense du silence (\esp{el hombre mandaba en casa y punto}) à l'adhésion au dialogue (\esp{el abuelo sonríe}). C'est le mouvement même du texte : la parole familiale transforme les regards.

\begin{exemplebox}[Formules pour analyser un personnage]
\begin{itemize}
\item \esp{El abuelo muestra la evolución de las mentalidades.} « Le grand-père montre l'évolution des mentalités. »
\item \esp{Carine representa a una juventud comprometida que rompe el tabú.} « Carine représente une jeunesse engagée qui brise le tabou. »
\item \esp{La madre hace de puente entre las dos generaciones.} « La mère fait le pont entre les deux générations. »
\end{itemize}
\end{exemplebox}

\courssub{Étape 3 : relever les procédés linguistiques}

Le correcteur attend que vous montriez \textbf{comment} la langue construit le sens. Trois procédés dominent dans ce texte :

\begin{enumerate}
\item \textbf{Les marqueurs d'opinion} : \esp{es terrible que..., me alegro de que..., ¿de verdad crees que...?} — ils expriment le jugement et le sentiment.
\item \textbf{Le subjonctif} après les expressions de sentiment et de jugement : \esp{me alegro de que las mujeres denuncien}, \esp{es terrible que haya mujeres que callen}.
\item \textbf{Les questions de relance} : \esp{¿De verdad crees que las cosas han cambiado tanto?}, \esp{¿Y qué podemos hacer nosotros?} — elles font avancer le dialogue et montrent l'écoute active.
\end{enumerate}

\begin{propriete}[Subjonctif après les expressions de sentiment et de jugement]
Après les verbes et expressions de \textbf{sentiment} (\esp{alegrarse de que, sentir que, temer que}) et de \textbf{jugement} (\esp{es terrible que, es importante que, es una lástima que}), la subordonnée introduite par \esp{que} se met au \cle{subjonctif} :
\begin{itemize}
\item \esp{Me alegro de que las mujeres denuncien.} « Je me réjouis que les femmes portent plainte. » (et non \esp{denuncian})
\item \esp{Es terrible que todavía haya mujeres que callen.} « Il est terrible qu'il y ait encore des femmes qui se taisent. »
\item \esp{Es importante que la familia escuche a las víctimas.} « Il est important que la famille écoute les victimes. »
\end{itemize}
En français, on hésite entre indicatif et subjonctif ; en espagnol, le subjonctif est \textbf{obligatoire} après ces tournures. Exception à connaître : après les verbes d'opinion à la forme \textbf{affirmative} (\esp{creo que, pienso que}), on utilise l'indicatif (\esp{creo que las cosas han cambiado}) ; c'est seulement à la forme \textbf{négative} ou interrogative dubitative que le subjonctif devient possible (\esp{no creo que las cosas hayan cambiado}).
\end{propriete}

\begin{attention}[Ne pas résumer : analyser et citer]
L'erreur la plus fréquente au baccalauréat est de \textbf{raconter le texte} (« d'abord Carine parle, puis le grand-père répond, puis la mère... »). Un commentaire n'est pas un résumé. Chaque idée doit être :
\begin{enumerate}
\item \textbf{énoncée} : « le texte dénonce le silence qui protège le maltraitant » ;
\item \textbf{appuyée par une courte citation} : \esp{el silencio es el mejor aliado del maltratador} ;
\item \textbf{expliquée} : la métaphore de l'« allié » montre que ne rien dire, c'est aider le coupable.
\end{enumerate}
Autre piège : ne recopiez jamais de longues phrases du texte ; une citation efficace compte de trois à huit mots.
\end{attention}

\begin{methode}[Rédiger l'introduction du commentaire]
\begin{enumerate}
\item \textbf{Présenter le type de texte et sa source} : \esp{Este texto es un diálogo familiar extraído de una lección sobre la vida familiar.}
\item \textbf{Annoncer le thème} : \esp{Trata de las violencias sufridas por las mujeres y del papel de la familia en esta lucha.}
\item \textbf{Présenter les personnages} : \esp{Los personajes son el abuelo, la madre y Carine, una joven estudiante.}
\item \textbf{Formuler la problématique} : \esp{La cuestión central es saber cómo el diálogo en familia puede romper el silencio que protege al maltratador.}
\item \textbf{Annoncer le plan} : \esp{Analizaremos primero el tema, después los puntos de vista de los personajes y, por último, los procedimientos lingüísticos.}
\end{enumerate}
\end{methode}

\courssub{Commentaire modèle rédigé}

Voici un commentaire complet, rédigé en espagnol, avec un plan apparent (introduction, deux axes d'analyse, conclusion). Lisez-le en repérant les formules étudiées plus haut.

\begin{textebox}[Comentario de texto — modèle rédigé]
\textbf{Introducción.} Este texto es un diálogo familiar situado en Yaoundé. Trata de un problema grave de nuestra sociedad: los malos tratos sufridos por las mujeres. Los personajes son el abuelo, la madre y Carine, una joven estudiante de liceo. La cuestión central es la siguiente: ¿cómo puede el diálogo en familia contribuir a la lucha contra la violencia de género? Analizaremos primero los puntos de vista de los personajes y después los procedimientos lingüísticos del diálogo.

\textbf{I. Tres generaciones, tres miradas.} El texto pone en escena la evolución de las mentalidades. El abuelo representa la tradición: «antes estas cosas no se hablaban», afirma al principio. Sin embargo, al final del diálogo, «el abuelo sonríe»: comprende que los tiempos han cambiado. La madre hace de puente entre las generaciones y recuerda los progresos sociales: «hoy hay centros que escuchan a las víctimas». Por último, Carine representa a una juventud comprometida que rompe el tabú y propone una acción concreta: acompañar a la vecina a denunciar.

\textbf{II. La lengua al servicio del debate.} El diálogo utiliza numerosos marcadores de opinión y de sentimiento seguidos de subjuntivo: «es terrible que todavía haya mujeres que callen», «me alegro de que cada vez más mujeres denuncien». Estas construcciones expresan el compromiso emocional de los personajes. Además, las preguntas que relanzan el intercambio —«¿de verdad crees que las cosas han cambiado tanto?»— muestran una escucha activa que hace avanzar la reflexión. Finalmente, la frase «el silencio es el mejor aliado del maltratador» resume la tesis del texto mediante una metáfora poderosa.

\textbf{Conclusión.} En conclusión, este diálogo muestra que hablar en familia es el primer paso para combatir la violencia de género. Cada generación aporta su mirada y, juntas, construyen una respuesta. Este texto nos invita a reflexionar sobre un problema más amplio: la educación de los niños y las niñas en la igualdad, condición indispensable para construir una sociedad más justa.
\end{textebox}

\textbf{Traduction des formules clés du modèle :} \esp{pone en escena} = met en scène ; \esp{hace de puente} = fait le pont ; \esp{compromiso emocional} = engagement émotionnel ; \esp{las preguntas que relanzan el intercambio} = les questions de relance ; \esp{hace avanzar la reflexión} = fait avancer la réflexion ; \esp{una metáfora poderosa} = une métaphore puissante ; \esp{condición indispensable} = condition indispensable.

\begin{attention}[Faux amis à éviter dans le commentaire]
\begin{itemize}
\item \esp{actualmente} = « actuellement », mais \esp{actual} = « actuel » ; « un événement actuel » se dit \esp{un acontecimiento actual}, jamais \esp{actualmente}.
\item \esp{asistir a} = « assister à » (être présent) ; « aider » se dit \esp{ayudar}. On dira donc : \esp{la familia ayuda a la vecina}.
\item \esp{comprometido} = « engagé » (politiquement, socialement) : \esp{una juventud comprometida} = « une jeunesse engagée » — c'est un bon usage, à ne pas confondre avec le français « compromis ».
\item « réaliser » au sens de « se rendre compte » = \esp{darse cuenta de}, et non \esp{realizar} (qui signifie « réaliser, accomplir » un projet).
\end{itemize}
\end{attention}

\courssub{Production guidée : rédiger une conclusion}

\begin{exoresolu}[Rédiger le paragraphe de conclusion d'un commentaire]
Énoncé. À partir du plan donné, rédigez en espagnol la conclusion (4 à 6 lignes) d'un commentaire sur le dialogue des Etoundi. Plan : (a) bilan — le dialogue familial brise le silence ; (b) bilan — les mentalités évoluent grâce aux jeunes ; (c) ouverture — le rôle de l'école dans l'éducation à l'égalité.
\tcblower
Solution détaillée. Une conclusion = bilan + ouverture, sans idée nouvelle d'analyse.

\esp{En conclusión, este diálogo demuestra que la familia es el primer lugar donde se puede romper el silencio que protege al maltratador. Gracias a jóvenes como Carine, las mentalidades evolucionan: el abuelo, que defendía la tradición, acepta al final que denunciar es un deber. Sin embargo, la lucha no termina en casa: este texto nos invita a reflexionar sobre el papel de la escuela, que debe educar a los niños y a las niñas en la igualdad y el respeto desde la más temprana edad.}

« En conclusion, ce dialogue démontre que la famille est le premier lieu où l'on peut briser le silence qui protège le maltraitant. Grâce à des jeunes comme Carine, les mentalités évoluent : le grand-père, qui défendait la tradition, accepte à la fin que dénoncer est un devoir. Cependant, la lutte ne s'arrête pas à la maison : ce texte nous invite à réfléchir au rôle de l'école, qui doit éduquer les garçons et les filles à l'égalité et au respect dès le plus jeune âge. »

Points de méthode : la phrase d'ouverture commence par un connecteur de transition (\esp{sin embargo}, « cependant ») et la formule \esp{nos invita a reflexionar sobre...} est la clausule type du commentaire au baccalauréat.
\end{exoresolu}

\begin{exercice}[À vous de commenter]
Rédigez en espagnol l'introduction complète (5 lignes) d'un commentaire sur le dialogue des Etoundi, en suivant les cinq étapes de la boîte méthode. Utilisez au moins deux des formules suivantes : \esp{este texto es un diálogo familiar...}, \esp{trata de...}, \esp{la cuestión central es...}, \esp{analizaremos primero...}.
\end{exercice}

\begin{corrige}[Exercice — introduction]
Éléments de réponse attendus (modèle) : \esp{Este texto es un diálogo familiar que tiene lugar un domingo por la tarde en Yaoundé. Trata de los malos tratos sufridos por una vecina y de la reacción de la familia Etoundi. Los personajes son el abuelo, la madre y Carine, una joven estudiante comprometida. La cuestión central es saber si el diálogo en familia puede ayudar a combatir la violencia de género. Analizaremos primero el tema, después los personajes y finalmente los procedimientos lingüísticos.}
Barème indicatif : type de texte (1 pt), thème (1 pt), personnages (1 pt), problématique (1 pt), annonce du plan et correction linguistique (1 pt).
\end{corrige}

\begin{aretenir}[L'essentiel de la méthode]
\begin{itemize}
\item Un commentaire = \textbf{introduction} (type, thème, personnages, problématique) + \textbf{analyse} (2 ou 3 axes, idées + courtes citations) + \textbf{conclusion} (bilan + ouverture).
\item On \cle{analyse}, on ne résume jamais : chaque idée est prouvée par une citation de trois à huit mots.
\item Les procédés à repérer dans un dialogue d'opinion : marqueurs de sentiment suivis du \textbf{subjonctif} (\esp{me alegro de que denuncien}) et questions de relance.
\item Formules à mémoriser : \esp{el texto plantea el problema de...}, \esp{el abuelo muestra la evolución de las mentalidades}, \esp{este texto nos invita a reflexionar sobre...}.
\end{itemize}
\end{aretenir}

\coursec{Grammaire : SER et ESTAR — emplois fondamentaux}

Après avoir commenté le dialogue familial et débattu de la violence de genre, nous avons rencontré à plusieurs reprises, dans le texte support, deux verbes qui traduisent tous les deux le français « être » : \esp{ser} et \esp{estar}. C'est l'une des difficultés majeures pour un francophone : là où le français n'a qu'un seul verbe, l'espagnol en impose deux, et le choix n'est jamais arbitraire. Cette section de grammaire va nous donner les clés pour choisir correctement, compétence indispensable pour la traduction (thème et version) comme pour l'expression orale.

\courssub{Observation : deux phrases, deux verbes}

Commençons par observer deux phrases simples :

\begin{exemplebox}[Observation guidée]
\esp{María es camerunesa.} « María est camerounaise. »

\esp{María está en Douala.} « María est à Douala. »
\end{exemplebox}

Dans les deux cas, le français utilise « est ». Pourtant l'espagnol choisit \esp{es} (verbe \esp{ser}) dans la première phrase et \esp{está} (verbe \esp{estar}) dans la seconde. Pourquoi ?

\begin{itemize}
\item \esp{María es camerunesa} : la nationalité définit qui est María ; c'est une caractéristique de son \cle{identité}, permanente, qui ne change pas d'un jour à l'autre.
\item \esp{María está en Douala} : sa présence à Douala est une \cle{circonstance} ; demain, elle peut être à Yaoundé ou à Bafoussam.
\end{itemize}

L'observation nous conduit donc à l'intuition fondamentale : \esp{ser} dit ce qu'on \textbf{est} (essence), \esp{estar} dit \textbf{comment} et \textbf{où} on est (état).

\begin{propriete}[Règle fondamentale]
\cle{SER} exprime l'\textbf{identité} et les \textbf{qualités permanentes} : \esp{Soy estudiante.} « Je suis élève. »

\cle{ESTAR} exprime la \textbf{localisation} et les \textbf{états temporaires} : \esp{Estoy cansado.} « Je suis fatigué. »
\end{propriete}

\courssub{Conjugaison au présent de l'indicatif}

Avant d'étudier les emplois, il faut connaître parfaitement les deux conjugaisons, car ce sont deux verbes irréguliers parmi les plus fréquents de la langue.

\begin{center}
\begin{tabular}{|l|l|l|l|}
\hline
\rowcolor{popblueL} \textbf{Personne} & \textbf{SER} & \textbf{ESTAR} & \textbf{Français} \\
\hline
yo & soy & estoy & je suis \\
\hline
tú & eres & estás & tu es \\
\hline
él / ella / usted & es & está & il / elle est \\
\hline
nosotros / nosotras & somos & estamos & nous sommes \\
\hline
vosotros / vosotras & sois & estáis & vous êtes \\
\hline
ellos / ellas / ustedes & son & están & ils / elles sont \\
\hline
\end{tabular}
\end{center}

\begin{attention}[Irrégularités à mémoriser]
\esp{ser} est fortement irrégulier (\esp{soy, eres, es, somos, sois, son}) : aucune forme ne ressemble à l'infinitif. \esp{estar} est régulier sauf à la première personne (\esp{estoy}) et porte un accent sur toutes les personnes sauf \esp{estoy} et \esp{estamos} : \esp{estás, está, estáis, están}. N'écrivez jamais \esp{*esta} sans accent pour « il est » : \esp{esta} sans accent signifie « cette » (adjectif démonstratif féminin).
\end{attention}

\courssub{Les emplois de SER}

Le verbe \esp{ser} s'emploie dans les cas suivants :

\begin{enumerate}
\item \textbf{L'identité} (qui on est) : \esp{Yo soy Aminatou.} « Je suis Aminatou. » \esp{Este es mi hermano Carlos.} « Voici mon frère Carlos. »
\item \textbf{La nationalité} : \esp{Somos cameruneses.} « Nous sommes camerounais. » \esp{Ellos son españoles.} « Ils sont espagnols. »
\item \textbf{La profession} : \esp{Mi padre es médico.} « Mon père est médecin. » Remarquez l'absence d'article en espagnol : on dit \esp{es médico} et non \esp{*es un médico} (sauf si l'on ajoute une qualification : \esp{es un médico excelente}).
\item \textbf{L'origine} (\esp{ser de}) : \esp{Soy de Yaoundé.} « Je suis de Yaoundé. » \esp{¿De dónde eres?} « D'où es-tu ? »
\item \textbf{Les caractéristiques permanentes} (physiques ou morales) : \esp{La casa es grande.} « La maison est grande. » \esp{Mi abuela es muy generosa.} « Ma grand-mère est très généreuse. »
\item \textbf{La possession} (\esp{ser de}) : \esp{El libro es de mi profesora.} « Le livre est à ma professeure. » \esp{¿De quién es este bolígrafo?} « À qui est ce stylo ? »
\item \textbf{La matière} : \esp{La mesa es de madera.} « La table est en bois. » \esp{El anillo es de oro.} « La bague est en or. »
\item \textbf{L'heure et la date} : \esp{Son las cinco.} « Il est cinq heures. » \esp{Es la una.} « Il est une heure. » \esp{Hoy es lunes.} « Aujourd'hui, c'est lundi. » \esp{Es el tres de mayo.} « C'est le trois mai. »
\item \textbf{Les événements} (où et quand ils ont lieu) : \esp{¿Dónde es la fiesta?} « Où est la fête ? » \esp{El partido es el sábado.} « Le match est samedi. »
\end{enumerate}

\courssub{Les emplois de ESTAR}

Le verbe \esp{estar} s'emploie dans les cas suivants :

\begin{enumerate}
\item \textbf{La localisation des personnes et des choses} : \esp{Mi padre está en el hospital.} « Mon père est à l'hôpital. » \esp{Douala está en el litoral.} « Douala est sur le littoral. » \esp{¿Dónde están mis llaves?} « Où sont mes clés ? »
\item \textbf{Les états temporaires physiques} : \esp{Estoy cansado.} « Je suis fatigué. » \esp{El niño está enfermo.} « L'enfant est malade. »
\item \textbf{Les états émotionnels} : \esp{Estamos contentos.} « Nous sommes contents. » \esp{Ella está triste hoy.} « Elle est triste aujourd'hui. » \esp{¿Estás nervioso por el examen?} « Es-tu nerveux à cause de l'examen ? »
\item \textbf{Le résultat d'un changement} : \esp{La sopa está fría.} « La soupe est froide » (elle a refroidi). \esp{La puerta está abierta.} « La porte est ouverte » (quelqu'un l'a ouverte).
\item \textbf{L'action en cours} (\esp{estar} + gérondif) : \esp{Estamos estudiando.} « Nous sommes en train d'étudier. » \esp{¿Qué estás haciendo?} « Qu'es-tu en train de faire ? »
\end{enumerate}

\begin{exemplebox}[Exemples traduits]
\esp{Mi padre es médico.} « Mon père est médecin. » (profession : \esp{ser})

\esp{Mi padre está en el hospital.} « Mon père est à l'hôpital. » (localisation : \esp{estar})

\esp{Son las cinco.} « Il est cinq heures. » (heure : \esp{ser})

\esp{Estamos estudiando.} « Nous sommes en train d'étudier. » (action en cours : \esp{estar} + gérondif)

\esp{La reunión es en el salón.} « La réunion est dans la salle. » (événement : \esp{ser})

\esp{El salón está al fondo.} « La salle est au fond. » (chose : \esp{estar})
\end{exemplebox}

\courssub{Texte de synthèse : une famille camerounaise}

Lisons maintenant un court texte où les deux verbes apparaissent en contexte. Soulignez mentalement chaque occurrence de \esp{ser} ou \esp{estar} et justifiez le choix.

\begin{textebox}[La familia Ngo Bell]
La familia Ngo Bell es de Douala, pero ahora está en Yaoundé porque el padre, que es ingeniero, está trabajando en un proyecto de construcción de carreteras. La madre es profesora de español en un liceo; es una mujer muy paciente y generosa, pero hoy está cansada porque son las diez de la noche y todavía está corrigiendo exámenes.

Los hijos son tres: Samuel, que es el mayor, es estudiante de medicina y está muy ocupado porque los exámenes son la próxima semana. Carmen, la mediana, es muy alegre por naturaleza, pero esta tarde está triste: su equipo de fútbol favorito está perdiendo el partido. El pequeño, Andrés, es travieso; ahora mismo está en el jardín y está jugando con el perro.

La casa es de ladrillo y está cerca del mercado. La abuela, que está viuda, está en el salón: está viendo una telenovela. La cena está lista y toda la familia está reunida alrededor de la mesa. « ¡La vida es bella cuando estamos juntos! », dice la abuela con una sonrisa.
\end{textebox}

\textbf{Glossaire} : \esp{ingeniero} = ingénieur ; \esp{carretera} = route ; \esp{liceo} = lycée ; \esp{corrigiendo} = en train de corriger ; \esp{mediano/a} = cadet(te) (du milieu) ; \esp{travieso} = espiègle ; \esp{ladrillo} = brique ; \esp{viudo/a} = veuf / veuve ; \esp{reunido} = réuni ; \esp{sonrisa} = sourire.

\textbf{Questions de compréhension} :
\begin{enumerate}
\item ¿Por qué está la familia en Yaoundé y no en Douala ?
\item ¿Por qué está cansada la madre ?
\item ¿Por qué está triste Carmen esta tarde ?
\item Identifiez trois emplois de \esp{ser} et trois emplois de \esp{estar} dans le texte et justifiez-les.
\end{enumerate}

\courssub{La règle de contraste et ses nuances}

\begin{aretenir}[SER = essence ; ESTAR = état]
\esp{ser} renvoie à ce qui définit durablement le sujet (essence, identité, qualité permanente) ; \esp{estar} renvoie à une circonstance, un état passager ou une localisation. Comparez :

\esp{Mi hermana es alta.} « Ma sœur est grande » (c'est sa taille, elle ne changera pas).

\esp{Mi hermana está delgada.} « Ma sœur est maigrie » (état actuel, résultat d'un changement).
\end{aretenir}

Mais attention : la nuance « permanent / temporaire » est un guide, pas une loi absolue. Certains emplois sont des \textbf{conventions de la langue} qu'il faut apprendre tels quels.

\begin{attention}[Piège n° 1 : la localisation d'un événement]
La localisation d'un \textbf{événement} (fête, réunion, match, concert) se dit avec \esp{ser} : \esp{La reunión es en el salón.} « La réunion est dans la salle. » \esp{El concierto es mañana en el estadio.} « Le concert est demain au stade. » En revanche, la localisation d'une \textbf{personne ou d'une chose} exige \esp{estar} : \esp{El salón está al fondo.} « La salle est au fond. »
\end{attention}

\begin{attention}[Piège n° 2 : l'état civil]
Erreur fréquente des francophones : \esp{*es casado}. Il faut dire \esp{está casado} « il est marié » : l'état civil est conçu comme un \textbf{état} résultant d'un changement (on n'est pas marié à la naissance). De même : \esp{está viudo} « il est veuf », \esp{está divorciado} « il est divorcé ». Exception admise : \esp{es soltero} « il est célibataire » avec \esp{ser} (le célibat est souvent perçu comme une identité), bien que \esp{está soltero} soit aussi possible.
\end{attention}

\begin{attention}[Piège n° 3 : calquer le français]
Le français « être » couvre tout ; l'espagnol vous oblige à \textbf{choisir} à chaque phrase. Avant de traduire, posez-vous la question : est-ce que je dis qui la personne \textbf{est} (identité) ou comment / où elle \textbf{se trouve} (état, lieu) ? Exemple : « La porte est ouverte » ne se traduit jamais par \esp{*la puerta es abierta} mais par \esp{la puerta está abierta} (résultat d'une action).
\end{attention}

\begin{center}
\begin{tabularx}{\linewidth}{|X|X|X|}
\hline
\rowcolor{popblueL} \textbf{Français} & \textbf{Español} & \textbf{Remarque} \\
\hline
Mon père est médecin. & Mi padre es médico. & profession : SER \\
\hline
Mon père est à l'hôpital. & Mi padre está en el hospital. & lieu : ESTAR \\
\hline
Il est cinq heures. & Son las cinco. & heure : SER \\
\hline
Nous sommes en train d'étudier. & Estamos estudiando. & action en cours : ESTAR \\
\hline
La réunion est dans la salle. & La reunión es en el salón. & événement : SER \\
\hline
La salle est au fond. & El salón está al fondo. & chose : ESTAR \\
\hline
Il est marié. & Está casado. & état civil : ESTAR \\
\hline
Elle est camerounaise. & Es camerunesa. & nationalité : SER \\
\hline
\end{tabularx}
\end{center}

\courssub{Carte mentale récapitulative}

\begin{popfigure}
\begin{tikzpicture}[mindmap, grow cyclic, every node/.style={concept, text=white, align=center},
root concept/.append style={concept color=popdark, font=\bfseries},
level 1/.append style={level distance=3.6cm, sibling angle=180},
level 2/.append style={level distance=2.6cm, sibling angle=45}]
\node[root concept]{SER / ESTAR}
  child[concept color=popblue] { node[align=center]{SER\\lo que eres}
    child[concept color=popblue!70] { node[align=center]{identidad\\nacionalidad} }
    child[concept color=popblue!70] { node[align=center]{origen\\profesión} }
    child[concept color=popblue!70] { node {hora y fecha} }
    child[concept color=popblue!70] { node[align=center]{cualidad\\permanente} }
    child[concept color=popblue!70] { node {eventos} }
  }
  child[concept color=poporange] { node[align=center]{ESTAR\\cómo y dónde estás}
    child[concept color=poporange!75] { node {lugar} }
    child[concept color=poporange!75] { node[align=center]{estado\\temporal} }
    child[concept color=poporange!75] { node[align=center]{estar +\\gerundio} }
  };
\end{tikzpicture}
\legende{Carte mentale : les emplois fondamentaux de SER et ESTAR.}
\end{popfigure}

\begin{savaistu}[Moyen mnémotechnique]
Retenez cette petite formule espagnole : \esp{SER = lo que eres} « ce que tu es » ; \esp{ESTAR = cómo y dónde estás} « comment et où tu es ». Si vous pouvez remplacer « être » par « se trouver » ou « se sentir » en français, c'est presque toujours \esp{estar} : « il est (se trouve) à Douala », « elle est (se sent) fatiguée ». Sinon, c'est généralement \esp{ser}.
\end{savaistu}

\courssub{Exercice résolu}

\begin{exoresolu}[SER ou ESTAR ? Justifiez votre choix]
Énoncé. Complétez avec la forme correcte de \esp{ser} ou \esp{estar} au présent, puis traduisez :
\begin{enumerate}
\item Mi tío \ldots\ profesor en el liceo de Bafoussam.
\item Hoy mi tío \ldots\ enfermo y \ldots\ en casa.
\item La fiesta de mi prima \ldots\ el sábado en el salón municipal.
\item Nosotros \ldots\ preparando la cena en este momento.
\item ¿De dónde \ldots\ tú? — Yo \ldots\ de Buea.
\item Mi hermano mayor \ldots\ casado desde hace dos años.
\end{enumerate}
\tcblower
Solution détaillée :
\begin{enumerate}
\item \esp{Mi tío \textbf{es} profesor en el liceo de Bafoussam.} « Mon oncle est professeur au lycée de Bafoussam. » Profession : \esp{ser}.
\item \esp{Hoy mi tío \textbf{está} enfermo y \textbf{está} en casa.} « Aujourd'hui, mon oncle est malade et il est à la maison. » État temporaire et localisation : \esp{estar}.
\item \esp{La fiesta de mi prima \textbf{es} el sábado en el salón municipal.} « La fête de ma cousine est samedi dans la salle municipale. » Événement (date et lieu) : \esp{ser}.
\item \esp{Nosotros \textbf{estamos} preparando la cena en este momento.} « Nous sommes en train de préparer le dîner en ce moment. » Action en cours : \esp{estar} + gérondif.
\item \esp{¿De dónde \textbf{eres} tú? — Yo \textbf{soy} de Buea.} « D'où es-tu ? — Je suis de Buea. » Origine : \esp{ser de}.
\item \esp{Mi hermano mayor \textbf{está} casado desde hace dos años.} « Mon frère aîné est marié depuis deux ans. » État civil : \esp{estar} (jamais \esp{*es casado}).
\end{enumerate}
\end{exoresolu}

\begin{exercice}[À vous de jouer]
Traduisez en espagnol : 1. Ma grand-mère est très généreuse. 2. Aujourd'hui elle est fatiguée et elle est dans sa chambre. 3. Il est huit heures et nous sommes en train de réviser la leçon. 4. Le match est au stade de Yaoundé. 5. Cette table est en bois et elle est près de la fenêtre. (Corrigé à la fin de la leçon.)
\end{exercice}

\begin{corrige}[Exercice « À vous de jouer »]
1. \esp{Mi abuela es muy generosa.} (qualité permanente : \esp{ser}) 2. \esp{Hoy está cansada y está en su habitación.} (état + lieu : \esp{estar}) 3. \esp{Son las ocho y estamos repasando la lección.} (heure : \esp{ser} ; action en cours : \esp{estar} + gérondif) 4. \esp{El partido es en el estadio de Yaoundé.} (événement : \esp{ser}) 5. \esp{Esta mesa es de madera y está cerca de la ventana.} (matière : \esp{ser} ; localisation : \esp{estar}).
\end{corrige}

La maîtrise de ces emplois fondamentaux acquise, nous pourrons aborder dans la section suivante un point plus délicat : les adjectifs qui \textbf{changent de sens} selon qu'ils sont employés avec \esp{ser} ou avec \esp{estar} (\esp{ser listo} / \esp{estar listo}, \esp{ser aburrido} / \esp{estar aburrido}...), puis nous appliquerons tout cela à la traduction guidée.

\coursec{SER/ESTAR + adjectif : changements de sens et expressions figées}

Dans la section précédente, nous avons posé les fondations : \esp{ser} exprime l'essence, l'identité, la définition durable ; \esp{estar} exprime l'état, la localisation, la circonstance passagère ou le résultat d'un changement. Nous allons maintenant voir la conséquence la plus spectaculaire de cette opposition : certains adjectifs \cle{changent de sens} selon le verbe qui les accompagne. C'est un point capital du programme de Terminale et un grand classique du Baccalauréat.

\courssub{Observation : un même adjectif, deux réalités}

\begin{textebox}[Dialogue : En la casa de los Mbarga, Yaoundé]
— Mamá, ¿ya está lista la cena? Tengo mucha hambre. \\
— Sí, hija, la sopa está lista, pero cuidado: está muy caliente. Espera un momento. \\
— ¡Qué bien! Tu sopa está buenísima, mamá. Eres la mejor cocinera del barrio. \\
— Gracias, cariño. Oye, ¿dónde está tu hermano? Estoy cansada de llamarlo. \\
— Está en su cuarto. Hoy está muy aburrido porque no hay electricidad y no puede ver el fútbol. \\
— Ese chico es muy listo para los estudios, pero es un poco vago para ayudar en casa... \\
— ¡Mamá! No seas mala con él. Además, hoy está malo, tiene fiebre. \\
— ¿Malo? ¡Ay, Dios mío! Voy a verlo enseguida. Tú pon la mesa, que ya es tarde y tu padre está de mal humor cuando espera.
\end{textebox}

\textbf{Glossaire :} \esp{estar listo} : être prêt ; \esp{cuidado} : attention ; \esp{cocinera} : cuisinière ; \esp{vago} : paresseux ; \esp{tener fiebre} : avoir de la fièvre ; \esp{poner la mesa} : mettre la table ; \esp{enseguida} : tout de suite ; \esp{de mal humor} : de mauvaise humeur.

\textbf{Questions d'observation :}
\begin{enumerate}
\item Relevez toutes les occurrences de \esp{ser} et \esp{estar} suivies d'un adjectif ou d'un participe.
\item Pourquoi dit-on \esp{la sopa está lista} mais \esp{ese chico es muy listo} ?
\item Pourquoi la mère dit-elle \esp{está malo} et non \esp{es malo} pour parler de la fièvre ?
\item Analysez \esp{está buenísima} (superlatif) et \esp{es tarde} (heure).
\end{enumerate}

On observe que \esp{listo}, \esp{bueno}, \esp{malo}, \esp{aburrido} apparaissent tantôt avec \esp{ser}, tantôt avec \esp{estar}, et que le sens change radicalement. Voilà la règle qui explique tout.

\courssub{Le principe fondamental}

\begin{propriete}[SER + adjectif / ESTAR + adjectif]
\begin{itemize}
\item \cle{SER + adjectif} qualifie la \textbf{nature durable} d'une personne ou d'une chose : un trait de caractère, une qualité inhérente, une définition, une classification. \esp{Ana es alegre.} « Ana est joyeuse (de caractère). »
\item \cle{ESTAR + adjectif} décrit un \textbf{état}, un \textbf{résultat} ou une \textbf{apparence} : quelque chose de circonstanciel, de temporaire, de subjectif ou de perçu à un moment précis. \esp{Ana está alegre hoy.} « Ana est joyeuse aujourd'hui. »
\end{itemize}
Conséquence : avec certains adjectifs, ce contraste est si fort que le français doit utiliser \textbf{deux mots différents} pour traduire.
\end{propriete}

\begin{exemplebox}[Le contraste en action]
\esp{La sopa está fría.} « La soupe est froide. » (état : elle a refroidi) \\
\esp{Mi tío es frío.} « Mon oncle est froid. » (caractère : il est distant, sans chaleur humaine) \\
\esp{El hielo es frío.} « La glace est froide. » (qualité inhérente : la glace est froide par nature)
\end{exemplebox}

\courssub{Les paires d'adjectifs à connaître par cœur}

\begin{center}
\begin{tabularx}{\linewidth}{|l|X|X|}
\hline
\rowcolor{popblueL} \textbf{Adjectif} & \textbf{SER + adj (nature, essence)} & \textbf{ESTAR + adj (état, résultat, apparence)} \\
\hline
\esp{listo} & intelligent, futé : \esp{Mi hermana es muy lista.} « Ma sœur est très intelligente. » & prêt, préparé : \esp{¿Estás listo? Salimos.} « Tu es prêt ? On part. » \\
\hline
\esp{aburrido} & ennuyeux (qui ennuie) : \esp{La película es aburrida.} « Le film est ennuyeux. » & qui s'ennuie (qui subit l'ennui) : \esp{Estoy aburrido en casa.} « Je m'ennuie à la maison. » \\
\hline
\esp{bueno} & bon (moralement, de nature) : \esp{Es una persona buena.} « C'est une personne bonne. » & bon (au goût), savoureux ; en bon état : \esp{Este plátano está bueno.} « Ce plantain est bon (délicieux). » \\
\hline
\esp{malo} & méchant, mauvais (moralement) : \esp{El villano es malo.} « Le méchant est mauvais. » & malade ; gâté, périmé (aliment) : \esp{Está malo, tiene fiebre.} « Il est malade. » \esp{La leche está mala.} « Le lait est tourné. » \\
\hline
\esp{vivo} & vif, éveillé, intelligent : \esp{Es un niño muy vivo.} « C'est un enfant très vif. » & vivant, en vie : \esp{Su abuelo todavía está vivo.} « Son grand-père est encore vivant. » \\
\hline
\esp{seguro} & sûr, fiable (lieu, chose) : \esp{Este barrio es seguro.} « Ce quartier est sûr. » & certain, convaincu (personne) : \esp{Estoy seguro de que vendrá.} « Je suis certain qu'il viendra. » \\
\hline
\esp{cansado} & fatigant (qui fatigue) : \esp{El viaje es cansado.} « Le voyage est fatigant. » & fatigué (qui subit la fatigue) : \esp{Estamos cansados después del partido.} « Nous sommes fatigués après le match. » \\
\hline
\esp{orgulloso} & orgueilleux, hautain (défaut) : \esp{Es un hombre orgulloso.} « C'est un homme orgueilleux. » & fier (sentiment positif) : \esp{Estoy orgulloso de ti.} « Je suis fier de toi. » \\
\hline
\esp{guapo} & beau, bel homme/belle femme (physique durable) : \esp{Es una chica guapa.} « C'est une belle fille. » & beau, bien mis (apparence du moment) : \esp{Estás guapa esta noche.} « Tu es belle ce soir. » \\
\hline
\esp{verde} & vert (couleur) : \esp{El campo es verde.} « La campagne est verte. » & pas mûr (fruit) : \esp{El mango está verde.} « La mangue n'est pas mûre. » \\
\hline
\esp{rico} & riche (argent) : \esp{Es un hombre rico.} « C'est un homme riche. » & délicieux, savoureux : \esp{El postre está riquísimo.} « Le dessert est délicieux. » \\
\hline
\esp{pálido} & pâle (teint habituel) : \esp{Ella es muy pálida.} « Elle est très pâle (de nature). » & pâle (signe de malaise, émotion) : \esp{Está pálido, ¿le duele algo?} « Il est pâle, a-t-il mal quelque part ? » \\
\hline
\end{tabularx}
\end{center}

\begin{popfigure}
\begin{tikzpicture}[mindmap, grow cyclic, every node/.style={concept, align=center}, concept color=popblueL,
root concept/.append style={concept color=popblue, font=\bfseries, text=white},
level 1/.append style={level distance=3.5cm, sibling angle=90},
level 2/.append style={level distance=2.8cm, sibling angle=50}]
\node[root concept] {SER / ESTAR + adjectif}
  child[concept color=popgreenL] { node {SER = Nature durable}
    child { node[concept color=popgreenL, align=center]{es listo\\ (intelligent)} }
    child { node[concept color=popgreenL, align=center]{es aburrido\\ (ennuyeux)} }
    child { node[concept color=popgreenL, align=center]{es bueno\\ (bon/moral)} }
  }
  child[concept color=poporangeL] { node {ESTAR = État / Résultat}
    child { node[concept color=poporangeL, align=center]{está listo\\ (prêt)} }
    child { node[concept color=poporangeL, align=center]{está aburrido\\ (s'ennuie)} }
    child { node[concept color=poporangeL, align=center]{está bueno\\ (délicieux)} }
  }
  child[concept color=poppinkL] { node {ESTAR = Apparence}
    child { node[concept color=poppinkL, align=center]{está guapa\\ esta noche} }
    child { node[concept color=poppinkL, align=center]{está pálido\\ (maintenant)} }
    child { node[concept color=poppinkL, align=center]{está caliente\\ (le café)} }
  }
  child[concept color=popgoldL] { node {Expressions figées}
    child { node[concept color=popgoldL, align=center]{estar de\\ acuerdo} }
    child { node[concept color=popgoldL, align=center]{es tarde\\ (heure)} }
    child { node[concept color=popgoldL, align=center]{estar de\\ vacaciones} }
  };
\end{tikzpicture}
\legende{Carte mentale : les quatre grands cas d'emploi de SER et ESTAR devant un adjectif.}
\end{popfigure}

\courssub{Expressions figées avec ESTAR}

Un certain nombre d'expressions très fréquentes se construisent avec \esp{estar} suivi de la préposition \esp{de} ou d'un adjectif. Le français utilise « être », mais l'espagnol impose \esp{estar} car il s'agit d'états circonstanciels, de situations ou de positions.

\begin{definition}[Locutions avec ESTAR + préposition / adjectif]
\begin{itemize}
\item \esp{estar de acuerdo (con)} : être d'accord (avec). \esp{Estoy de acuerdo contigo.} « Je suis d'accord avec toi. »
\item \esp{estar de moda} : être à la mode. \esp{Los pantalones anchos están de moda.} « Les pantalons larges sont à la mode. »
\item \esp{estar de vacaciones} : être en vacances. \esp{En diciembre estamos de vacaciones.} « En décembre, nous sommes en vacances. »
\item \esp{estar de buen / mal humor} : être de bonne / mauvaise humeur. \esp{Hoy el profesor está de buen humor.} « Aujourd'hui, le professeur est de bonne humeur. »
\item \esp{estar de pie} : être debout. \esp{Los pasajeros están de pie en el autobús.} « Les passagers sont debout dans le bus. »
\item \esp{estar sentado / acostado} : être assis / couché. \esp{Está sentado en el sofá.} « Il est assis sur le canapé. »
\item \esp{estar enfermo} : être malade. \esp{Mi abuela está enferma.} « Ma grand-mère est malade. »
\item \esp{estar a favor de / en contra de} : être pour / contre. \esp{Estamos en contra de la violencia de género.} « Nous sommes contre les violences faites aux femmes. »
\item \esp{estar seguro de} : être sûr de. \esp{¿Estás seguro de tu respuesta?} « Es-tu sûr de ta réponse ? »
\end{itemize}
\end{definition}

\courssub{Expressions figées avec SER}

À l'inverse, certaines constructions impersonnelles, identitaires ou de classification exigent \esp{ser} :

\begin{definition}[Locutions avec SER]
\begin{itemize}
\item \esp{ser de + origine / possession / matière} : \esp{Soy de Yaoundé.} « Je suis de Yaoundé. » \esp{El libro es de María.} « Le livre est à María. » \esp{La mesa es de madera.} « La table est en bois. »
\item \esp{ser hora de + infinitif} : être l'heure de. \esp{Es hora de estudiar.} « Il est l'heure d'étudier. »
\item \esp{es que...} : « c'est que... » (pour justifier). \esp{No vine porque... es que estaba enfermo.} « Je ne suis pas venu parce que... c'est que j'étais malade. »
\item \esp{es decir} : « c'est-à-dire ». \esp{Es bilingüe, es decir, habla español y francés.} « Il est bilingue, c'est-à-dire qu'il parle espagnol et français. »
\item \esp{ser + profession / nationalité / religion} (sans article) : \esp{Es médico.} « Il est médecin. » \esp{Somos cameruneses.} « Nous sommes Camerounais. »
\end{itemize}
\end{definition}

\courssub{Participe passé comme adjectif : ESTAR + participio}

C'est une extension majeure de la valeur « résultat » de \esp{estar}. Le participe passé (accordé en genre et nombre) fonctionne comme un adjectif d'état résultant d'une action achevée.

\begin{propriete}[Estar + participio = état résultant]
\begin{itemize}
\item L'action est terminée, on en voit le résultat présent.
\item \esp{La puerta está abierta.} (Quelqu'un l'a ouverte $\rightarrow$ résultat : elle est ouverte).
\item \esp{El libro está escrito en español.} (Résultat : il est écrit en espagnol).
\item \esp{Estoy cansado.} (Participe lexicalisé de \esp{cansar}, mais structure identique).
\end{itemize}
\textbf{Attention :} Ne pas confondre avec la voix passive (\esp{ser + participio} = action subie par le sujet). \esp{La puerta es abierta por el portero} (action habituelle) vs \esp{La puerta está abierta} (état actuel).
\end{propriete}

\begin{exemplebox}[Participe passé : SER (passif) vs ESTAR (état)]
\esp{El museo está cerrado los lunes.} (État : il est fermé aujourd'hui, on le voit). \\
\esp{El museo es cerrado por el director a las 20h.} (Passif : action de fermer, faite par le directeur). \\
\esp{La cena está preparada.} (Résultat : elle est prête sur la table). \\
\esp{La cena es preparada por mi madre.} (Passif : c'est ma mère qui la prépare d'habitude).
\end{exemplebox}

\courssub{Pièges de traduction et erreurs des francophones}

\begin{attention}[« Il est tard » : attention à l'heure]
Pour parler de l'heure, l'espagnol utilise \textbf{toujours} \esp{ser} : \esp{Es tarde.} « Il est tard. » \esp{Es temprano.} « Il est tôt. » \esp{Son las ocho.} « Il est huit heures. » Ne dites jamais \esp{*está tarde}.
\end{attention}

\begin{attention}[État ou nature : le choix change tout]
\esp{El café está caliente.} « Le café est chaud. » (état : il vient d'être servi ; demain il sera froid). Mais \esp{El sol es caliente.} « Le soleil est chaud. » (par nature). De même : \esp{Está guapa esta noche.} « Elle est belle ce soir » (apparence du moment, maquillée, bien habillée) ; \esp{Es guapa.} « Elle est belle » (de nature, physiquement).
\end{attention}

\begin{attention}[\esp{estar muerto} : l'exception célèbre]
Même si la mort est définitive, on dit \esp{estar muerto} : \esp{Su bisabuelo está muerto.} « Son arrière-grand-père est mort. » \esp{estar} exprime ici le \textbf{résultat d'un changement} (le verbe \esp{morir}). Ne dites jamais \esp{*es muerto}. Même logique pour tous les participes d'état : \esp{estar roto} (cassé), \esp{estar casado} (marié), \esp{estar cerrado} (fermé), \esp{estar abierto} (ouvert), \esp{estar quemado} (brûlé).
\end{attention}

\begin{attention}[Le français « être » n'est pas toujours \esp{ser}]
\begin{itemize}
\item « Être en vacances » = \esp{estar de vacaciones} ; « être d'accord » = \esp{estar de acuerdo} ; « être de bonne humeur » = \esp{estar de buen humor} ; « être debout/assis/couché » = \esp{estar de pie/sentado/acostado}.
\item « Être malade » = \esp{estar enfermo} / \esp{estar malo}.
\item Inversement, « il est étudiant » = \esp{es estudiante} (profession, identité : \esp{ser} sans article). « Il est Camerounais » = \esp{es camerunés}.
\end{itemize}
\end{attention}

\begin{methode}[Choisir entre SER et ESTAR devant un adjectif]
\begin{enumerate}
\item \textbf{Identifiez l'adjectif.} Est-il dans la liste des "changeurs de sens" (\esp{listo, aburrido, bueno, malo, vivo, seguro, cansado, orgulloso, guapo, verde, rico, pálido...}) ?
\item \textbf{Testez la nature (SER) :} Peut-on dire « c'est un être/une chose \emph{intrinsèquement} [adjectif] » ? $\rightarrow$ \esp{ser}. Ex: \esp{Es listo} (il a de l'intelligence en lui).
\item \textbf{Testez l'état/résultat/apparence (ESTAR) :} Est-ce une condition \emph{actuelle}, \emph{temporaire}, le \emph{résultat d'une action} ou une \emph{impression visuelle} ? $\rightarrow$ \esp{estar}. Ex: \esp{Está listo} (il a fini ses préparatifs, il est prêt maintenant).
\item \textbf{Vérifiez les expressions figées :} \esp{estar de acuerdo, estar de vacaciones, ser de, es hora de, es que...}
\item \textbf{Vérifiez l'heure :} Toujours \esp{ser}.
\item \textbf{Vérifiez le participe passé :} État visible = \esp{estar} ; Action subie (passif) = \esp{ser}.
\end{enumerate}
\end{methode}

\courssub{Exercice résolu et entraînement}

\begin{exoresolu}[Thème : choisir entre SER et ESTAR]
Énoncé : Traduisez en espagnol.
1. Mon frère est très intelligent, mais aujourd'hui il n'est pas prêt pour l'examen.
2. Ce film est ennuyeux ; les enfants s'ennuient.
3. Je suis d'accord avec toi : la soupe est froide.
4. Elle est belle ce soir avec sa robe verte.
5. Il est tard, il est l'heure de rentrer.
\tcblower
Solution détaillée :
\begin{enumerate}
\item \esp{Mi hermano es muy listo, pero hoy no está listo para el examen.} (« intelligent » = nature : \esp{ser} ; « prêt » = état : \esp{estar}.)
\item \esp{Esta película es aburrida; los niños están aburridos.} (le film ennuie : qualité \esp{ser} ; les enfants subissent l'ennui : état \esp{estar}.)
\item \esp{Estoy de acuerdo contigo: la sopa está fría.} (expression figée \esp{estar de acuerdo} ; état de la soupe \esp{estar}.)
\item \esp{Está guapa esta noche con su vestido verde.} (apparence du moment : \esp{estar}.)
\item \esp{Es tarde, es hora de volver a casa.} (l'heure : toujours \esp{ser} ; expression figée \esp{ser hora de}.)
\end{enumerate}
\end{exoresolu}

\begin{exercice}[Thème : à vous de jouer]
Traduisez en espagnol :
1. Mon grand-père est mort en 2010 (état actuel).
2. Ce quartier est sûr, mais je ne suis pas certain du chemin.
3. Nous sommes contre la violence et nous sommes fiers de notre engagement.
4. Le lait est tourné ; ne le bois pas.
5. En août, ma famille est en vacances à Kribi.
6. La porte est ouverte (résultat) vs La porte est ouverte par le gardien (passif).
7. Ce gâteau est délicieux ! / Cet homme est riche.
\end{exercice}

\begin{exercice}[Version : traduisez en français]
1. \esp{Mi abuelo está muerto desde 2010.}
2. \esp{Este barrio es seguro, pero no estoy seguro del camino.}
3. \esp{Estamos en contra de la violencia y estamos orgullosos de nuestro compromiso.}
4. \esp{La leche está mala; no la bebas.}
5. \esp{En agosto, mi familia está de vacaciones en Kribi.}
6. \esp{La ventana está rota.} / \esp{La ventana es rota por el viento.}
7. \esp{El mango está verde.} / \esp{El pasto es verde.}
\end{exercice}

\begin{corrige}[Exercices : Thème et Version]
\textbf{Thème :}
1. \esp{Mi abuelo está muerto desde 2010.} (résultat d'un changement : \esp{estar} + participio). \textit{Note : pour l'événement « il est mort en 2010 », on dirait \esp{Mi abuelo murió en 2010}.}
2. \esp{Este barrio es seguro, pero no estoy seguro del camino.}
3. \esp{Estamos en contra de la violencia y estamos orgullosos de nuestro compromiso.}
4. \esp{La leche está mala; no la bebas.} (subjonctif \esp{bebas} pour impératif négatif).
5. \esp{En agosto, mi familia está de vacaciones en Kribi.}
6. \esp{La puerta está abierta.} (état) / \esp{La puerta es abierta por el portero.} (passif).
7. \esp{¡Este pastel está riquísimo! / Este hombre es rico.}

\textbf{Version :}
1. Mon grand-père est mort depuis 2010. (État résultant).
2. Ce quartier est sûr, mais je ne suis pas certain du chemin.
3. Nous sommes contre la violence et nous sommes fiers de notre engagement.
4. Le lait est tourné ; ne le bois pas.
5. En août, ma famille est en vacances à Kribi.
6. La fenêtre est cassée (état). / La fenêtre est cassée par le vent (passif).
7. La mangue n'est pas mûre. / L'herbe est verte.
\end{corrige}

\begin{savaistu}[Un classique du Bac]
Au baccalauréat, la distinction \esp{ser}/\esp{estar} + adjectif est un grand classique des questions de grammaire et des exercices de traduction (thème/version). Les correcteurs attendent que vous sachiez les paires \esp{ser listo}/\esp{estar listo}, \esp{ser aburrido}/\esp{estar aburrido}, \esp{ser malo}/\esp{estar malo}, \esp{ser rico}/\esp{estar rico} et le cas \esp{estar muerto} / \esp{estar + participio}. 

\textbf{Méthode efficace :} Apprenez chaque paire par cœur avec \textbf{un exemple complet traduit}, et non avec une simple liste de mots. C'est la phrase entière (\esp{« La sopa está buena »} vs \esp{« La niña es buena »}) qui vous reviendra en mémoire le jour de l'épreuve. Entraînez-vous à justifier votre choix en une phrase : « \esp{Ser} car c'est une qualité intrinsèque » ou « \esp{Estar} car c'est un état résultant/apparence ».
\end{savaistu}

\coursec{Traduction guidée : thème et version}

Dans la section précédente, nous avons vu que \esp{ser} et \esp{estar} peuvent changer le sens d'un adjectif : \esp{ser aburrido} (« être ennuyeux ») ne veut pas dire la même chose que \esp{estar aburrido} (« s'ennuyer »). Il est temps maintenant de mettre ce savoir en pratique dans l'exercice roi de l'évaluation en Terminale A : la \cle{traduction}. Traduire le verbe français « être » en espagnol est l'une des opérations les plus délicates du thème, car le français n'a qu'un seul verbe là où l'espagnol en a deux. Cette section vous donne une méthode infaillible en quatre étapes, puis vous entraîne à la version et au thème, en contexte, avec une correction collective argumentée.

\courssub{Observer avant de traduire : quatre phrases, quatre réalités}

Lisons d'abord ces quatre phrases françaises. Elles contiennent toutes le même verbe « être », et pourtant\ldots

\begin{exemplebox}[Quatre « être » français, quatre traductions espagnoles]
\begin{itemize}
\item \esp{Él es profesor.} « Il est professeur. » (identité, profession)
\item \esp{Él está en la escuela.} « Il est à l'école. » (lieu)
\item \esp{Él está enfermo.} « Il est malade. » (état passager)
\item \esp{Son las ocho.} « Il est huit heures. » (heure)
\end{itemize}
\end{exemplebox}

Observation : le français dit « il est » partout ; l'espagnol, lui, \cle{classe la réalité} avant de parler. Est-ce une définition permanente ? \esp{ser}. Est-ce une localisation ou un état temporaire ? \esp{estar}. Traduire « être », c'est donc d'abord \emph{penser}, ensuite seulement écrire.

\begin{methode}[Traduire « être » en quatre étapes]
\begin{enumerate}
\item \textbf{Repérer le sens de « être »} dans la phrase française : s'agit-il d'une identité, d'une définition, d'une origine, d'une profession, de l'heure ? Ou bien d'un lieu, d'un état physique ou moral, d'une action en cours ?
\item \textbf{Choisir le verbe} : identité / origine / heure / possession / caractéristique essentielle $\rightarrow$ \esp{ser} ; lieu / état passager / action en cours $\rightarrow$ \esp{estar}.
\item \textbf{Vérifier l'accord de l'adjectif} qui suit : genre et nombre avec le sujet (\esp{Mi madre está cansada}, \esp{Mis hermanos son altos}).
\item \textbf{Relire} la phrase espagnole complète : conjugaison correcte ? accents ? ponctuation ¿ ? ¡ ! ?
\end{enumerate}
\end{methode}

\begin{popfigure}
\begin{tikzpicture}[
  node distance=0.9cm and 1.4cm,
  startstop/.style={rectangle, rounded corners, draw=popdark, fill=popblueL, text width=3.2cm, align=center, minimum height=1cm, font=\small\bfseries},
  decision/.style={diamond, draw=popdark, fill=popgoldL, text width=3.4cm, align=center, inner sep=1pt, font=\small},
  result/.style={rectangle, rounded corners, draw=popdark, fill=popgreenL, text width=3.4cm, align=center, minimum height=1cm, font=\small\bfseries},
  arrow/.style={-{Stealth[length=3mm]}, thick, popdark}
]
\node[startstop] (start) {Le français dit « être »};
\node[decision, below=of start] (q1) {Identité ? origine ? heure ? profession ?};
\node[result, right=of q1, align=center] (ser){SER\\ \esp{Es profesor. Son las ocho.}};
\node[decision, below=of q1] (q2) {Lieu ? état passager ? action en cours ?};
\node[result, right=of q2, align=center] (estar){ESTAR\\ \esp{Está enfermo. Está en casa.}};
\node[result, below=of q2, fill=poporangeL] (accord) {Accorder l'adjectif puis relire};
\draw[arrow] (start) -- (q1);
\draw[arrow] (q1) -- node[above, font=\small]{oui} (ser);
\draw[arrow] (q1) -- node[left, font=\small]{non} (q2);
\draw[arrow] (q2) -- node[above, font=\small]{oui} (estar);
\draw[arrow] (q2) -- (accord);
\end{tikzpicture}
\legende{Organigramme de décision : comment traduire « être » en espagnol}
\end{popfigure}

\courssub{Version guidée : de l'espagnol vers le français}

En \cle{version}, le travail est inversé : l'espagnol a déjà choisi entre \esp{ser} et \esp{estar}, et c'est à vous de \emph{justifier} ce choix pour bien comprendre le texte avant de le traduire. Cette justification est souvent demandée à l'examen : elle prouve que vous avez compris, pas seulement décodé.

\begin{exoresolu}[Version guidée : traduire et justifier]
Traduisez en français et justifiez l'emploi de \esp{ser} ou \esp{estar} :
\begin{enumerate}
\item \esp{Mi cuñada es muy generosa.}
\item \esp{La sopa está fría.}
\item \esp{El debate sobre la violencia de género es importante.}
\item \esp{Estamos muy preocupados por la situación.}
\end{enumerate}
\tcblower
\textbf{Solution détaillée :}
\begin{enumerate}
\item « Ma belle-sœur est très généreuse. » $\rightarrow$ \esp{ser} : trait de caractère permanent, qualité essentielle de la personne.
\item « La soupe est froide. » $\rightarrow$ \esp{estar} : état passager (la soupe peut être réchauffée ; ce n'est pas sa nature).
\item « Le débat sur les violences faites aux femmes est important. » $\rightarrow$ \esp{ser} : jugement qui définit le débat, caractéristique essentielle.
\item « Nous sommes très inquiets face à la situation. » $\rightarrow$ \esp{estar} : état émotionnel passager ; attention à l'accord pluriel \esp{preocupados}.
\end{enumerate}
\end{exoresolu}

\begin{attention}[Faux amis et calques en version]
\begin{itemize}
\item \esp{estar constipado} signifie « être enrhumé », et non « être constipé » !
\item \esp{ser largo} = « être long » ; « être large » se dit \esp{ser ancho}.
\item \esp{estar embarazada} = « être enceinte » (et non « être embarrassée »).
\item Le français « il est tard » se traduit \esp{es tarde} (l'heure relève de \esp{ser}), jamais \esp{está tarde}.
\end{itemize}
\end{attention}

\courssub{Thème guidé : du français vers l'espagnol, pas à pas}

Appliquons maintenant la méthode des quatre étapes à un thème complet, phrase par phrase.

\begin{exoresolu}[Thème guidé : le portrait de ma sœur]
Traduisez : « Ma sœur est infirmière. Elle est à l'hôpital aujourd'hui. Elle est fatiguée mais elle est contente de son travail. »
\tcblower
\textbf{Analyse phrase par phrase :}
\begin{enumerate}
\item « Ma sœur \textbf{est} infirmière » : profession, identité $\rightarrow$ \esp{ser}. Pas d'article indéfini devant la profession en espagnol. $\rightarrow$ \esp{Mi hermana \textbf{es} enfermera.}
\item « Elle \textbf{est} à l'hôpital aujourd'hui » : lieu, localisation $\rightarrow$ \esp{estar}. $\rightarrow$ \esp{Hoy \textbf{está} en el hospital.}
\item « Elle \textbf{est} fatiguée » : état physique passager $\rightarrow$ \esp{estar} ; accord féminin : \esp{cansada}. $\rightarrow$ \esp{\textbf{Está} cansada}
\item « mais elle \textbf{est} contente de son travail » : état moral passager $\rightarrow$ \esp{estar} ; accord féminin : \esp{contenta}. $\rightarrow$ \esp{pero \textbf{está} contenta con su trabajo.}
\end{enumerate}
\textbf{Traduction complète :} \esp{Mi hermana es enfermera. Hoy está en el hospital. Está cansada pero está contenta con su trabajo.}
\end{exoresolu}

Remarquez la beauté de l'exercice : le français répète quatre fois « est », et l'espagnol répond une fois \esp{es} (identité) puis trois fois \esp{está} (lieu et états). C'est exactement le raisonnement attendu à l'examen.

\begin{attention}[L'accord de l'adjectif : la faute qui coûte cher]
Après \esp{ser} comme après \esp{estar}, l'adjectif s'accorde en genre et en nombre avec le sujet. Les francophones oublient souvent la marque du féminin :
\begin{itemize}
\item \esp{Mi madre está cansad\textbf{a}.} « Ma mère est fatiguée. » (et non *\esp{cansado})
\item \esp{Las chicas son muy trabajadoras.} « Les filles sont très travailleuses. »
\item \esp{El padre está preocupado y la hija está preocupada.} « Le père est inquiet et la fille est inquiète. »
\end{itemize}
\end{attention}

\begin{propriete}[« Être en train de » = ESTAR + gérondif]
Pour exprimer une \cle{action en cours}, le français utilise « être en train de + infinitif ». L'espagnol dit \esp{estar} + \cle{gérondif} (en \esp{-ando} pour les verbes en \esp{-ar}, en \esp{-iendo} pour les verbes en \esp{-er}/\esp{-ir}) :
\begin{itemize}
\item \esp{Estamos debatiendo.} « Nous sommes en train de débattre. »
\item \esp{Ella está cocinando.} « Elle est en train de cuisiner. »
\item \esp{¿Qué estás haciendo?} « Qu'es-tu en train de faire ? »
\end{itemize}
Gérondifs irréguliers à connaître : \esp{leer} $\rightarrow$ \esp{leyendo}, \esp{dormir} $\rightarrow$ \esp{durmiendo}, \esp{decir} $\rightarrow$ \esp{diciendo}, \esp{ir} $\rightarrow$ \esp{yendo}.
\end{propriete}

\courssub{Thème en contexte : retour au dialogue familial}

Replaçons maintenant la traduction dans le contexte de notre module : la vie familiale et le débat sur les violences faites aux femmes. Voici des phrases inspirées du dialogue familial étudié dans cette leçon.

\begin{exercice}[Thème en contexte : dans la famille de Ndoumbe]
Traduisez en espagnol en justifiant chaque choix de verbe :
\begin{enumerate}
\item Mon père est mécanicien à Douala.
\item La réunion de famille est à huit heures du soir.
\item Ma mère est très en colère aujourd'hui.
\item Nous sommes en train de discuter d'un problème grave.
\item Les violences faites aux femmes sont inacceptables.
\item Ta tante est à l'hôpital parce qu'elle est malade.
\item Ce débat est très utile pour les jeunes.
\item Mes cousins sont en train d'écouter le débat à la radio.
\end{enumerate}
\end{exercice}

\begin{corrige}[Thème en contexte : dans la famille de Ndoumbe]
\begin{enumerate}
\item \esp{Mi padre es mecánico en Douala.} $\rightarrow$ profession : \esp{ser} ; pas d'article devant le métier.
\item \esp{La reunión familiar es a las ocho de la tarde.} $\rightarrow$ heure d'un événement : \esp{ser}.
\item \esp{Mi madre está muy enfadada hoy.} $\rightarrow$ état passager : \esp{estar} ; accord féminin \esp{enfadada}.
\item \esp{Estamos discutiendo un problema grave.} $\rightarrow$ action en cours : \esp{estar} + gérondif.
\item \esp{La violencia contra las mujeres es inaceptable.} $\rightarrow$ jugement, définition : \esp{ser}.
\item \esp{Tu tía está en el hospital porque está enferma.} $\rightarrow$ lieu puis état : deux fois \esp{estar}.
\item \esp{Este debate es muy útil para los jóvenes.} $\rightarrow$ caractéristique essentielle : \esp{ser}.
\item \esp{Mis primos están escuchando el debate en la radio.} $\rightarrow$ action en cours : \esp{estar} + gérondif.
\end{enumerate}
\end{corrige}

\courssub{Correction collective : justifier pour progresser}

La \cle{correction collective} est un moment capital : chaque élève doit pouvoir dire \emph{pourquoi} il a choisi \esp{ser} ou \esp{estar}, pas seulement donner la bonne réponse. Voici le tableau récapitulatif qui sert de référence pendant la correction.

\begin{center}
\begin{tabularx}{\linewidth}{|X|X|X|}
\hline
\rowcolor{popblueL} Sens de « être » en français & Verbe espagnol & Exemple traduit \\
\hline
Identité, définition, nationalité & SER & \esp{Soy camerunés.} « Je suis camerounais. » \\
\hline
Profession, religion, fonction & SER & \esp{Es enfermera.} « Elle est infirmière. » \\
\hline
Heure, date, prix & SER & \esp{Son las ocho.} « Il est huit heures. » \\
\hline
Possession, matière, origine & SER & \esp{El anillo es de oro.} « La bague est en or. » \\
\hline
Caractère, qualité essentielle & SER & \esp{Es muy generosa.} « Elle est très généreuse. » \\
\hline
Lieu, localisation & ESTAR & \esp{Está en el mercado.} « Elle est au marché. » \\
\hline
État physique ou moral passager & ESTAR & \esp{Está cansada.} « Elle est fatiguée. » \\
\hline
Action en cours (+ gérondif) & ESTAR & \esp{Está hablando.} « Elle est en train de parler. » \\
\hline
Résultat d'un changement & ESTAR & \esp{La sopa está fría.} « La soupe est froide. » \\
\hline
\end{tabularx}
\end{center}

\begin{experience}[Correction collective argumentée]
Le professeur écrit au tableau huit phrases à traduire. Par groupes de quatre, les élèves traduisent, puis chaque groupe vient justifier \emph{oralement} un choix de verbe devant la classe, en complétant la phrase rituelle : \esp{Uso «ser/estar» porque es una cuestión de\ldots} (« J'utilise ser/estar parce qu'il s'agit de\ldots »). Les autres groupes valident ou corrigent. Le groupe qui justifie le plus de choix correctement gagne le titre de \esp{campeón de la traducción}.
\end{experience}

\courssub{Bilan : ma grille d'auto-évaluation}

Terminez la leçon en remplissant cette grille. Si vous pouvez cocher chaque case et donner un exemple, vous êtes prêt pour l'épreuve de traduction.

\begin{aretenir}[Je choisis SER quand\ldots / Je choisis ESTAR quand\ldots]
\begin{itemize}
\item Je choisis \esp{SER} quand « être » définit une \cle{identité}, une profession, une origine, l'heure, une possession ou une qualité permanente. Exemple : \esp{Mi hermana es enfermera.}
\item Je choisis \esp{ESTAR} quand « être » indique un \cle{lieu}, un état passager ou une action en cours. Exemple : \esp{Está en el hospital. Está cansada.}
\item Je vérifie toujours l'\cle{accord} de l'adjectif : \esp{cansada}, \esp{contenta}, \esp{preocupados}\ldots
\item Je me méfie des pièges : \esp{es tarde} (« il est tard »), \esp{estar constipado} (« être enrhumé »), \esp{estar + gérondif} pour « être en train de ».
\item Je sais justifier mon choix à l'oral : \esp{Uso «estar» porque es un estado pasajero.}
\end{itemize}
\end{aretenir}

En résumé : traduire « être », ce n'est pas chercher un mot dans un dictionnaire, c'est \cle{classer la réalité} comme le fait spontanément tout hispanophone. Identité ? \esp{ser}. Lieu ou état ? \esp{estar}. Action en cours ? \esp{estar} + gérondif. Avec l'organigramme en tête et la grille d'auto-évaluation sous les yeux, l'exercice de thème n'a plus de secret pour vous.

\newpage

\coursec{Activité d'intégration}

\courssub{Objectifs de l'activité}

À l'issue de cette activité d'intégration, l'élève sera capable de :
\begin{itemize}
\item mener un \cle{dialogue familial} en espagnol : prendre la parole, écouter activement, relancer un interlocuteur et reformuler ses propos ;
\item exprimer et défendre une \cle{opinion argumentée} sur les violences faites aux femmes, dans le respect de l'autre ;
\item employer correctement \cle{SER} et \cle{ESTAR} dans une production écrite authentique, y compris avec des adjectifs à changement de sens (\esp{ser listo / estar listo}) et des expressions figées (\esp{estar de acuerdo}, \esp{estar harto de}) ;
\item rédiger une courte \cle{carta de opinión} destinée à un journal scolaire ;
\item s'autoévaluer et corriger collectivement les erreurs de grammaire.
\end{itemize}

\courssub{Situation}

Tu es élève en Terminale A au lycée de Nkol-Eton, à Yaoundé. Ce matin, dans le bus qui te menait au lycée, tu as écouté à la radio un reportage sur les violences faites aux femmes au Cameroun. Le soir, à la maison, toute la famille en parle autour de la table : le grand-père, la mère, le père et la fille adolescente. Chacun a son avis... et les avis ne sont pas toujours les mêmes ! Le lendemain, le club de journalisme du lycée lance un appel à contributions pour son numéro spécial « \esp{No a la violencia de género} ». Avec ton groupe, tu décides d'écrire une lettre d'opinion.

Voici la transcription du début du reportage entendu à la radio :

\begin{textebox}[Reportage radio — « Voces de Camerún » (texte original)]
\textsc{Periodista} : Buenas tardes, oyentes. Hoy hablamos de un tema que está presente en muchas familias camerunesas : la violencia contra las mujeres. En los mercados de Yaundé, en las escuelas, en las casas... muchas mujeres sufren en silencio.\\
\textsc{Testigo} : Mi vecina está cansada de sufrir. Su marido es violento y ella tiene miedo de hablar.\\
\textsc{Periodista} : Las asociaciones de defensa de los derechos de la mujer piden más educación, más justicia y más valentía para denunciar. Porque la violencia no es una tradición : es un crimen.
\end{textebox}

\textbf{Glossaire :} \esp{oyentes} : auditeurs ; \esp{sufrir en silencio} : souffrir en silence ; \esp{estar cansada de} : être fatiguée de ; \esp{tener miedo de} : avoir peur de ; \esp{denunciar} : dénoncer ; \esp{un crimen} : un crime.

\courssub{Consignes}

Travail par groupes de 4 élèves. L'activité se déroule en cinq étapes guidées.

\begin{enumerate}
\item \textbf{Compréhension du document.} Lis le reportage radio à voix haute dans ton groupe. Réponds oralement : ¿\cle{De qué} habla el periodista ? ¿Qué problema tiene la vecina del testigo ? ¿Qué piden las asociaciones ? Souligne dans le texte tous les emplois de \esp{ser} et \esp{estar} et justifie chaque choix.
\item \textbf{Préparation du dialogue.} Répartissez les rôles : \esp{abuelo}, \esp{madre}, \esp{padre}, \esp{hija adolescente}. Chacun prépare deux ou trois idées (une opinion, un argument, un exemple) en utilisant le lexique de la leçon. Attention : les quatre personnages ne doivent pas être d'accord sur tout !
\item \textbf{Jeu de rôle : el debate en familia.} Jouez le dialogue (5 minutes minimum). Règles du jeu : chaque personnage doit \cle{relancer} au moins une fois un autre membre (\esp{¿Y tú qué piensas, papá?}) et \cle{reformuler} au moins une fois les propos d'un autre (\esp{Es decir que, según tú, ...} / \esp{¿Quieres decir que...?}). Utilisez les expressions d'accord et de désaccord.
\item \textbf{Production écrite : la carta de opinión.} Rédigez ensemble une lettre d'opinion de 8 à 10 lignes adressée au journal scolaire. Contraintes obligatoires : employer au moins \cle{5 fois} \esp{SER} ou \esp{ESTAR} correctement, dont \cle{2 cas} avec changement de sens ou expression figée ; utiliser au moins 6 mots du lexique de la violence de genre ; structurer la lettre (salutation, opinion, arguments, conclusion, formule de politesse).
\item \textbf{Restitution et correction collective.} Deux groupes jouent leur dialogue devant la classe. Les lettres sont affichées ou lues : la classe repère les emplois de \esp{SER}/\esp{ESTAR} et corrige ensemble les erreurs éventuelles au tableau.
\end{enumerate}

\courssub{Ressources à mobiliser}

\begin{aretenir}[Boîte à outils de l'élève]
\textbf{Écoute active et relance :} \esp{¿Y tú qué piensas?} ; \esp{¿Qué opinas tú, abuelo?} ; \esp{¿Estás de acuerdo conmigo?}\\
\textbf{Reformulation :} \esp{Es decir...} ; \esp{¿Quieres decir que...?} ; \esp{En otras palabras, ...} ; \esp{Si te entiendo bien, ...}\\
\textbf{Opinion et argumentation :} \esp{Pienso / Creo / Me parece que...} ; \esp{Estoy de acuerdo / No estoy de acuerdo} ; \esp{En mi opinión...} ; \esp{Porque..., ya que..., por eso...}\\
\textbf{Lexique :} \esp{la violencia de género}, \esp{los derechos de la mujer}, \esp{denunciar}, \esp{maltratar}, \esp{el maltrato}, \esp{la igualdad}, \esp{respetar}, \esp{proteger}, \esp{la víctima}, \esp{el agresor}, \esp{la justicia}.\\
\textbf{SER :} identité, profession, origine, caractéristiques permanentes, opinions générales (\esp{La violencia es un crimen}).\\
\textbf{ESTAR :} localisation, état temporaire, résultat d'un changement, expressions figées (\esp{estar de acuerdo}, \esp{estar harto de}, \esp{estar cansado de}).\\
\textbf{Changements de sens :} \esp{ser listo} (être intelligent) / \esp{estar listo} (être prêt) ; \esp{ser bueno} (être bon, gentil) / \esp{estar bueno} (être bon au goût, être en forme) ; \esp{ser aburrido} (être ennuyeux) / \esp{estar aburrido} (s'ennuyer).
\end{aretenir}

\begin{attention}[Pièges à éviter dans ta production]
Ne traduis jamais « être » mécaniquement ! « Elle est fatiguée » = \esp{Ella está cansada} (état temporaire), mais « Elle est courageuse » = \esp{Ella es valiente} (caractère). Attention aussi à l'accord de l'adjectif : \esp{Las mujeres están cansadas de sufrir}. Et n'oublie pas : \esp{estar de acuerdo} est une expression figée — jamais \esp{ser de acuerdo} !
\end{attention}

\courssub{Proposition de production et corrigé commenté}

\begin{exoresolu}[Modèle de « carta de opinión »]
Rédige avec ton groupe une lettre d'opinion de 8 à 10 lignes pour le journal scolaire, en respectant les contraintes de la consigne 4.
\tcblower
\textbf{Production modèle :}

\esp{Estimada redacción del periódico escolar :}

\esp{Somos alumnos de Terminale A y queremos denunciar la violencia de género en nuestra sociedad. En nuestra opinión, la violencia contra las mujeres es un crimen, no una tradición. Muchas mujeres están cansadas de sufrir en silencio y tienen miedo de hablar. Nosotros estamos hartos de esta injusticia : es hora de actuar. La escuela es un lugar donde todos debemos estar seguros y respetados. Además, los jóvenes somos el futuro del país : si somos valientes y estamos listos para denunciar, las cosas van a cambiar. Pedimos más educación, más justicia y más respeto por los derechos de la mujer.}

\esp{Atentamente,}

\esp{El club de español del Lycée de Nkol-Eton, Yaundé.}

\textbf{Commentaire des choix de langue :}
\begin{itemize}
\item \esp{Somos alumnos} : \esp{SER} pour l'identité et la qualité permanente.
\item \esp{la violencia ... es un crimen} : \esp{SER} pour une définition, un jugement général.
\item \esp{están cansadas de sufrir} : \esp{ESTAR} + expression figée \esp{estar cansado de} (état résultant d'une situation) ; accord féminin pluriel.
\item \esp{estamos hartos de} : expression figée avec \esp{ESTAR} (« nous en avons assez de »).
\item \esp{estar seguros} : état temporaire, condition.
\item \esp{somos el futuro} : \esp{SER} d'identité ; \esp{somos valientes} : trait de caractère.
\item \esp{estamos listos para denunciar} : changement de sens — \esp{estar listo} = « être prêt » (et non \esp{ser listo} = « être intelligent »).
\item Lexique mobilisé : \esp{denunciar}, \esp{violencia de género}, \esp{derechos de la mujer}, \esp{justicia}, \esp{respeto}, \esp{víctimas} (implicite), \esp{injusticia}.
\item Structure respectée : salutation, opinion, arguments, appel à l'action, formule de politesse.
\end{itemize}
\end{exoresolu}

\courssub{Bilan de l'activité}

Cette activité t'a permis de réinvestir l'ensemble des savoir-faire de la leçon dans une situation de communication complète et réaliste. Vérifie que tu sais désormais :

\begin{itemize}
\item[\textbf{Oral :}] prendre la parole dans un dialogue familial, relancer (\esp{¿Y tú qué piensas?}) et reformuler (\esp{Es decir...}) sans bloquer la conversation ;
\item[\textbf{Argumentation :}] défendre une opinion sur un sujet de société sensible avec respect et arguments (\esp{En mi opinión..., porque..., por eso...}) ;
\item[\textbf{Grammaire :}] choisir correctement entre \esp{SER} (identité, caractère, définition) et \esp{ESTAR} (état, lieu, changement, expressions figées), et expliquer ton choix ;
\item[\textbf{Écrit :}] rédiger une lettre d'opinion structurée, avec accords corrects et lexique thématique approprié.
\end{itemize}

\begin{experience}[Pour aller plus loin : mini-débat de classe]
Organisez un débat parlementaire en classe entière sur la motion : \esp{« La educación es la mejor arma contra la violencia de género »}. Un groupe défend la motion (\esp{a favor}), un autre la combat (\esp{en contra}), un troisième joue les journalistes et pose des questions. Le professeur compte les emplois corrects de \esp{SER} et \esp{ESTAR} : chaque emploi juste rapporte un point à l'équipe !
\end{experience}

\begin{savaistu}[Le saviez-vous ?]
En Guinée équatoriale, voisin hispanophone du Cameroun, l'espagnol est langue officielle : les débats sur les droits des femmes s'y déroulent aussi en espagnol. Le 25 novembre est la Journée internationale pour l'élimination de la violence à l'égard des femmes (\esp{Día Internacional de la Eliminación de la Violencia contra la Mujer}), en mémoire des sœurs Mirabal, assassinées en République dominicaine en 1960 pour leur combat pour la liberté.
\end{savaistu}

\newpage

\coursec{Méthodes et exercices résolus}

\coursec{Leçon E02 : Diálogo familiar, debate sobre la violencia de género y traducción : SER y ESTAR}

\begin{textebox}[Dialogue familial : un projet de vie à Yaoundé]
\esp{--- Mamá, quiero dejar los estudios y abrir una peluquería.}\\
\esp{--- ¿De verdad? ¡Pero hija, si eres la mejor de tu clase!}\\
\esp{--- Es que ya no me gusta el instituto. Me siento perdida.}\\
\esp{--- O sea que te sientes perdida... ¿Por qué no me lo dijiste antes?}\\
\esp{--- Porque siempre estás ocupada en el mercado.}\\
\esp{--- Entiendo... Tienes razón. Esta noche hablamos con calma, ¿vale?}\\
\esp{--- Vale, máma. Gracias por escucharme.}
\end{textebox}

\begin{definition}[Lexique thématique : La familia y la violencia de género]
\textbf{La familia :} \esp{los padres}, \esp{la madre}, \esp{el padre}, \esp{los hijos}, \esp{la hija}, \esp{el hijo}, \esp{los abuelos}, \esp{la abuela}, \esp{el abuelo}, \esp{los hermanos}, \esp{la hermana}, \esp{el hermano}, \esp{el tío}, \esp{la tía}, \esp{el primo}, \esp{la prima}, \esp{el esposo}, \esp{la esposa}, \esp{el matrimonio}, \esp{el divorcio}, \esp{la custodia}, \esp{la pensión alimenticia}, \esp{la convivencia}, \esp{el diálogo}, \esp{la confianza}, \esp{el respeto}, \esp{la autoridad}, \esp{el apoyo}, \esp{la comprensión}, \esp{el proyecto de vida}, \esp{los estudios}, \esp{el instituto}, \esp{la peluquería}, \esp{el mercado}, \esp{el bendskin}.

\textbf{La violencia de género :} \esp{la violencia de género}, \esp{el maltrato}, \esp{la violencia física}, \esp{la violencia psicológica}, \esp{la violencia económica}, \esp{la víctima}, \esp{el agresor}, \esp{denunciar}, \esp{la denuncia}, \esp{proteger}, \esp{la protección}, \esp{el refugio}, \esp{el centro de acogida}, \esp{la orden de alejamiento}, \esp{la igualdad de derechos}, \esp{la igualdad de oportunidades}, \esp{la educación afectivo-sexual}, \esp{el machismo}, \esp{el feminicidio}, \esp{el acoso}, \esp{el consentimiento}, \esp{la sensibilización}, \esp{la prevención}, \esp{la impunidad}, \esp{la justicia}.
\end{definition}

\courssub{Méthode 1 : Commenter un texte (dialogue ou article de presse)}

\begin{methode}[Étapes du commentaire composé]
\begin{enumerate}
\item \cle{Introduction} : Présenter le document (nature, source, date, thème), annoncer le problème (problématique) et le plan (deux ou trois parties).
\item \cle{Développement} (2--3 parties) : Chaque partie = une idée force. \textbf{Citer} le texte (guillemets espagnols \esp{« ... »}), \textbf{analyser} (vocabulaire, grammaire, figures de style, ton), \textbf{interpréter} (portée, intention, contexte social).
\item \cle{Conclusion} : Répondre à la problématique, ouvrir sur un autre texte, l'actualité (Cameroun, Espagne, Guinée équatoriale) ou une réflexion personnelle.
\end{enumerate}
\textbf{Connecteurs logiques :} \esp{En primer lugar}, \esp{Además}, \esp{Por otro lado}, \esp{En consecuencia}, \esp{Por tanto}, \esp{Finalmente}, \esp{En conclusión}.
\end{methode}

\begin{exoresolu}[Commentaire modèle : le dialogue familial]
\textbf{Sujet :} Commentez le dialogue ci-dessus en montrant comment la mère gère le conflit par l'écoute active.

\textbf{Production modèle :}

\esp{El texto presentado es un diálogo familiar entre una madre y su hija en Yaundé. El tema central es el conflicto generacional: la hija quiere abandonar los estudios para abrir una peluquería, mientras que la madre valora el éxito escolar. La problemática es: ¿cómo resolver un desacuerdo familiar sin romper el vínculo afectivo?}

\esp{En primer lugar, la madre utiliza marcadores de escucha activa. Ante el anuncio sorpresivo, reacciona con « ¿De verdad? », mostrando su asombro sin agresividad. Luego reformula los sentimientos de la hija: « O sea que te sientes perdida... », lo que valida la emoción de la joven. Esta técnica, llamada escucha activa, evita la escalada del conflicto.}

\esp{Además, la madre reconoce su propia responsabilidad: « Tienes razón. Siempre estás ocupada en el mercado. » Esta autocrítica crea un clima de confianza. Finalmente, propone un espacio de diálogo diferido: « Esta noche hablamos con calma », aplazando la decisión para evitar la precipitación.}

\esp{En conclusión, el diálogo ilustra una resolución pacífica basada en la empatía, la reformulación y el compromiso. Este modelo es transferible a contextos escolares o comunitarios en Camerún, donde el diálogo intergeneracional es clave para la cohesión social.}
\end{exoresolu}

\courssub{Méthode 2 : Comprendre et mener un dialogue familial (écoute active, relance, reformulation)}

\begin{methode}[Réussir l'écoute active et la relance dans un dialogue]
\begin{enumerate}
\item \cle{Identifier la situation} : qui parle ? à qui ? où ? de quoi ? (famille, école, conflit, projet). Repérer le registre : \esp{tú} (familier) ou \esp{usted} (formel), ton affectueux ou tendu.
\item \cle{Repérer les marqueurs d'écoute active} : \esp{ya}, \esp{claro}, \esp{entiendo}, \esp{¿de verdad?}, \esp{¡no me digas!}, \esp{ajá}, \esp{vaya}. Ils montrent que l'interlocuteur suit la conversation.
\item \cle{Relancer} avec des questions ouvertes : \esp{¿Y tú qué opinas?}, \esp{¿Por qué dices eso?}, \esp{¿Qué pasó después?}, \esp{¿Me puedes explicar mejor?}, \esp{¿Cómo te sientes al respecto?}
\item \cle{Reformuler} pour vérifier la compréhension : \esp{O sea que...}, \esp{Entonces, lo que quieres decir es que...}, \esp{Si te entiendo bien, piensas que...}, \esp{En resumen...}
\item \cle{Conclure} par un accord, un compromis ou une proposition : \esp{De acuerdo}, \esp{Vale, lo hablamos esta noche}, \esp{Propongo que...}, \esp{Busquemos una solución juntos.}
\end{enumerate}
\textbf{Structures à retenir} : \esp{Me parece que...} / \esp{No estoy de acuerdo porque...} / \esp{Escúchame un momento} / \esp{No te enfades, por favor.} / \esp{Te entiendo perfectamente.}
\end{methode}

\begin{experience}[Jeu de rôle : conflit familial]
\textbf{Consigne :} Par binômes. L'un joue le parent, l'autre l'adolescent(e). Situation : l'adolescent(e) veut quitter l'école pour travailler / se marier / faire du sport de haut niveau. Le parent est inquiet.
\textbf{Étapes :} 1. L'adolescent annonce sa décision. 2. Le parent utilise \textbf{au moins 3 marqueurs d'écoute active}, \textbf{1 reformulation}, \textbf{1 relance}. 3. Ils trouvent un compromis. 4. Présentation orale devant la classe. Évaluation par les pairs (grille : écoute, reformulation, politesse, accord).
\end{experience}

\courssub{Méthode 3 : Participer à un débat sur les violences faites aux femmes}

\begin{methode}[Exprimer son opinion et argumenter dans un débat guidé]
\begin{enumerate}
\item \cle{Prendre position clairement} : \esp{En mi opinión...}, \esp{Estoy convencido/a de que...}, \esp{A mi juicio...}, \esp{Desde mi punto de vista...}, \esp{Considero que...}
\item \cle{Structurer l'argumentation} : \esp{En primer lugar}, \esp{Además}, \esp{Por otra parte}, \esp{Por último} ; illustrer avec \esp{Por ejemplo}, \esp{Un caso claro es...}
\item \cle{Nuancer ou concéder} : \esp{Es cierto que..., pero...}, \esp{Aunque reconozco que..., pienso que...}, \esp{No niego que..., sin embargo...}
\item \cle{Réagir à autrui} : \esp{Estoy de acuerdo con... porque...} / \esp{No comparto tu opinión, ya que...} / \esp{Permíteme que no esté de acuerdo.} / \esp{¿Me permites una objeción?}
\item \cle{Employer le vocabulaire précis du thème} (voir Lexique ci-dessus) : \esp{la violencia de género}, \esp{el maltrato}, \esp{la víctima}, \esp{el agresor}, \esp{denunciar}, \esp{proteger}, \esp{la igualdad de derechos}, \esp{la educación}, \esp{el machismo}, \esp{el feminicidio}.
\item \cle{Conclure par une proposition concrète} : \esp{Por eso, propongo que se creen centros de acogida...} (subjonctif après \esp{propongo que}), \esp{Exijo que el gobierno aplique la ley...} (subjonctif), \esp{Es necesario que la sociedad reaccione...} (subjonctif).
\end{enumerate}
\end{methode}

\begin{experience}[Débat guidé : La escuela y la violencia de género]
\textbf{Sujet :} \esp{¿La escuela debe participar activamente en la lucha contra la violencia de género?}
\textbf{Organisation :} 4 groupes (Pour / Contre / Nuancé / Jury). Préparation 15 min (arguments, exemples, vocabulaire). Débat 20 min (modéré par le prof). Jury évalue : justesse linguistique (ser/estar, subjonctif), richesse lexicale, cohérence, respect du tour de parole.
\end{experience}

\begin{exoresolu}[Débat : rédiger une prise de position écrite]
\textbf{Énoncé.} Rédigez une prise de position de 8 à 10 lignes en espagnol : opinion, deux arguments, un exemple contextualisé (Cameroun ou Espagne), une concession, une conclusion avec proposition (subjonctif).
\tcblower
\textbf{Production modèle :}

\esp{En mi opinión, la escuela debe ser un pilar fundamental en la lucha contra la violencia de género. En primer lugar, la educación en igualdad desde la infancia previene la interiorización de roles machistas: los niños aprenden que el respeto no tiene género. Además, el profesorado está en posición privilegiada para detectar signos de maltrato en el alumnado o en sus familias y activar los protocolos de protección. Por ejemplo, en Camerún, el MINEDUB ha integrado módulos de género en los nuevos programas de secundaria. Es cierto que la familia es el primer agente socializador, pero no todas las familias tienen las herramientas para educar en igualdad. Por eso, propongo que se cree una asignatura obligatoria de "Educación para la Ciudadanía e Igualdad" en todos los ciclos y que se forme al profesorado en detección precoz.}

\textbf{Justifications :} \esp{debe ser} (obligation morale) ; \esp{previene} (indicatif, certitude) ; \esp{está en posición} (localisation/état) ; \esp{ha integrado} (passé composé) ; \esp{Es cierto que... pero...} (concession) ; \esp{propongo que se cree / se forme} (subjonctif présent \esp{cree}, \esp{forme} après verbe de volonté).
\end{exoresolu}

\courssub{Grammaire : SER y ESTAR — Emplois fondamentaux}

\begin{propriete}[L'arbre de décision SER / ESTAR (Présent de l'indicatif)]
\cle{Règle d'or :} \textbf{SER} = Identité, essence, classification, origine, heure, date, possession, matière, événement (lieu/date). \textbf{ESTAR} = Localisation (personnes, choses), état temporaire, résultat d'action, action en cours (gérondif).
\begin{center}
\begin{tabularx}{\linewidth}{|l|X|X|}\hline
\rowcolor{popblueL} \textbf{Critère} & \textbf{SER} & \textbf{ESTAR} \\ \hline
\cle{Localisation} & \textbf{Événement} : \esp{La boda es en la iglesia.} & \textbf{Personne/Chose} : \esp{La iglesia está en el centro.} \\ \hline
\cle{Temps} & \textbf{Heure/Date} : \esp{Son las tres. Hoy es lunes.} & \textbf{Durée/Moment} : \esp{Estoy aquí desde las tres.} \\ \hline
\cle{Identité} & \textbf{Nationalité, profession, religion} : \esp{Soy camerunés, soy médico, soy católico.} & --- \\ \hline
\cle{Caractère} & \textbf{Essentiel, permanent} : \esp{El cacao es dulce. Ella es alta.} & \textbf{Accidentel, passager} : \esp{El cacao está dulce (aujourd'hui). Ella está pálida.} \\ \hline
\cle{État/Santé} & \textbf{Caractère} : \esp{Es un hombre enfermo (malade chronique).} & \textbf{État momentané} : \esp{Está enfermo (grippe).} \\ \hline
\cle{Action} & \textbf{Passif (être fait par)} : \esp{El libro es escrito por él.} & \textbf{En cours (gérondif)} : \esp{Está escribiendo el libro.} \\ \hline
\cle{Matière} & \esp{La mesa es de madera.} & --- \\ \hline
\cle{Possession} & \esp{El coche es de mi padre.} & --- \\ \hline
\end{tabularx}
\end{center}
\textbf{Conjugaison présent (rappel) :}
\begin{center}
\begin{tabular}{|l|l|l|}\hline
\rowcolor{popblueL} \textbf{Personne} & \textbf{SER} & \textbf{ESTAR} \\ \hline
yo & soy & estoy \\ \hline
tú & eres & estás \\ \hline
él/ella/usted & es & está \\ \hline
nosotros/as & somos & estamos \\ \hline
vosotros/as & sois & estáis \\ \hline
ellos/ellas/ustedes & son & están \\ \hline
\end{tabular}
\end{center}
\cle{Attention accents :} \esp{estás, está, estamos, estáis, están} (accent sur la dernière syllabe sauf \emph{nosotros/as}). \esp{ésta} (démonstratif) n'existe plus (norme RAE 2010 : \esp{esta} sans accent), mais \esp{está} (verbe) garde l'accent.
\end{propriete}

\begin{exoresolu}[QCM justifié : SER ou ESTAR au présent]
\textbf{Énoncé.} Complétez avec la forme correcte, puis justifiez en français (critère de l'arbre de décision).
\begin{enumerate}
\item Mi padre \_\_\_ médico en el hospital de Douala.
\item Los niños \_\_\_ en el patio, jugando al fútbol.
\item Hoy \_\_\_ lunes y \_\_\_ las ocho de la mañana.
\item La sopa \_\_\_ muy caliente; espera un poco.
\item El partido de fútbol \_\_\_ en el estadio Ahmadou Ahidjo.
\item \_\_\_ tú de acuerdo con la nueva ley?
\item Mi hermana \_\_\_ muy lista; siempre saca buenas notas.
\item \_\_\_ listos para salir? El bendskin \_\_\_ llegando.
\end{enumerate}
\tcblower
\textbf{Solutions et justifications :}
\begin{enumerate}
\item \esp{es} --- Profession $\rightarrow$ \textbf{SER}.
\item \esp{están} --- Localisation personnes $\rightarrow$ \textbf{ESTAR}.
\item \esp{es ; son} --- Date (\esp{es}) + Heure pluriel (\esp{son las ocho}) $\rightarrow$ \textbf{SER}. Rappel : \esp{es la una}, \esp{son las dos}.
\item \esp{está} --- État temporaire (température) $\rightarrow$ \textbf{ESTAR}.
\item \esp{es} --- Lieu d'un \textbf{événement} $\rightarrow$ \textbf{SER}. (Contraste : \esp{El estadio está en Yaoundé} = localisation du bâtiment).
\item \esp{¿Estás} --- Expression figée \esp{estar de acuerdo} $\rightarrow$ \textbf{ESTAR}.
\item \esp{es} --- Qualité intellectuelle permanente (intelligent) $\rightarrow$ \textbf{SER} + \esp{listo}.
\item \esp{¿Están} (prêts) / \esp{está} (action en cours : \esp{está llegando}) $\rightarrow$ \textbf{ESTAR} pour les deux.
\end{enumerate}
\end{exoresolu}

\courssub{SER / ESTAR + Adjectif : Changements de sens et expressions figées}

\begin{propriete}[Adjectifs à double sens : la règle sémantique]
Avec \textbf{SER}, l'adjectif qualifie l'\emph{essence} (caractère permanent, nature). Avec \textbf{ESTAR}, il qualifie l'\emph{état} (situation passagère, résultat, impression subjective).
\begin{center}
\begin{tabularx}{\linewidth}{|l|X|X|}\hline
\rowcolor{popblueL} \textbf{Adjectif} & \textbf{SER + adj. (Essence)} & \textbf{ESTAR + adj. (État/Résultat)} \\ \hline
\esp{aburrido} & ennuyeux (personne/film) & ennuyé, qui s'ennuie \\ \hline
\esp{bueno / malo} & bon/mauvais (moral, qualité, compétence) & bon au goût / délicieux ; malade / gâté (nourriture) \\ \hline
\esp{listo} & intelligent, vif d'esprit & prêt, préparé \\ \hline
\esp{rico} & riche (argent) & délicieux, savoureux (nourriture) \\ \hline
\esp{verde} & vert (couleur) & pas mûr (fruit) ; novice (personne) \\ \hline
\esp{vivo} & vif, éveillé, malin & vivant, en vie \\ \hline
\esp{seguro} & sûr, fiable (personne/objet) & certain, convaincu (que) \\ \hline
\esp{orgulloso} & fier (caractère, orgueilleux) & fier de (sentiment passager) $\rightarrow$ \esp{estar orgulloso de} \\ \hline
\esp{interesado} & intéressé (égoïste) & intéressé par $\rightarrow$ \esp{estar interesado en} \\ \hline
\esp{pálido} & pâle (teint naturel) & pâle (émotion, maladie) \\ \hline
\end{tabularx}
\end{center}
\cle{Expressions figées avec ESTAR (obligatoires) :}
\begin{itemize}
\item \esp{estar de acuerdo} (être d'accord) --- \textbf{Jamais \esp{ser de acuerdo}.}
\item \esp{estar de vacaciones / de viaje / de baja} (être en vacances / voyage / arrêt maladie)
\item \esp{estar de moda} (être à la mode)
\item \esp{estar enfermo} (être malade)
\item \esp{estar vivo / muerto} (être vivant / mort)
\item \esp{estar a punto de + infinitif} (être sur le point de)
\item \esp{estar para + infinitif} (être prêt à / devoir être fait : \esp{la cena está para servir})
\item \esp{estar por + infinitif} (avoir envie de / rester à faire : \esp{estoy por irme})
\end{itemize}
\end{propriete}

\begin{exoresolu}[Choisir le sens correct : traduction thématique]
\textbf{Énoncé.} Traduisez en espagnol en justifiant le choix SER/ESTAR.
\begin{enumerate}
\item Mon frère est très ennuyeux : il parle toujours de politique.
\item Les enfants sont ennuyés parce qu'il pleut.
\item Ce plat de ndolé est délicieux.
\item Ma sœur est très intelligente : elle est première de sa classe.
\item Es-tu prête ? Le bendskin nous attend.
\item Malgré l'accident, le chauffeur est vivant.
\item Ce jeune est très intéressé par la politique (il veut en faire son métier).
\item Ce vendeur est intéressé : il veut nous arnaquer.
\item La porte est ouverte (action terminée).
\item La porte est ouverte (caractéristique : elle n'a pas de serrure).
\end{enumerate}
\tcblower
\textbf{Solutions :}
\begin{enumerate}
\item \esp{Mi hermano es muy aburrido: siempre habla de política.} (Qualité permanente $\rightarrow$ SER).
\item \esp{Los niños están aburridos porque llueve.} (État passager $\rightarrow$ ESTAR).
\item \esp{Este plato de ndolé está riquísimo / está muy bueno.} (Goût $\rightarrow$ ESTAR + \esp{rico/bueno}). \esp{Es rico} = « Il est riche ».
\item \esp{Mi hermana es muy lista: es la primera de su clase.} (Intelligence $\rightarrow$ SER + \esp{listo}).
\item \esp{¿Estás lista? El bendskin nos espera.} (« Prête » $\rightarrow$ ESTAR + \esp{lista}, accord féminin).
\item \esp{A pesar del accidente, el conductor está vivo.} (« En vie » $\rightarrow$ ESTAR + \esp{vivo}). \esp{Es vivo} = « Il est vif/malin ».
\item \esp{Este joven está muy interesado en la política.} (Intérêt pour $\rightarrow$ ESTAR + \esp{interesado en}).
\item \esp{Este vendedor es muy interesado.} (Égoïste $\rightarrow$ SER + \esp{interesado}).
\item \esp{La puerta está abierta.} (Résultat d'action $\rightarrow$ ESTAR + participe \esp{abierta}).
\item \esp{La puerta es abierta.} (Caractéristique structurelle $\rightarrow$ SER + adjectif \esp{abierta}).
\end{enumerate}
\end{exoresolu}

\courssub{Méthode 4 : Traduire « être » en espagnol (Thème et Version)}

\begin{methode}[Procédure rigoureuse pour le thème (FR $\rightarrow$ ES)]
\begin{enumerate}
\item \cle{Segmenter} la phrase française : isoler chaque verbe « être » (ou locution verbale équivalente).
\item \cle{Analyser} chaque occurrence avec l'arbre de décision (Méthode 3) : Essence ? $\rightarrow$ SER. État/Lieu/Action en cours ? $\rightarrow$ ESTAR.
\item \cle{Traiter les cas spéciaux} :
\begin{itemize}
\item « Être en train de + infinitif » $\rightarrow$ \esp{ESTAR + gérondif} (\esp{está leyendo}).
\item « Il est + adjectif + que » (impersonnel) $\rightarrow$ \esp{ES + adj + QUE + SUBJONCTIF} (\esp{Es importante que estudies}).
\item « Il est + heure/date » $\rightarrow$ \esp{SER} (\esp{Son las cinco}, \esp{Es lunes}).
\item « Être d'accord / en vacances / malade / à la mode » $\rightarrow$ \esp{ESTAR} (expressions figées).
\item Passif : « Il a été construit » $\rightarrow$ \esp{SER + participe} (\esp{Fue construido}) ou \esp{SE + verbe} (\esp{Se construyó}). « Il est construit » (état) $\rightarrow$ \esp{ESTAR + participe} (\esp{Está construido}).
\end{itemize}
\item \cle{Accorder} le participe/adjectif en genre et nombre avec le sujet espagnol.
\item \cle{Vérifier} : accents (\esp{está, están, médico, inglés}), ponctuation (\esp{¿ ? ¡ !}), contraction (\esp{del, al}), subjonctif si déclencheur.
\end{enumerate}
\end{methode}

\begin{methode}[Procédure pour la version (ES $\rightarrow$ FR)]
\begin{enumerate}
\item Identifier le verbe : \esp{ser} ou \esp{estar}.
\item \esp{Ser} $\rightarrow$ généralement « être » (identité, définition, heure, passif d'action).
\item \esp{Estar} $\rightarrow$ « se trouver » (lieu), « être » (état, santé, humeur), « être en train de » (gérondif), expressions figées.
\item \esp{Ser/Estar + participe} : \esp{Ser} = passif d'action (par qui?); \esp{Estar} = état résultant (pas d'agent).
\item Attention aux adjectifs à double sens : traduire le sens, pas le mot.
\end{enumerate}
\end{methode}

\begin{exoresolu}[Thème guidé : phrases complexes avec « être »]
\textbf{Énoncé.} Traduisez en espagnol.
\begin{enumerate}
\item Ma grand-mère est malade : elle est à l'hôpital depuis lundi.
\item Nous sommes camerounais et nous sommes fiers de notre pays.
\item Il est important que les femmes soient protégées contre la violence.
\item La maison de mes parents est près du marché de Mokolo.
\item Quelle heure est-il ? Il est sept heures et demie.
\item Le match a lieu demain au stade. Le stade est à Yaoundé.
\item Elles sont en train de préparer le dîner.
\item Ce gâteau est fait par ma tante (action). Ce gâteau est fait (terminé).
\end{enumerate}
\tcblower
\textbf{Solutions détaillées :}
\begin{enumerate}
\item \esp{Mi abuela está enferma: está en el hospital desde el lunes.} (Santé + Localisation $\rightarrow$ ESTAR). Accent sur \esp{está}.
\item \esp{Somos cameruneses y estamos orgullosos de nuestro país.} (Nationalité $\rightarrow$ SER ; Sentiment $\rightarrow$ ESTAR + \esp{orgullosos}, accord masc. pl.).
\item \esp{Es importante que las mujeres estén protegidas contra la violencia.} (Impersonnel $\rightarrow$ SER ; Subjonctif \esp{estén} après \esp{es importante que} ; \esp{protegidas} accord fém. pl.).
\item \esp{La casa de mis padres está cerca del mercado de Mokolo.} (Localisation chose $\rightarrow$ ESTAR ; contraction \esp{del} = \esp{de + el}).
\item \esp{¿Qué hora es? Son las siete y media.} (Heure $\rightarrow$ SER ; pluriel \esp{son} ; \esp{¿ ?}).
\item \esp{El partido es mañana en el estadio. El estadio está en Yaundé.} (Événement $\rightarrow$ SER ; Localisation bâtiment $\rightarrow$ ESTAR).
\item \esp{Están preparando la cena.} (Action en cours $\rightarrow$ ESTAR + gérondif).
\item \esp{Este pastel es hecho por mi tía.} (Passif action $\rightarrow$ SER + participe \esp{hecho}). \esp{Este pastel está hecho.} (État/résultat $\rightarrow$ ESTAR + participe \esp{hecho}).
\end{enumerate}
\end{exoresolu}

\begin{attention}[Les 5 erreurs fatales au Bac (SER/ESTAR + Traduction)]
\begin{enumerate}
\item \cle{Accent manquant sur ESTAR conjugué} : \esp{está, estás, están} portent l'accent tonique. \esp{Esta} (sans accent) = adjectif démonstratif « cette ». \esp{Mi madre esta cansada} = \textbf{Faute majeure}.
\item \cle{Calque sur « être d'accord / en train de »} : \textbf{Jamais} \esp{ser de acuerdo}, \textbf{jamais} \esp{ser + gérondif}. Toujours \esp{estar de acuerdo}, \esp{estar + gérondif}.
\item \cle{Confusion adjectifs double sens} : \esp{Juan es listo} (intelligent) $\neq$ \esp{Juan está listo} (prêt). \esp{La sopa es buena} (qualité nutritive) $\neq$ \esp{La sopa está buena} (bon goût). \esp{María es viva} (maline) $\neq$ \esp{María está viva} (vivante).
\item \cle{Oubli du subjonctif} après : \esp{Es importante/necesario/urgente que...}, \esp{Propongo/exijo/quiero que...}, \esp{No creo que...}, \esp{Dudo que...}. Ex: \esp{Propongo que se proteja} (et non \esp{protege}).
\item \cle{Faute d'accord du participe/adjectif} : Après SER/ESTAR, l'attribut s'accorde avec le sujet. \esp{Las víctimas están protegidas} (fém. pl.), \esp{Los libros son interesantes} (masc. pl.). \esp{La puerta está abierta} (fém. sg.).
\end{enumerate}
\end{attention}

\courssub{Entraînement : Exercices de synthèse}

\begin{exercice}[Exercices de révision — Leçon E02]
\textbf{Exercice 1 (Dialogue) :} Complétez le dialogue entre un père et son fils à Douala avec une formule d'écoute active, de relance ou de reformulation.
\esp{--- Papá, he decidido no presentarme al examen de conducir.}\\
\esp{--- ¿(1)? Pero si llevas meses practicando...}\\
\esp{--- Es que tengo mucho miedo de tener un accidente.}\\
\esp{--- (2), tienes miedo. ¿Es por el tráfico de aquí?}\\
\esp{--- Sí, los bendskins no respetan nada.}\\
\esp{--- (3)... Es normal tener miedo. Propongo que tomes clases con un profesor particular antes de decidir.}

\textbf{Exercice 2 (Vocabulaire/Débat) :} Associez les termes espagnols à leur définition française.
\begin{tabular}{ll}
1. \esp{El agresor} & a. Personne qui subit la violence \\
2. \esp{La víctima} & b. Personne qui exerce la violence \\
3. \esp{Denunciar} & c. Signaler aux autorités \\
4. \esp{La orden de alejamiento} & d. Mesure judiciaire d'interdiction d'approcher \\
5. \esp{El machismo} & e. Idéologie de supériorité masculine
\end{tabular}

\textbf{Exercice 3 (SER/ESTAR) :} Conjuguez le verbe entre parenthèses au présent et justifiez en une phrase.
1. Los libros \_\_\_ (ser/estar) en la mesa. (Localisation chose)
2. Mi primo \_\_\_ (ser/estar) ingeniero. (Profession)
3. Hoy \_\_\_ (ser/estar) muy cansado. (État passager)
4. \_\_\_ (ser/estar) ustedes de acuerdo? (Expression figée)
5. La conferencia \_\_\_ (ser/estar) en el aula magna. (Événement)
6. El mango \_\_\_ (ser/estar) verde. (Pas mûr / Couleur ? Choisir « pas mûr »)

\textbf{Exercice 4 (Thème) :} Traduisez en espagnol.
1. Il est tard, nous sommes pressés.
2. Mon oncle est médecin, il est très intelligent.
3. Les filles sont en vacances à Kribi.
4. Il est interdit de fumer ici. (Indice : \esp{Está prohibido fumar...})
5. Je propose que l'on crée une association pour les victimes.

\textbf{Exercice 5 (Version) :} Traduisez en français.
1. \esp{Mi abuelo está muy viejo, pero su mente está muy clara.}
2. \esp{Es fundamental que los hombres cambien su mentalidad.}
3. \esp{El feminicidio es un problema grave en nuestra sociedad.}
4. \esp{¿Estás de acuerdo en que la educación es la clave?}
5. \esp{El partido fue jugado ayer, pero el estadio está vacío hoy.}
\end{exercice}

\begin{corrige}[Exercices de révision — Leçon E02]
\textbf{Exercice 1 :}
(1) \esp{¿De verdad?} / \esp{¡No me digas!} (Étonnement/Écoute active).
(2) \esp{O sea que...} / \esp{Entonces, lo que dices es que...} (Reformulation).
(3) \esp{Entiendo / Te comprendo / Si te entiendo bien...} (Empathie avant proposition).

\textbf{Exercice 2 :} 1-b, 2-a, 3-c, 4-d, 5-e.

\textbf{Exercice 3 :}
1. \esp{están} (Localisation chose $\rightarrow$ ESTAR).
2. \esp{es} (Profession $\rightarrow$ SER).
3. \esp{está} (État passager $\rightarrow$ ESTAR).
4. \esp{¿Están} (Expression figée \esp{estar de acuerdo} $\rightarrow$ ESTAR).
5. \esp{es} (Événement $\rightarrow$ SER).
6. \esp{está} (État « pas mûr » $\rightarrow$ ESTAR + \esp{verde}).

\textbf{Exercice 4 :}
1. \esp{Es tarde, tenemos prisa.} (\esp{Ser} heure ; \esp{Tener prisa} expression avoir hâte/être pressé).
2. \esp{Mi tío es médico, es muy listo.} (Profession + Qualité $\rightarrow$ SER).
3. \esp{Las chicas están de vacaciones en Kribi.} (État/Expression \esp{estar de vacaciones} $\rightarrow$ ESTAR).
4. \esp{Está prohibido fumar aquí.} (Impersonnel passif état $\rightarrow$ ESTAR + participe).
5. \esp{Propongo que se cree una asociación para las víctimas.} (Volonté $\rightarrow$ Subjonctif \esp{se cree}).

\textbf{Exercice 5 :}
1. Mon grand-père est très vieux (âge/état), mais son esprit est très clair (état).
2. Il est fondamental que les hommes changent de mentalité (Subjonctif \esp{cambien}).
3. Le féminicide est un problème grave dans notre société.
4. Es-tu d'accord pour dire que l'éducation est la clé ?
5. Le match a eu lieu hier (passif action \esp{fue jugado}), mais le stade est vide aujourd'hui (état \esp{está vacío}).
\end{corrige}

\begin{savaistu}[Culture : La Guinée équatoriale et la lutte contre les violences]
La Guinée équatoriale (\esp{Guinea Ecuatorial}), seul pays africain hispanophone, a ratifié la Convention de Belém do Pará (1994) et la Convention CEDAW. Depuis 2020, un \esp{Protocolo de Actuación} coordonne police, justice et santé pour protéger les victimes. À Malabo et Bata, des \esp{Oficinas de Atención a la Víctima} offrent aide juridique et psychologique. Cependant, le machisme culturel et l'impunité persistent. Le Cameroun, pays voisin francophone/anglophone, partage ces défis : la loi de 2016 sur le Code pénal réprime les violences conjugales, mais l'application reste inégale. L'espagnol, langue officielle en Guinée, devient un outil régional de plaidoyer pour les droits des femmes en Afrique centrale.
\end{savaistu}

\newpage

\coursec{Exercices}

\courssub{Je vérifie mes connaissances}

\begin{exercice}[QCM : ser ou estar ?]
Pour chaque phrase, choisis la bonne forme du verbe et entoure-la.

\begin{enumerate}
\item Mi hermano \textbf{es / está} médico en el Hospital Central de Yaoundé.
\item Hoy \textbf{es / está} lunes 12 de mayo.
\item La sopa \textbf{es / está} fría : hay que calentarla.
\item ¿Dónde \textbf{es / está} el mercado de Mokolo, por favor?
\item Mi padre \textbf{es / está} muy cansado después del trabajo.
\item Ese chico \textbf{es / está} muy listo : siempre encuentra la solución.
\item ¡Despierta! Ya \textbf{somos / estamos} listos para salir.
\item Nosotros \textbf{somos / estamos} de acuerdo con tu propuesta.
\item La fiesta de mi prima \textbf{es / está} en Douala, el sábado próximo.
\item Los derechos de las mujeres \textbf{son / están} un tema importante del debate.
\end{enumerate}
\end{exercice}

\begin{exercice}[Vrai ou faux ?]
Indique si chaque affirmation est \textbf{vraie (V)} ou \textbf{fausse (F)} et justifie ta réponse en une phrase.

\begin{enumerate}
\item On utilise \esp{ser} pour exprimer la profession d'une personne.
\item \esp{Estar} sert à exprimer l'identité et la nationalité.
\item \esp{Ser aburrido} signifie « être ennuyé ».
\item Pour dire l'heure, on utilise toujours \esp{ser}.
\item \esp{Estar de acuerdo} est une expression figée qui signifie « être d'accord ».
\item On dit \esp{estar enfermo} pour parler d'une maladie passagère.
\end{enumerate}
\end{exercice}

\begin{exercice}[Vocabulaire : la famille et la violence de genre]
Complète chaque phrase avec le mot espagnol qui convient, choisi dans la liste suivante :

\begin{center}
\esp{maltratador -- denunciar -- el hogar -- la nuera -- los derechos -- el maltrato}
\end{center}

\begin{enumerate}
\item Mi cuñada y mi \dots\ vienen a cenar esta noche a \dots\ de mis suegros.
\item Es importante \dots\ los casos de violencia a la policía.
\item El \dots\ de las mujeres es un problema grave en muchas sociedades.
\item Un hombre que pega a su esposa es un \dots.
\item Hay que defender \dots\ de las mujeres y de los niños.
\end{enumerate}
\end{exercice}

\begin{exercice}[Conjugaison à compléter]
Complète le tableau avec les formes correctes de \esp{ser} et \esp{estar} au présent de l'indicatif.

\begin{center}
\begin{tabular}{|l|l|l|}
\hline
\rowcolor{popblueL} \textbf{Personne} & \textbf{ser} & \textbf{estar} \\
\hline
yo & \dots & \dots \\
\hline
tú & \dots & \dots \\
\hline
él / ella / usted & \dots & \dots \\
\hline
nosotros / nosotras & \dots & \dots \\
\hline
vosotros / vosotras & \dots & \dots \\
\hline
ellos / ellas / ustedes & \dots & \dots \\
\hline
\end{tabular}
\end{center}
\end{exercice}

\courssub{Je m'entraîne}

\begin{exercice}[Thème : traduis en espagnol]
Traduis les phrases suivantes en espagnol et justifie ton choix entre \esp{ser} et \esp{estar}.

\begin{enumerate}
\item Ma sœur est institutrice à Bafoussam.
\item Il est huit heures du matin.
\item Les enfants sont fatigués après le match de football.
\item Le marché est très loin de notre quartier.
\item Mon grand-père est malade depuis lundi.
\end{enumerate}
\end{exercice}

\begin{exercice}[Transformations : changement de sens]
Transforme chaque phrase en remplaçant \esp{ser} par \esp{estar} (ou l'inverse) devant l'adjectif, puis explique en français le changement de sens obtenu.

\begin{enumerate}
\item \esp{Juan es muy listo.}
\item \esp{La sopa está fría.}
\item \esp{Mi prima es aburrida.}
\item \esp{Estamos listos para el examen.}
\item \esp{Pedro es muy guapo.}
\end{enumerate}
\end{exercice}

\begin{exercice}[Texte à trous]
Complète le dialogue familial suivant avec la forme correcte de \esp{ser} ou \esp{estar} au présent.

\begin{textebox}[Un dialogue à la maison]
--- Mamá, ¿dónde \dots\ mis zapatos? ¡No los encuentro!\\
--- \dots\ en la cocina, hijo. ¿Por qué \dots\ tan nervioso esta mañana?\\
--- Porque hoy \dots\ el día del debate en el liceo. El tema \dots\ la violencia contra las mujeres.\\
--- \dots\ un tema muy serio. ¿Tú qué opinas?\\
--- Yo \dots\ de acuerdo con castigar a los maltratadores. La violencia \dots\ inaceptable.\\
--- Muy bien, hijo. Tu padre y yo \dots\ orgullosos de ti.
\end{textebox}
\end{exercice}

\begin{exercice}[Appariements : expressions figées]
Relie chaque expression espagnole (colonne de gauche) à sa traduction française (colonne de droite).

\begin{center}
\begin{tabularx}{\linewidth}{|X|X|}
\hline
\rowcolor{popblueL} \textbf{Español} & \textbf{Français} \\
\hline
1. \esp{estar de acuerdo} & a. être en train de \\
\hline
2. \esp{estar de vacaciones} & b. être d'accord \\
\hline
3. \esp{estar de pie} & c. être de retour \\
\hline
4. \esp{estar de vuelta} & d. être en vacances \\
\hline
5. \esp{estar + gérondif} & e. être debout \\
\hline
\end{tabularx}
\end{center}
\end{exercice}

\courssub{Je me prépare au Bac}

\begin{exercice}[Compréhension et questions de langue]
Lis attentivement le texte suivant, puis réponds aux questions.

\begin{textebox}[Un debate en el barrio]
En un barrio de Douala, la asociación « Mujeres Unidas » organizó ayer un debate sobre la violencia de género. María, una joven activista, abrió la sesión con estas palabras : « La violencia contra las mujeres no es un problema privado : es un problema de toda la sociedad. Muchas víctimas están calladas porque tienen miedo o porque dependen económicamente de sus maltratadores. »

Un vecino, don Samuel, respondió : « Estoy de acuerdo, pero también hay que educar a los niños desde pequeños. Si un niño ve que su padre es violento con su madre, mañana él será violento también. »

Al final del debate, todos los participantes estaban convencidos de una cosa : denunciar es un deber, y el silencio es cómplice. La próxima reunión será el mes que viene, en la escuela del barrio.
\end{textebox}

\textbf{Questions de compréhension} (réponds en espagnol, avec des phrases complètes) :

\begin{enumerate}
\item ¿Quién organiza el debate y dónde tiene lugar?
\item Según María, ¿por qué muchas víctimas están calladas?
\item ¿Qué propone don Samuel para luchar contra la violencia?
\item ¿Cuándo será la próxima reunión?
\end{enumerate}

\textbf{Questions de langue} :

\begin{enumerate}
\item Releve dans le texte trois emplois de \esp{estar} et explique chacun (état, localisation, expression figée).
\item Justifie l'emploi de \esp{ser} dans : \esp{« La violencia contra las mujeres no es un problema privado »}.
\item Traduis en français la dernière phrase du texte.
\end{enumerate}
\end{exercice}

\begin{exercice}[Thème et version type Bac]
\textbf{Thème.} Traduis en espagnol les cinq phrases suivantes en justifiant à chaque fois ton choix entre \esp{ser} et \esp{estar} (identité, lieu, état, heure, action en cours).

\begin{enumerate}
\item Nous sommes camerounais et nous sommes fiers de notre pays.
\item Il est neuf heures et demie ; la réunion est dans la salle des professeurs.
\item Ta mère est en train de préparer le repas dans la cuisine.
\item Ces femmes sont courageuses : elles dénoncent les violences.
\item La maison de mes grands-parents est près de l'église.
\end{enumerate}

\textbf{Version.} Traduis en français le paragraphe suivant, extrait d'un article de presse espagnol.

\begin{textebox}[Extracto de prensa]
La violencia de género es una realidad que todavía está presente en muchos países. Las víctimas suelen estar solas y asustadas, pero cada vez más mujeres son valientes y denuncian a sus agresores. Las asociaciones de ayuda están abiertas día y noche, y su trabajo es fundamental. Estar informado es el primer paso : la sociedad entera debe estar atenta y comprometida.
\end{textebox}
\end{exercice}

\begin{exercice}[Expression : débat et rédaction type Bac]
\textbf{Partie A — Expression orale (en binôme, 5 minutes de préparation).} Avec ton camarade, organise un débat guidé sur la question suivante :

\begin{center}
\esp{¿Estás a favor de castigar más duramente a los maltratadores?}
\end{center}

Consignes obligatoires :
\begin{itemize}
\item chaque élève exprime au moins deux opinions avec \esp{creo que...}, \esp{me parece que...}, \esp{estoy convencido/a de que...} ;
\item utilise au moins deux relances (\esp{¿Y tú qué piensas?}, \esp{¿De verdad?}, \esp{¿Puedes explicar eso?}) et une reformulation (\esp{Es decir que...}, \esp{¿Quieres decir que...?}) ;
\item emploie correctement au moins trois fois \esp{ser} ou \esp{estar}.
\end{itemize}

\textbf{Partie B — Rédaction (150 à 180 mots).} Tu es journaliste au journal de ton lycée. Rédige un article en espagnol intitulé \esp{« No a la violencia de género »} dans lequel tu :
\begin{itemize}
\item présentes le problème des violences faites aux femmes ;
\item donnes ton opinion et deux arguments ;
\item proposes deux solutions concrètes pour ton quartier ou ton école.
\end{itemize}
\end{exercice}

\newpage

\coursec{Corrigés des exercices}

\begin{corrige}[Exercice 1 — QCM : ser ou estar ?]
\begin{enumerate}
\item \esp{Mi hermano \textbf{es} médico en el Hospital Central de Yaoundé.} — Profession : \esp{ser} (identité).
\item \esp{Hoy \textbf{es} lunes 12 de mayo.} — Date : \esp{ser}.
\item \esp{La sopa \textbf{está} fría : hay que calentarla.} — État passager : \esp{estar}.
\item \esp{¿Dónde \textbf{está} el mercado de Mokolo, por favor?} — Localisation d'un lieu : \esp{estar}.
\item \esp{Mi padre \textbf{está} muy cansado después del trabajo.} — État temporaire : \esp{estar}.
\item \esp{Ese chico \textbf{es} muy listo : siempre encuentra la solución.} — \esp{ser listo} = « être intelligent » (caractéristique).
\item \esp{¡Despierta! Ya \textbf{estamos} listos para salir.} — \esp{estar listo} = « être prêt ».
\item \esp{Nosotros \textbf{estamos} de acuerdo con tu propuesta.} — Expression figée : \esp{estar de acuerdo}.
\item \esp{La fiesta de mi prima \textbf{es} en Douala, el sábado próximo.} — Lieu d'un \textbf{événement} : \esp{ser} (exception à la règle de localisation).
\item \esp{Los derechos de las mujeres \textbf{son} un tema importante del debate.} — Définition, identité : \esp{ser}.
\end{enumerate}
\end{corrige}

\begin{corrige}[Exercice 2 — Vrai ou faux ?]
\begin{enumerate}
\item \textbf{V.} La profession s'exprime avec \esp{ser} : \esp{Ella es profesora.} « Elle est professeure. »
\item \textbf{F.} C'est \esp{ser} qui exprime l'identité et la nationalité : \esp{Soy camerunés.} ; \esp{estar} exprime l'état ou le lieu.
\item \textbf{F.} \esp{Ser aburrido} signifie « être ennuyeux » (caractère) ; « être ennuyé » se dit \esp{estar aburrido}.
\item \textbf{V.} L'heure s'exprime toujours avec \esp{ser} : \esp{Son las ocho.} « Il est huit heures. »
\item \textbf{V.} \esp{Estar de acuerdo} = « être d'accord » ; c'est une expression figée avec \esp{estar}.
\item \textbf{V.} Une maladie est un état passager : \esp{estar enfermo} « être malade ».
\end{enumerate}
\end{corrige}

\begin{corrige}[Exercice 3 — Vocabulaire : la famille et la violence de genre]
\begin{enumerate}
\item \esp{Mi cuñada y mi \textbf{nuera} vienen a cenar esta noche a \textbf{el hogar} de mis suegros.} — « Ma belle-sœur et ma belle-fille viennent dîner ce soir au foyer de mes beaux-parents. » (On peut aussi dire \esp{a casa de mis suegros}.)
\item \esp{Es importante \textbf{denunciar} los casos de violencia a la policía.} — « Il est important de dénoncer les cas de violence à la police. »
\item \esp{El \textbf{maltrato} de las mujeres es un problema grave en muchas sociedades.} — « La maltraitance des femmes est un problème grave dans de nombreuses sociétés. »
\item \esp{Un hombre que pega a su esposa es un \textbf{maltratador}.} — « Un homme qui frappe sa femme est un maltraitant. »
\item \esp{Hay que defender \textbf{los derechos} de las mujeres y de los niños.} — « Il faut défendre les droits des femmes et des enfants. »
\end{enumerate}
\end{corrige}

\begin{corrige}[Exercice 4 — Conjugaison à compléter]
\begin{center}
\begin{tabular}{|l|l|l|}
\hline
\rowcolor{popblueL} \textbf{Personne} & \textbf{ser} & \textbf{estar} \\
\hline
yo & \esp{soy} & \esp{estoy} \\
\hline
tú & \esp{eres} & \esp{estás} \\
\hline
él / ella / usted & \esp{es} & \esp{está} \\
\hline
nosotros / nosotras & \esp{somos} & \esp{estamos} \\
\hline
vosotros / vosotras & \esp{sois} & \esp{estáis} \\
\hline
ellos / ellas / ustedes & \esp{son} & \esp{están} \\
\hline
\end{tabular}
\end{center}
Points de vigilance : les deux verbes sont irréguliers ; \esp{estar} prend l'accent sur la dernière syllabe sauf à la 1\iere{} personne du pluriel (\esp{estamos}) ; \esp{soy} et \esp{estoy} sont les seules formes en \esp{-oy}.
\end{corrige}

\begin{corrige}[Exercice 5 — Thème : traduis en espagnol]
\begin{enumerate}
\item \esp{Mi hermana es profesora (de primaria) en Bafoussam.} — Profession : \esp{ser}. « institutrice » = \esp{profesora de primaria} ou \esp{maestra}.
\item \esp{Son las ocho de la mañana.} — Heure : \esp{ser} ; attention, « il est huit heures » = \esp{son las ocho} (pluriel, sauf \esp{es la una}).
\item \esp{Los niños están cansados después del partido de fútbol.} — État passager : \esp{estar} ; « match » = \esp{partido} (faux ami : \esp{match} n'existe pas).
\item \esp{El mercado está muy lejos de nuestro barrio.} — Localisation : \esp{estar} ; « quartier » = \esp{barrio}.
\item \esp{Mi abuelo está enfermo desde el lunes.} — État de santé : \esp{estar} ; « depuis lundi » = \esp{desde el lunes}.
\end{enumerate}
\end{corrige}

\begin{corrige}[Exercice 6 — Transformations : changement de sens]
\begin{enumerate}
\item \esp{Juan es muy listo.} « Juan est très intelligent. » $\rightarrow$ \esp{Juan está muy listo.} « Juan est très prêt. » : \esp{ser listo} = intelligent (qualité), \esp{estar listo} = prêt (état).
\item \esp{La sopa está fría.} « La soupe est froide » (elle a refroidi, état) $\rightarrow$ \esp{La sopa es fría.} « C'est une soupe froide » (caractéristique du plat, par nature).
\item \esp{Mi prima es aburrida.} « Ma cousine est ennuyeuse » (caractère) $\rightarrow$ \esp{Mi prima está aburrida.} « Ma cousine s'ennuie / est ennuyée » (état passager).
\item \esp{Estamos listos para el examen.} « Nous sommes prêts pour l'examen » $\rightarrow$ \esp{Somos listos para el examen.} « Nous sommes doués/intelligents pour l'examen » : le sens devient une qualité intellectuelle, ce qui est peu naturel ici ; la phrase correcte au sens courant est avec \esp{estar}.
\item \esp{Pedro es muy guapo.} « Pedro est très beau » (de nature) $\rightarrow$ \esp{Pedro está muy guapo.} « Pedro est très beau (aujourd'hui, dans cette tenue) » : \esp{estar} souligne une apparence momentanée.
\end{enumerate}
\end{corrige}

\begin{corrige}[Exercice 7 — Texte à trous]
\begin{textebox}[Un dialogue à la maison — corrigé]
--- Mamá, ¿dónde \textbf{están} mis zapatos? ¡No los encuentro!\\
--- \textbf{Están} en la cocina, hijo. ¿Por qué \textbf{estás} tan nervioso esta mañana?\\
--- Porque hoy \textbf{es} el día del debate en el liceo. El tema \textbf{es} la violencia contra las mujeres.\\
--- \textbf{Es} un tema muy serio. ¿Tú qué opinas?\\
--- Yo \textbf{estoy} de acuerdo con castigar a los maltratadores. La violencia \textbf{es} inaceptable.\\
--- Muy bien, hijo. Tu padre y yo \textbf{estamos} orgullosos de ti.
\end{textebox}
Justifications : \esp{están} (lieu, pluriel) ; \esp{estás} (état passager) ; \esp{es} (date, puis définition du thème, puis jugement de valeur = caractéristique) ; \esp{estoy de acuerdo} (expression figée) ; \esp{estamos orgullosos} (sentiment, état).
Traduction : « — Maman, où sont mes chaussures ? Je ne les trouve pas ! — Elles sont dans la cuisine, mon fils. Pourquoi es-tu si nerveux ce matin ? — Parce qu'aujourd'hui c'est le jour du débat au lycée. Le thème, c'est la violence contre les femmes. — C'est un sujet très sérieux. Et toi, qu'en penses-tu ? — Je suis d'accord pour punir les maltraitants. La violence est inacceptable. — Très bien, mon fils. Ton père et moi sommes fiers de toi. »
\end{corrige}

\begin{corrige}[Exercice 8 — Appariements : expressions figées]
\begin{center}
\begin{tabularx}{\linewidth}{|X|X|}
\hline
\rowcolor{popblueL} \textbf{Español} & \textbf{Français} \\
\hline
1. \esp{estar de acuerdo} & \textbf{b.} être d'accord \\
\hline
2. \esp{estar de vacaciones} & \textbf{d.} être en vacances \\
\hline
3. \esp{estar de pie} & \textbf{e.} être debout \\
\hline
4. \esp{estar de vuelta} & \textbf{c.} être de retour \\
\hline
5. \esp{estar + gérondif} & \textbf{a.} être en train de \\
\hline
\end{tabularx}
\end{center}
Exemples : \esp{Estamos de vacaciones en Kribi.} « Nous sommes en vacances à Kribi. » ; \esp{Está hablando con el director.} « Il est en train de parler avec le directeur. »
\end{corrige}

\begin{corrige}[Exercice 9 — Compréhension et questions de langue]
\textbf{Questions de compréhension} (réponses modèles) :
\begin{enumerate}
\item \esp{La asociación « Mujeres Unidas » organiza el debate, y tiene lugar en un barrio de Douala.} « L'association "Femmes Unies" organise le débat, qui a lieu dans un quartier de Douala. »
\item \esp{Según María, muchas víctimas están calladas porque tienen miedo o porque dependen económicamente de sus maltratadores.}
\item \esp{Don Samuel propone educar a los niños desde pequeños, porque un niño que ve la violencia en casa puede ser violento mañana.}
\item \esp{La próxima reunión será el mes que viene, en la escuela del barrio.}
\end{enumerate}
\textbf{Questions de langue} :
\begin{enumerate}
\item Trois emplois de \esp{estar} : \esp{« están calladas »} (état : les victimes sont silencieuses) ; \esp{« Estoy de acuerdo »} (expression figée : être d'accord) ; \esp{« estaban convencidos »} (état d'esprit des participants). On peut aussi citer \esp{« está presente »} dans l'exercice 10 (état).
\item Dans \esp{« La violencia contra las mujeres no es un problema privado »}, \esp{ser} exprime une définition, une caractéristique essentielle et permanente : on définit la nature du problème, ce n'est ni un lieu ni un état passager.
\item Traduction de la dernière phrase : « La prochaine réunion aura lieu le mois prochain, à l'école du quartier. »
\end{enumerate}
Barème indicatif : compréhension 4 $\times$ 1 pt ; langue 2 + 1 + 1 pt ; total sur 8 points. Points de vigilance : répondre par des phrases complètes ; ne pas confondre \esp{estar callado} (se taire, état) et un lieu ; pour la traduction, \esp{será} est un futur de \esp{ser} (« aura lieu / sera »).
\end{corrige}

\begin{corrige}[Exercice 10 — Thème et version type Bac]
\textbf{Thème} (corrigé et justifications) :
\begin{enumerate}
\item \esp{Somos cameruneses y estamos orgullosos de nuestro país.} — Nationalité : \esp{ser} ; fierté ressentie (état) : \esp{estar}.
\item \esp{Son las nueve y media ; la reunión es en la sala de profesores.} — Heure : \esp{ser} ; lieu d'un \textbf{événement} : \esp{ser} (et non \esp{estar}).
\item \esp{Tu madre está preparando la comida en la cocina.} — Action en cours : \esp{estar} + gérondif.
\item \esp{Estas mujeres son valientes : denuncian las violencias.} — Qualité de caractère : \esp{ser valiente} (« être courageux » de nature).
\item \esp{La casa de mis abuelos está cerca de la iglesia.} — Localisation d'un lieu fixe : \esp{estar}.
\end{enumerate}
\textbf{Version} (traduction modèle) : « La violence de genre est une réalité encore présente dans de nombreux pays. Les victimes sont souvent seules et effrayées, mais de plus en plus de femmes sont courageuses et dénoncent leurs agresseurs. Les associations d'aide sont ouvertes jour et nuit, et leur travail est fondamental. Être informé est le premier pas : la société entière doit être attentive et engagée. »
Barème indicatif : thème 5 $\times$ 2 pts (1 pt exactitude, 1 pt choix justifié ser/estar) ; version sur 10 pts (respect du sens 6 pts, qualité du français 4 pts).
Points de vigilance : \esp{camerunés, camerunesa} prend l'accent au pluriel (\esp{cameruneses}) ; « il est neuf heures et demie » = \esp{son las nueve y media} ; « près de » = \esp{cerca de} ; \esp{suelen estar} = « ont l'habitude d'être / sont souvent ».
\end{corrige}

\begin{corrige}[Exercice 11 — Expression : débat et rédaction type Bac]
\textbf{Partie A — Débat guidé (grille d'évaluation orale).} Barème indicatif sur 10 : opinions argumentées (3 pts), relances et reformulation (2 pts), emploi correct de \esp{ser}/\esp{estar} (3 pts), prononciation et fluidité (2 pts).
Exemple de dialogue modèle :
\begin{textebox}[Modelo de debate]
--- ¿Estás a favor de castigar más duramente a los maltratadores?\\
--- Sí, creo que la violencia es un crimen y estoy convencida de que las penas deben ser más fuertes. ¿Y tú qué piensas?\\
--- Me parece que el castigo no es suficiente. ¿Quieres decir que la cárcel lo soluciona todo?\\
--- No exactamente. Es decir que el castigo es necesario, pero también lo es la educación.\\
--- ¿De verdad? ¿Puedes explicar eso?\\
--- Claro : si educamos a los niños desde pequeños, mañana serán hombres respetuosos.\\
--- Estoy de acuerdo contigo : la educación es la base de todo.
\end{textebox}
\textbf{Partie B — Rédaction modèle} (environ 165 mots) :
\begin{textebox}[No a la violencia de género]
La violencia de género es un problema grave que todavía está presente en nuestra sociedad. Muchas mujeres son víctimas de maltratos en su propio hogar y están calladas por miedo o por vergüenza.

En mi opinión, esta situación es inaceptable. Primero, la violencia es una violación de los derechos humanos : ninguna persona debe ser golpeada ni humillada. Segundo, los niños que ven la violencia en casa están traumatizados y pueden repetirla mañana.

Para luchar contra este problema, propongo dos soluciones concretas. En mi barrio, hay que crear un club de sensibilización donde las víctimas puedan hablar libremente y recibir ayuda. Además, en nuestro liceo, los profesores deben organizar debates y talleres para educar a los jóvenes en el respeto mutuo.

En conclusión, estoy convencido de que denunciar es un deber de todos. El silencio es cómplice : digamos juntos « no » a la violencia de género.
\end{textebox}
Traduction du modèle : « La violence de genre est un problème grave encore présent dans notre société. Beaucoup de femmes sont victimes de maltraitance dans leur propre foyer et se taisent par peur ou par honte. À mon avis, cette situation est inacceptable. D'abord, la violence est une violation des droits humains : personne ne doit être frappé ni humilié. Ensuite, les enfants qui voient la violence à la maison sont traumatisés et peuvent la reproduire demain. Pour lutter contre ce problème, je propose deux solutions concrètes. Dans mon quartier, il faut créer un club de sensibilisation où les victimes puissent parler librement et recevoir de l'aide. De plus, dans notre lycée, les professeurs doivent organiser des débats et des ateliers pour éduquer les jeunes au respect mutuel. En conclusion, je suis convaincu que dénoncer est un devoir pour tous. Le silence est complice : disons tous ensemble "non" à la violence de genre. »
Barème indicatif sur 20 : contenu et respect des consignes (6 pts), cohérence et organisation (4 pts), correction grammaticale dont \esp{ser}/\esp{estar} (6 pts), richesse du lexique (4 pts).
Points de vigilance : \esp{maltratador} (et non « maltratante ») ; \esp{estar callado} « se taire » ; \esp{denunciar} + nom sans préposition ; éviter le calque « las violencias son un problema » avec \esp{estar} : ici c'est une définition, donc \esp{es}.
\end{corrige}

\newpage

\coursec{Fiche bilan}

\begin{popfigure}
\begin{tikzpicture}[mindmap, grow cyclic, every node/.style={concept}, root concept/.append style={concept color=popblue, font=\bfseries\small, align=center, text=white}, level 1 concept/.append style={sibling angle=60, level distance=3.6cm, font=\scriptsize, align=center}, level 2 concept/.append style={level distance=2.2cm, font=\tiny, align=center}]
\node[root concept, align=center]{SER y ESTAR\\Diálogo familiar\\y debate}
  child[concept color=poporange]{ node[align=center]{Diálogo\\familiar}
    child[concept color=poporange!70]{ node[align=center]{Escucha\\activa} }
    child[concept color=poporange!70]{ node[align=center]{Relanzar,\\reformular} }
  }
  child[concept color=poppink]{ node[align=center]{Debate:\\violencia\\de género}
    child[concept color=poppink!70]{ node[align=center]{Opinar y\\argumentar} }
    child[concept color=poppink!70]{ node[align=center]{Conectores:\\sin embargo...} }
  }
  child[concept color=popgreen]{ node[align=center]{Léxico\\clé}
    child[concept color=popgreen!70]{ node[align=center]{Familia:\\el hogar...} }
    child[concept color=popgreen!70]{ node[align=center]{Igualdad,\\maltrato} }
  }
  child[concept color=popblue]{ node[align=center]{SER:\\esencia}
    child[concept color=popblue!70]{ node[align=center]{Identidad,\\origen, hora} }
    child[concept color=popblue!70]{ node[align=center]{Profesión,\\posesión} }
  }
  child[concept color=popgold]{ node[align=center]{ESTAR:\\estado}
    child[concept color=popgold!70]{ node[align=center]{Lugar,\\estado de ánimo} }
    child[concept color=popgold!70]{ node[align=center]{estar +\\gerundio} }
  }
  child[concept color=poppurple]{ node[align=center]{Cambios\\de sentido}
    child[concept color=poppurple!70]{ node[align=center]{ser listo /\\estar listo} }
    child[concept color=poppurple!70]{ node[align=center]{ser bueno /\\estar bueno} }
  }
  child[concept color=popdark]{ node {Traducción}
    child[concept color=popdark!70]{ node[align=center]{Thème :\\fr $\to$ esp} }
    child[concept color=popdark!70]{ node[align=center]{Version :\\esp $\to$ fr} }
  };
\end{tikzpicture}
\legende{Carte mentale de la leçon : dialogue familial, débat et les verbes \esp{ser} et \esp{estar}.}
\end{popfigure}

\begin{bilan}
\textbf{1. Lexique clé (espagnol -- français)}
\begin{itemize}
  \item \esp{el hogar} : le foyer ; \esp{los parientes / los familiares} : les parents (famille) ; \esp{la convivencia} : la vie commune ; \esp{la crianza} : l'éducation des enfants.
  \item \esp{escuchar con atención} : écouter activement ; \esp{¿Qué quieres decir?} : que veux-tu dire ? ; \esp{es decir...} : c'est-à-dire ; \esp{¿me explico?} : je me fais comprendre ?
  \item \esp{la violencia de género} : la violence de genre ; \esp{el maltrato} : la maltraitance ; \esp{la igualdad} : l'égalité ; \esp{denunciar} : dénoncer ; \esp{proteger a las víctimas} : protéger les victimes ; \esp{los derechos de la mujer} : les droits de la femme.
  \item Argumenter : \esp{en mi opinión} (à mon avis), \esp{estoy de acuerdo / no estoy de acuerdo} (je suis / ne suis pas d'accord), \esp{sin embargo} (cependant), \esp{por eso} (c'est pourquoi), \esp{en primer lugar} (d'abord).
\end{itemize}

\textbf{2. Grammaire à connaître par cœur : SER ou ESTAR}
\begin{center}
\begin{tabularx}{\linewidth}{|X|X|X|}\hline
\rowcolor{popblueL} Emploi & Verbe & Exemple \\ \hline
Identité, définition, nationalité, profession & \esp{ser} & \esp{Soy camerunés. Soy profesora.} \\ \hline
Origine, matière, possession & \esp{ser} & \esp{Somos de Yaoundé. El libro es de Ana.} \\ \hline
Date, heure, caractère permanent & \esp{ser} & \esp{Son las cinco. Ella es amable.} \\ \hline
Lieu, position (même pour les choses) & \esp{estar} & \esp{Douala está en el litoral.} \\ \hline
État passager, humeur, santé & \esp{estar} & \esp{Estoy cansado. Está enferma.} \\ \hline
Action en cours : \esp{estar} + gérondif & \esp{estar} & \esp{Estamos estudiando.} \\ \hline
\end{tabularx}
\end{center}

\textbf{3. Changements de sens avec l'adjectif}
\begin{itemize}
  \item \esp{ser listo} : être intelligent / \esp{estar listo} : être prêt.
  \item \esp{ser bueno} : être bon (de nature) / \esp{estar bueno} : être bon (au goût), être en forme.
  \item \esp{ser aburrido} : être ennuyeux / \esp{estar aburrido} : s'ennuyer.
  \item \esp{ser vivo} : être vif / \esp{estar vivo} : être vivant.
  \item Expressions figées : \esp{estar de acuerdo}, \esp{estar de vacaciones}, \esp{estar de pie}, \esp{ser de noche}.
\end{itemize}

\textbf{4. Méthodes express}
\begin{itemize}
  \item \textbf{Dialogue familial} : écouter $\to$ reformuler (\esp{¿Quieres decir que...?}) $\to$ relancer (\esp{¿Y tú qué piensas?}) $\to$ conclure.
  \item \textbf{Débat} : opinion (\esp{creo que...}) $\to$ argument (\esp{porque...}) $\to$ exemple $\to$ concession (\esp{sin embargo...}) $\to$ conclusion.
  \item \textbf{Traduire « être »} : 1) identité ou état ? 2) permanent ou passager ? 3) lieu ? $\to$ choisir \esp{ser} ou \esp{estar}, puis accorder l'adjectif.
\end{itemize}
\end{bilan}

\begin{aretenir}[Checklist avant le Bac]
\begin{itemize}
  \item[$\square$] Je sais conjuguer \esp{ser} (\esp{soy, eres, es, somos, sois, son}) et \esp{estar} (\esp{estoy, estás, está, estamos, estáis, están}) au présent sans hésiter.
  \item[$\square$] Je distingue identité / essence (\esp{ser}) et état passager / lieu (\esp{estar}).
  \item[$\square$] Je sais que le lieu s'exprime toujours avec \esp{estar}, même pour une ville ou un objet.
  \item[$\square$] Je connais les paires \esp{ser/estar + listo, bueno, aburrido, vivo, cansado} et leurs changements de sens.
  \item[$\square$] Je sais former l'action en cours : \esp{estar} + gérondif (\esp{está lloviendo}).
  \item[$\square$] Je maîtrise le lexique de la famille et de la violence de genre (\esp{maltrato, igualdad, denunciar, víctima}).
  \item[$\square$] Je sais relancer et reformuler dans un dialogue : \esp{¿Qué quieres decir? Es decir...}
  \item[$\square$] Je sais donner mon opinion et argumenter : \esp{en mi opinión, porque, sin embargo, por eso}.
  \item[$\square$] Je n'oublie pas les accents : \esp{está} (il est) $\neq$ \esp{esta} (cette).
  \item[$\square$] En thème, je vérifie l'accord de l'adjectif après \esp{ser/estar} : \esp{Las mujeres están cansadas.}
  \item[$\square$] Je n'utilise jamais la ponctuation espagnole (¿ ¡) dans une phrase française.
  \item[$\square$] Je relis ma traduction : un seul verbe « être » espagnol par phrase, choisi à bon escient.
\end{itemize}
\end{aretenir}

\findecours

\end{document}
